% !TEX program = xelatex
% Engineering Mathematics — lecture notes (Persian)
% Instructor: Abbas Ebrahimi Moghadam — compiled and edited by Aidin Shekari.
% Compile with XeLaTeX (or tectonic) from this folder.
\documentclass[11pt]{article}
\def\EMroot{../../template}
\input{\EMroot/em-fa}
\begin{document}
\EMcoverpage{ریاضی مهندسی}{آنالیز مختلط \ \textbullet\ آنالیز فوریه \ \textbullet\ معادلات دیفرانسیل با مشتقات جزئی}{Engineering Mathematics}
\pagenumbering{roman}\setcounter{page}{1}
\EMcolophon
\EMsetmodule{0}{\EMlangContents}{}
\begin{EMbooktoc}
\EMtocpart{۱}{آنالیز مختلط}
\EMtocchapter{1}{۱}{اعداد مختلط}
\EMtocitemk{1}{1}{۱.۱}{مجموعهٔ اعداد مختلط}
\EMtocitemk{1}{2}{۱.۲}{چهار عمل اصلی}
\EMtocitemk{1}{3}{۱.۳}{خواص جمع و ضرب}
\EMtocitemk{1}{4}{۱.۴}{مزدوج مختلط}
\EMtocitemk{1}{5}{۱.۵}{نامساوی مثلثی}
\EMtocitemk{1}{6}{۱.۶}{فرم قطبی اعداد مختلط}
\EMtocitemk{1}{7}{۱.۷}{ضرب دو عدد مختلط در فرم قطبی}
\EMtocitemk{1}{8}{۱.۸}{ریشهٔ \EMim{m}-ام یک عدد مختلط}
\EMtocchapter{2}{۲}{توابع مختلط، مشتق و شرایط کشی--ریمان}
\EMtocitemk{2}{1}{۲.۱}{توابع مختلط}
\EMtocitemk{2}{2}{۲.۲}{چند مکان هندسی مهم}
\EMtocitemk{2}{3}{۲.۳}{حد توابع مختلط}
\EMtocitemk{2}{4}{۲.۴}{پیوستگی توابع مختلط}
\EMtocitemk{2}{5}{۲.۵}{مشتق توابع مختلط}
\EMtocitemk{2}{6}{۲.۶}{شرایط کشی--ریمان}
\EMtocitemk{2}{7}{۲.۷}{مشتق‌پذیری توابع مختلط: مثال‌ها}
\EMtocitemk{2}{8}{۲.۸}{توابع مختلط تحلیلی}
\EMtocitemk{2}{9}{۲.۹}{شرایط کشی--ریمان در مختصات قطبی}
\EMtocchapter{3}{۳}{توابع مختلط مقدماتی}
\EMtocitemk{3}{1}{۳.۱}{تابع نمایی}
\EMtocitemk{3}{2}{۳.۲}{توابع مثلثاتی}
\EMtocitemk{3}{3}{۳.۳}{توابع هذلولی‌گون}
\EMtocitemk{3}{4}{۳.۴}{تابع لگاریتم}
\EMtocchapter{4}{۴}{نگاشت‌های خطی}
\EMtocitemk{4}{1}{۴.۱}{انتقال}
\EMtocitemk{4}{2}{۴.۲}{دوران}
\EMtocitemk{4}{3}{۴.۳}{انبساط یا انقباض}
\EMtocitemk{4}{4}{۴.۴}{نگاشت خطی کلی}
\EMtocchapter{5}{۵}{انتگرال توابع مختلط}
\EMtocitemk{5}{1}{۵.۱}{انتگرال معین و مسیر انتگرال‌گیری}
\EMtocitemk{5}{2}{۵.۲}{انتگرال معین توابع مختلط}
\EMtocitemk{5}{3}{۵.۳}{منحنی ساده بسته و ناحیهٔ همبند ساده}
\EMtocitemk{5}{4}{۵.۴}{انتگرال با تابع اولیه}
\EMtocitemk{5}{5}{۵.۵}{انتگرال توابع مختلط روی مسیر بسته}
\EMtocitemk{5}{6}{۵.۶}{فرمول انتگرال کشی}
\EMtocitemk{5}{7}{۵.۷}{کاربرد: محاسبهٔ انتگرال‌های حقیقی}
\EMtocitemk{5}{8}{۵.۸}{قضیهٔ انتگرال کشی برای چند نقطهٔ غیرتحلیلی}
\EMtocitemk{5}{9}{۵.۹}{مشتق‌های تابع تحلیلی}
\EMtocchapter{6}{۶}{سری لورنت، نقاط تکین و مانده‌ها}
\EMtocitemk{6}{1}{۶.۱}{سری لورنت}
\EMtocitemk{6}{2}{۶.۲}{مانده}
\EMtocitemk{6}{3}{۶.۳}{قضیهٔ مانده‌ها}
\EMtocitemk{6}{4}{۶.۴}{مثال‌هایی از بسط تیلور و لورنت}
\EMtocitemk{6}{5}{۶.۵}{مثال‌هایی از نقاط تکین و مانده}
\EMtocchapter{7}{۷}{محاسبهٔ انتگرال‌های معین به روش مانده‌ها}
\EMtocitemk{7}{1}{۷.۱}{انتگرال‌های دستهٔ اول}
\EMtocitemk{7}{2}{۷.۲}{انتگرال‌های دستهٔ دوم}
\EMtocpart{۲}{آنالیز فوریه}
\EMtocchapter{8}{۸}{سری فوریه}
\EMtocitemk{8}{1}{۸.۱}{توابع متناوب}
\EMtocitemk{8}{2}{۸.۲}{یادآوری از مبحث مثلثات}
\EMtocitemk{8}{3}{۸.۳}{یادآوری: انتگرال‌ها و تعامد}
\EMtocitemk{8}{4}{۸.۴}{سری فوریه}
\EMtocitemk{8}{5}{۸.۵}{یادآوری: انتگرال معین حاصل‌ضرب دو تابع}
\EMtocitemk{8}{6}{۸.۶}{مثال‌های سری فوریه}
\EMtocitemk{8}{7}{۸.۷}{سری فوریهٔ توابع متناوب با دورهٔ تناوب دلخواه}
\EMtocitemk{8}{8}{۸.۸}{قضیهٔ دریشله}
\EMtocitemk{8}{9}{۸.۹}{حل چند مسئله}
\EMtocchapter{9}{۹}{خواص سری فوریه، قضیهٔ پارسوال و صورت مختلط}
\EMtocitemk{9}{1}{۹.۱}{تقریب یک تابع متناوب توسط یک سری مثلثاتی با تعداد جملات محدود}
\EMtocitemk{9}{2}{۹.۲}{مشتق و انتگرال سری فوریه}
\EMtocitemk{9}{3}{۹.۳}{قضیهٔ رایلی}
\EMtocitemk{9}{4}{۹.۴}{قضیهٔ پارسوال}
\EMtocitemk{9}{5}{۹.۵}{صورت مختلط سری فوریه}
\EMtocchapter{10}{۱۰}{انتگرال فوریه}
\EMtocitemk{10}{1}{۱۰.۱}{تحلیل فوریه برای توابع غیرتناوبی}
\EMtocitemk{10}{2}{۱۰.۲}{انتگرال فوریه}
\EMtocitemk{10}{3}{۱۰.۳}{صورت مختلط انتگرال فوریه (تبدیل فوریه)}
\EMtocitemk{10}{4}{۱۰.۴}{معرفی تابع \EMim{\operatorname{Si}(u)}}
\EMtocitemk{10}{5}{۱۰.۵}{مثال‌های انتگرال فوریه}
\EMtocitemk{10}{6}{۱۰.۶}{بعضی خواص انتگرال فوریه}
\EMtocpart{۳}{معادلات دیفرانسیل با مشتقات جزئی و تبدیل لاپلاس}
\EMtocchapter{11}{۱۱}{معادلهٔ موجی یک‌بعدی و روش جداسازی متغیرها}
\EMtocitemk{11}{1}{۱۱.۱}{مقدمه}
\EMtocitemk{11}{2}{۱۱.۲}{معادلهٔ دیفرانسیل موجی یک‌بعدی}
\EMtocitemk{11}{3}{۱۱.۳}{حل معادلهٔ دیفرانسیل موجی یک‌بعدی}
\EMtocitemk{11}{4}{۱۱.۴}{روش جداسازی متغیرها: مثال‌ها}
\EMtocitemk{11}{5}{۱۱.۵}{معادلهٔ موجی یک‌بعدی با شرایط مرزی غیرصفر}
\EMtocitemk{11}{6}{۱۱.۶}{معادلهٔ موجی یک‌بعدی با سرعت اولیهٔ صفر (امواج متحرک)}
\EMtocchapter{12}{۱۲}{روش دالامبر و معادلهٔ انتقال حرارت یک‌بعدی}
\EMtocitemk{12}{1}{۱۲.۱}{روش دالامبر در حل معادلهٔ موجی یک‌بعدی}
\EMtocitemk{12}{2}{۱۲.۲}{معادلهٔ دیفرانسیل انتقال حرارت}
\EMtocitemk{12}{3}{۱۲.۳}{حل معادلهٔ دیفرانسیل انتقال حرارت یک‌بعدی}
\EMtocitemk{12}{4}{۱۲.۴}{معادلهٔ انتقال حرارت یک‌بعدی بر روی مفتول نامتناهی}
\EMtocitemk{12}{5}{۱۲.۵}{تابع خطا}
\EMtocchapter{13}{۱۳}{معادلات دو‌بعدی موج و گرما: مستطیل و دایره}
\EMtocitemk{13}{1}{۱۳.۱}{معادلهٔ موجی دو‌بعدی}
\EMtocitemk{13}{2}{۱۳.۲}{معادلهٔ موجی دو‌بعدی با شرایط مرزی مستطیل}
\EMtocitemk{13}{3}{۱۳.۳}{معادلهٔ گرمای دو‌بعدی با شرایط مرزی مستطیل}
\EMtocitemk{13}{4}{۱۳.۴}{بیان لاپلاسین در مختصات قطبی}
\EMtocitemk{13}{5}{۱۳.۵}{معادلهٔ موجی دو‌بعدی با شرایط مرزی دایره}
\EMtocchapter{14}{۱۴}{تبدیل لاپلاس}
\EMtocitemk{14}{1}{۱۴.۱}{تبدیل لاپلاس یک‌طرفه}
\EMtocitemk{14}{2}{۱۴.۲}{استفاده از تبدیل لاپلاس در حل معادلات دیفرانسیل}
\end{EMbooktoc}
\clearpage\pagenumbering{arabic}\setcounter{page}{1}
\EMpartpage{1}{۱}{آنالیز مختلط}{Complex Analysis}{فصل‌های ۱ تا ۷}
\EMsetmodule{1}{اعداد مختلط}{Complex Numbers}
\EMchapterpage{1}{1}{اعداد مختلط}{Complex Numbers}
\begin{EMminitoc}
\EMtocitem{1}{مجموعهٔ اعداد مختلط}
\EMtocitem{2}{چهار عمل اصلی}
\EMtocitem{3}{خواص جمع و ضرب}
\EMtocitem{4}{مزدوج مختلط}
\EMtocitem{5}{نامساوی مثلثی}
\EMtocitem{6}{فرم قطبی اعداد مختلط}
\EMtocitem{7}{ضرب دو عدد مختلط در فرم قطبی}
\EMtocitem{8}{ریشهٔ \EMim{m}-ام یک عدد مختلط}
\end{EMminitoc}
\EMchapterend
\EMhold

\EMsection{مجموعهٔ اعداد مختلط}{مجموعهٔ اعداد مختلط}

\begin{EMdefinition}{اعداد مختلط}

مجموعهٔ اعداد مختلط تعمیم اعداد واقع بر محور اعداد حقیقی به صفحه‌ای از اعداد دارای دو مولفهٔ حقیقی و موهومی می‌باشد.
\EMdisp{\mathbb{R}\to\mathbb{C}}

\end{EMdefinition}

\EMdisp{\mathbb{C}\EMbr =\{(x,y)\mid x,y\in\mathbb{R}\},\qquad\EMbr  z\EMbr =(x,y),\qquad\EMbr  x\EMbr =\operatorname{Re}\{z\},\qquad\EMbr  y\EMbr =\operatorname{Im}\{z\}}

\EMfigure{m01-complex-plane}{نمایش عدد مختلط \EMim{z=(x,y)} در صفحهٔ مختلط؛ اندازهٔ \EMim{|z|} و
آرگومان \EMim{\theta}}

اندازه یا نرم (\lr{norm}) عدد مختلط:
\EMdisp{|z|\EMbr =\sqrt{x^2+y^2}}

فاز یا آرگومان:
\EMdisp{\theta\EMbr =\measuredangle z\EMbr =\tan^{-1}\!\left(\frac{y}{x}\right),\qquad\EMbr  z\EMbr =|z|\measuredangle\theta}
\EMdisp{x\EMbr =|z|\cos\theta,\qquad\EMbr  y\EMbr =|z|\sin\theta}

\EMdisp{z\EMbr =x\times(1,0)+y\times(0,1)\EMbr =x+iy}

\begin{EMimportant}{}

\EMdisp{i^2\EMbr =-1}

\end{EMimportant}

\EMsection{چهار عمل اصلی}{چهار عمل اصلی}

فرض کنید \EMim{z_1=(a_1,b_1)=a_1+ib_1} و \EMim{z_2=(a_2,b_2)=a_2+ib_2}.

\textbf{(\lr{I}) جمع:}
\EMdisp{z_1+z_2\EMbr =(a_1+a_2,\,b_1+b_2)\EMbr =a_1+a_2+i(b_1+b_2)}

\textbf{(\lr{II}) تفریق:}
\EMdisp{z_1-z_2\EMbr =(a_1-a_2,\,b_1-b_2)\EMbr =a_1-a_2+i(b_1-b_2)}

\textbf{(\lr{III}) ضرب:}
\EMdisp{\begin{aligned}
z_1\times z_2&=(a_1+ib_1)\times(a_2+ib_2)\\
&=a_1\times a_2+a_1\times ib_2+ib_1\times a_2+ib_1\times ib_2\\
&=(a_1a_2-b_1b_2)+i(a_1b_2+b_1a_2)
\end{aligned}}

ضرب در یک عدد حقیقی:
\EMdisp{\forall\alpha\in\mathbb{R}:\quad\EMbr  \alpha z_1\EMbr =\alpha a_1+i\alpha b_1}

\textbf{(\lr{IV}) تقسیم:}
\EMdisp{\begin{aligned}
z_1\div z_2&=(a_1+ib_1)\div(a_2+ib_2)=\frac{a_1+ib_1}{a_2+ib_2}=\frac{(a_1+ib_1)(a_2-ib_2)}{(a_2+ib_2)(a_2-ib_2)}\\
&=\frac{(a_1a_2+b_1b_2)+i(b_1a_2-a_1b_2)}{a_2^2+b_2^2}
=\frac{a_1a_2+b_1b_2}{a_2^2+b_2^2}+i\,\frac{b_1a_2-a_1b_2}{a_2^2+b_2^2}
\end{aligned}}

\EMhold

\EMsection{خواص جمع و ضرب}{خواص جمع و ضرب}

\begin{EMtheorem}{میدان بودن اعداد مختلط}

مجموعهٔ اعداد مختلط همراه با دو عمل جمع و ضرب تشکیل یک میدان غیرترتیب‌پذیر می‌دهد.

\end{EMtheorem}

برای \EMim{z_1,z_2,z_3\in\mathbb{C}}:

\EMdisp{z_1+z_2\in\mathbb{C},\qquad\EMbr  z_1\times z_2\in\mathbb{C}}
\EMdisp{z_1+z_2\EMbr =z_2+z_1,\qquad\EMbr  z_1\times z_2\EMbr =z_2\times z_1}
\EMdisp{(z_1+z_2)+z_3\EMbr =z_1+(z_2+z_3),\qquad\EMbr  (z_1z_2)z_3\EMbr =z_1(z_2z_3)}
\EMdisp{z_1(z_2+z_3)\EMbr =z_1z_2+z_1z_3}
\EMdisp{z_1+(0,0)\EMbr =z_1+0\EMbr =z_1,\qquad\EMbr  z_1+(-z_1)\EMbr =(0,0)\EMbr =0}
\EMdisp{-z_1\EMbr =-(a_1,b_1)\EMbr =(-a_1,-b_1)\EMbr =-a_1-ib_1}
\EMdisp{z_1\times(1,0)\EMbr =z_1\times1\EMbr =z_1,\qquad\EMbr  z_1\times z_1^{-1}\EMbr =(1,0)\EMbr =1}
\EMdisp{z_1^{-1}\EMbr =\left(\frac{a_1}{a_1^2+b_1^2},\frac{-b_1}{a_1^2+b_1^2}\right)\EMbr =\frac{a_1}{a_1^2+b_1^2}-i\,\frac{b_1}{a_1^2+b_1^2}}

\begin{EMremark}{عدم ترتیب‌پذیری}

در مجموعهٔ اعداد مختلط نمی‌توان رابطهٔ ترتیب \EMim{z_1<z_2} تعریف کرد.

\end{EMremark}

\EMhold

\EMsection{مزدوج مختلط}{مزدوج مختلط}

\begin{EMdefinition}{مزدوج مختلط یک عدد}

\EMdisp{z\EMbr =x+iy\quad\EMbr \EMbr \Longrightarrow\quad\EMbr  \bar z\triangleq x-iy}

\end{EMdefinition}

\begin{EMexample}{}

مزدوج عدد مختلط \EMim{z=5+2i} را در صفحهٔ مختلط رسم کنید.

\end{EMexample}

\EMfigure{m01-conjugate}{عدد مختلط \EMim{z=5+2i} و مزدوج آن \EMim{\bar z=5-2i}}

\begin{EMexercise}{}

عبارت‌های زیر را ثابت کنید.
\EMdisp{\operatorname{Re}\{z\}\EMbr =\frac12\,(z+\bar z),\qquad\EMbr  \operatorname{Im}\{z\}\EMbr =\frac{1}{2i}\,(z-\bar z),\qquad\EMbr  |z|^2\EMbr =z\cdot\bar z}
\EMdisp{\overline{z_1+z_2}\EMbr =\bar z_1+\bar z_2,\qquad\EMbr  \overline{z_1\cdot z_2}\EMbr =\bar z_1\cdot\bar z_2,\qquad\EMbr  \overline{\left(\frac{z_1}{z_2}\right)}\EMbr =\frac{\bar z_1}{\bar z_2}}
\EMdisp{|z|\EMbr =|\bar z|,\qquad\EMbr  |\operatorname{Re}\{z\}|\le|z|,\qquad\EMbr  |\operatorname{Im}\{z\}|\le|z|}

\end{EMexercise}

\EMkeepfig{m01-triangle}{3}

\EMsection{نامساوی مثلثی}{نامساوی مثلثی}

\EMfigure{m01-triangle}{نامساوی مثلثی: مثلثی با رئوس \EMim{0}، \EMim{z_1} و \EMim{z_1+z_2}}

\begin{EMexercise}{}

\begin{itemize}
\tightlist
\item
  نامساوی مثلثی را برای هر دو عدد مختلط ثابت کنید: \EMim{|z_1+z_2|\le|z_1|+|z_2|}.
\item
  فرم کلی نامساوی مثلثی را برای \EMim{n} عدد مختلط دلخواه ثابت کنید:
  \EMdisp{|z_1+z_2+\cdots+z_n|\le|z_1|+|z_2|+\cdots+|z_n|}
\item
  ثابت کنید: \EMim{|z_1+z_2|\ge|z_1|-|z_2|}.
\end{itemize}

\end{EMexercise}

\EMsection{فرم قطبی اعداد مختلط}{فرم قطبی اعداد مختلط}

\EMdisp{z\EMbr =x+iy,\qquad\EMbr  \left.\begin{aligned}x&=|z|\cos\theta\\ y&=|z|\sin\theta\end{aligned}\right\}\;\EMbr \Longrightarrow\; z\EMbr =|z|\,(\cos\theta+i\sin\theta)}

\EMdisp{\theta\EMbr =\operatorname{Arg}(z)\EMbr =\tan^{-1}(y/x),\qquad\EMbr  -\pi<\theta\le\pi}
\EMdisp{\arg(z)\EMbr =\operatorname{Arg}(z)+2k\pi,\qquad\EMbr  k\in\mathbb{Z}}

\EMsection{ضرب دو عدد مختلط در فرم قطبی}{ضرب دو عدد مختلط در فرم قطبی}

\EMdisp{z_1\EMbr =|z_1|(\cos\theta_1+i\sin\theta_1),\qquad\EMbr  z_2\EMbr =|z_2|(\cos\theta_2+i\sin\theta_2)}

\EMdisp{\begin{aligned}
z_1\cdot z_2&=|z_1||z_2|(\cos\theta_1+i\sin\theta_1)(\cos\theta_2+i\sin\theta_2)\\
&=|z_1||z_2|\big[(\cos\theta_1\cos\theta_2-\sin\theta_1\sin\theta_2)\\
&\qquad\qquad+\,i(\sin\theta_1\cos\theta_2+\cos\theta_1\sin\theta_2)\big]
\end{aligned}}

\begin{EMimportant}{}

\EMdisp{z_1\cdot z_2\EMbr =|z_1||z_2|\big[\cos(\theta_1+\theta_2)+i\sin(\theta_1+\theta_2)\big]}

\end{EMimportant}

\begin{EMremark}{نتیجهٔ ۱}

در هنگام ضرب دو عدد مختلط، اندازه‌ها در هم ضرب و فازها با هم جمع می‌شوند.

\end{EMremark}

\begin{EMremark}{نتیجهٔ ۲}

اگر یک عدد مختلط را به توان یک عدد طبیعی برسانیم، اندازه را به توان آن عدد طبیعی رسانده، آرگومان (فاز) را در آن عدد طبیعی ضرب می‌کنیم (به روش استقراء ثابت کنید). به عبارت دیگر:
\EMdisp{z^k\EMbr =|z|^k(\cos k\theta+i\sin k\theta)}

\end{EMremark}

\EMdisp{z^k\EMbr =\big[|z|(\cos\theta+i\sin\theta)\big]^k\EMbr =|z|^k(\cos\theta+i\sin\theta)^k\EMbr =|z|^k(\cos k\theta+i\sin k\theta)}

\begin{EMimportant}{رابطهٔ دموار}

\EMdisp{(\cos\theta+i\sin\theta)^k\EMbr =\cos k\theta+i\sin k\theta}

\end{EMimportant}

\begin{EMexercise}{}

با استفاده از رابطهٔ دموار \EMim{\cos3\theta} و \EMim{\cos4\theta} را بر حسب \EMim{\cos\theta} و \EMim{\sin\theta} به دست آورید.

\end{EMexercise}

\EMsection{ریشهٔ \EMim{m}-ام یک عدد مختلط}{ریشهٔ m-ام یک عدد مختلط}

\EMdisp{w\EMbr =\sqrt[m]{z}\EMbr =z^{1/m},\qquad\EMbr  z\EMbr =|z|(\cos\theta+i\sin\theta)}

\EMdisp{z\EMbr =w^m,\qquad\EMbr  w\EMbr =|w|(\cos\varphi+i\sin\varphi)}

\EMdisp{z\EMbr =|z|(\cos\theta+i\sin\theta)\EMbr =|w|^m(\cos m\varphi+i\sin m\varphi)}

\EMdisp{|w|^m\EMbr =|z|\;\EMbr \Longrightarrow\;|w|\EMbr =\sqrt[m]{|z|}}

\EMdisp{m\varphi\EMbr =\theta+2k\pi\;\EMbr \Longrightarrow\;\varphi\EMbr =\frac{\theta+2k\pi}{m},\qquad\EMbr  k\EMbr =0,1,2,\dots,m-1}

\begin{EMimportant}{}

\EMdisp{w\EMbr =\sqrt[m]{z}\EMbr =\sqrt[m]{|z|}\left[\cos\!\left(\frac{\theta+2k\pi}{m}\right)+i\sin\!\left(\frac{\theta+2k\pi}{m}\right)\right],\qquad\EMbr  k\EMbr =0,1,2,\dots,m-1}

\end{EMimportant}

\begin{EMexample}{}

ریشه‌های معادلهٔ \EMim{z^n=1} چیست؟

\EMdisp{z\EMbr =\sqrt[n]{1},\qquad\EMbr  1\EMbr =1\measuredangle0\;\EMbr \Longrightarrow\; z\EMbr =\cos\!\left(\frac{2k\pi}{n}\right)+i\sin\!\left(\frac{2k\pi}{n}\right),\qquad\EMbr  k\EMbr =0,1,\dots,n-1}
با تعریف \EMim{\omega=1\measuredangle\dfrac{2\pi}{n}} ریشه‌ها عبارت‌اند از
\EMdisp{z\EMbr =1,\ \omega,\ \omega^2,\ \dots,\ \omega^{n-1}}

\end{EMexample}

\EMfigure{m01-roots-of-unity}{ریشه‌های \EMim{z=\sqrt[3]{1}}، \EMim{z=\sqrt[4]{1}} و
\EMim{z=\sqrt[5]{1}} روی دایرهٔ واحد}

\begin{EMexercise}{}

\begin{itemize}
\tightlist
\item
  ریشه‌های معادلهٔ \EMim{z^6+6+8i=0} را به دست آورده در صفحهٔ اعداد مختلط رسم کنید.
\item
  درستی اتحاد \EMim{1+z+z^2+\cdots+z^n=\dfrac{1-z^{n+1}}{1-z}} را تحقیق کنید و سپس اتحاد زیر (اتحاد لاگرانژ) را نتیجه بگیرید.
  \EMdisp{1+\cos\theta+\cos2\theta+\cdots+\cos n\theta\EMbr =\frac12+\frac{\sin\!\left[\left(n+\frac12\right)\theta\right]}{2\sin(\theta/2)}}
\end{itemize}

\end{EMexercise}

\EMsetmodule{2}{توابع مختلط، مشتق و شرایط کشی--ریمان}{Complex Functions, the Derivative and the Cauchy--Riemann Conditions}
\EMchapterpage{2}{2}{توابع مختلط، مشتق و شرایط کشی--ریمان}{Complex Functions, the Derivative and the Cauchy--Riemann Conditions}
\begin{EMminitoc}
\EMtocitem{1}{توابع مختلط}
\EMtocitem{2}{چند مکان هندسی مهم}
\EMtocitem{3}{حد توابع مختلط}
\EMtocitem{4}{پیوستگی توابع مختلط}
\EMtocitem{5}{مشتق توابع مختلط}
\EMtocitem{6}{شرایط کشی--ریمان}
\EMtocitem{7}{مشتق‌پذیری توابع مختلط: مثال‌ها}
\EMtocitem{8}{توابع مختلط تحلیلی}
\EMtocitem{9}{شرایط کشی--ریمان در مختصات قطبی}
\end{EMminitoc}
\EMchapterend
\EMkeepfig{m02-mapping}{3}

\EMsection{توابع مختلط}{توابع مختلط}

\EMfigure{m02-mapping}{تابع مختلط \EMim{w=f(z)} نقاط دامنه در صفحهٔ \EMim{z} را به نقاط برد در
صفحهٔ \EMim{w} می‌نگارد}

\EMdisp{z\EMbr =x+iy,\qquad\EMbr  w\EMbr =u+iv}
\EMdisp{f:\mathbb{C}\to\mathbb{C},\qquad\EMbr  f:\begin{cases}u=u(x,y)\\ v=v(x,y)\end{cases}}

\begin{EMexample}{}

تابع \EMim{w=z^2} را در نظر بگیرید:
\EMdisp{u+iv\EMbr =(x+iy)^2\EMbr =(x^2-y^2)+i(2xy)\quad\EMbr \EMbr \Longrightarrow\quad\EMbr \begin{cases}u=x^2-y^2\\ v=2xy\end{cases}}

\end{EMexample}

\EMkeepfig{m02-loci}{3}

\EMsection{چند مکان هندسی مهم}{چند مکان هندسی مهم}

\EMfigure{m02-loci}{مکان‌های هندسی مهم: دایره، \EMim{\rho}-همسایگی نقطهٔ \EMim{a} و ناحیهٔ
حلقوی}

\EMdisp{|z|\EMbr =1,\qquad\EMbr  |z-a|\EMbr =\rho,\qquad\EMbr  |z-a|<\rho,\qquad\EMbr  \rho_1<|z-a|<\rho_2}

مجموعهٔ \EMim{|z-a|<\rho} را \EMim{\rho}-همسایگی \EMim{a} می‌نامند.

\EMhold

\EMsection{حد توابع مختلط}{حد توابع مختلط}

\begin{EMdefinition}{حد}

اگر به ازاء هر \EMim{\varepsilon>0} بتوان \EMim{\delta>0} به‌دست آورد، به طوری که از نامساوی \EMim{|z-z_0|<\delta} نتیجه شود \EMim{|f(z)-l|<\varepsilon}، آنگاه
\EMdisp{\lim_{z\to z_0}f(z)\EMbr =l}
بنابراین به ازاء کلیهٔ نقاط داخل دایره‌ای به مرکز \EMim{z_0} و به شعاع \EMim{\delta} لازم است که \EMim{f(z)} داخل دایره‌ای به مرکز \EMim{l} و به شعاع \EMim{\varepsilon} باشد. در اینجا \EMim{\varepsilon} داده می‌شود و \EMim{\delta(\varepsilon)} را بایستی به‌دست آورد. واضح است مفهوم حد در اینجا گسترده‌تر از حد در توابع حقیقی است.

\end{EMdefinition}

\EMfigure{m02-limit}{تعریف حد: دایرهٔ \EMim{\delta}-همسایگی \EMim{z_0} در صفحهٔ \EMim{z} و
دایرهٔ \EMim{\varepsilon}-همسایگی \EMim{l} در صفحهٔ \EMim{w}}

\EMhold

\EMsection{پیوستگی توابع مختلط}{پیوستگی توابع مختلط}

\begin{EMdefinition}{پیوستگی در یک نقطه}

تابع \EMim{w=f(z)} در نقطهٔ \EMim{z=z_0} پیوسته است هرگاه سه شرط زیر برقرار باشد:

\begin{enumerate}
\def\labelenumi{\arabic{enumi}.}
\tightlist
\item
  \EMim{f(z)} در نقطهٔ \EMim{z=z_0} معین باشد.
\item
  \EMim{\displaystyle\lim_{z\to z_0}f(z)} وجود داشته باشد.
\item
  \EMim{\displaystyle\lim_{z\to z_0}f(z)=f(z_0)} باشد.
\end{enumerate}

\end{EMdefinition}

\begin{EMexample}{}

آیا تابع \EMim{f(z)=\sin(1/z)} در نقطهٔ \EMim{z=0} پیوسته است؟

\end{EMexample}

\begin{EMexample}{}

آیا تابع
\EMdisp{f(z)\EMbr =\begin{cases}\sin(1/z)&z\neq0\\ 5+4i&z=0\end{cases}}
در نقطهٔ \EMim{z=0} پیوسته است؟

\end{EMexample}

\begin{EMexample}{}

آیا تابع
\EMdisp{f(z)\EMbr =\begin{cases}3z^3+4&z\neq0\\ 17&z=0\end{cases}}
در نقطهٔ \EMim{z=0} پیوسته است؟

\end{EMexample}

\EMhold

\EMsubsection{خواص توابع مختلط پیوسته}{خواص توابع مختلط پیوسته}

\begin{EMtheorem}{قضایا}

\begin{enumerate}
\def\labelenumi{\arabic{enumi}.}
\tightlist
\item
  اگر توابع \EMim{f(z),g(z)} در نقطهٔ \EMim{z_0} پیوسته باشند، آنگاه \EMim{f(z)\pm g(z)} در \EMim{z_0} پیوسته است.
\item
  اگر توابع \EMim{f(z),g(z)} در نقطهٔ \EMim{z_0} پیوسته باشند، آنگاه \EMim{f(z)\times g(z)} در \EMim{z_0} پیوسته است.
\item
  اگر توابع \EMim{f(z),g(z)} در نقطهٔ \EMim{z_0} پیوسته باشند و \EMim{g(z_0)\neq0} آنگاه \EMim{\dfrac{f(z)}{g(z)}} در \EMim{z_0} پیوسته است.
\end{enumerate}

\end{EMtheorem}

\begin{EMdefinition}{پیوستگی در یک میدان}

تابع \EMim{w=f(z)} در میدان \EMim{D} پیوسته است هرگاه تابع \EMim{f} در تمام نقاط میدان \EMim{D} پیوسته باشد.

میدان: تمام نقاط داخل یک منحنی بسته.

\end{EMdefinition}

\EMfigure{m02-domain}{میدان \EMim{D} در صفحهٔ مختلط}

\EMhold

\EMsection{مشتق توابع مختلط}{مشتق توابع مختلط}

\begin{EMdefinition}{مشتق}

تابع \EMim{w=f(z)} در نقطهٔ \EMim{z=z_0} مشتق‌پذیر است هرگاه حد زیر موجود باشد:
\EMdisp{f'(z_0)\EMbr =\lim_{|\Delta z|\to0}\frac{f(z_0+\Delta z)-f(z_0)}{\Delta z}}

\end{EMdefinition}

\begin{EMremark}{}

وجود مشتق در توابع مختلط نیازمند شرایط سخت‌تری نسبت به توابع حقیقی است، زیرا اگر در صفحهٔ \EMim{z} از \emph{تمام} جهات به نقطهٔ \EMim{z_0} نزدیک شویم، بایستی این حد موجود بوده به مسیر انتخابی بستگی به سمت \EMim{z_0} بستگی نداشته باشد.

\end{EMremark}

\EMfigure{m02-approach}{نزدیک شدن به نقطهٔ \EMim{z_0} از تمام جهات}

\EMhold

\EMsection{شرایط کشی--ریمان}{شرایط کشی–ریمان}

\begin{EMtheorem}{شرط لازم مشتق‌پذیری}

شرط لازم برای آنکه تابع مختلط \EMim{w=f(z)} مشتق‌پذیر باشد آن‌است که شرایط کشی--ریمان (شرایط مقابل) برقرار باشند:
\EMdisp{\begin{cases}\dfrac{\partial u}{\partial x}=\dfrac{\partial v}{\partial y}\\\noalign{\vskip 2mm} \dfrac{\partial u}{\partial y}=-\dfrac{\partial v}{\partial x}\end{cases}\qquad\EMbr  w\EMbr =u+iv\EMbr =f(x+iy)\EMbr =f(z)}

\end{EMtheorem}

\begin{EMproof}

\EMdisp{f'(z)\EMbr =\lim_{|\Delta z|\to0}\frac{f(z+\Delta z)-f(z)}{\Delta z},\qquad\EMbr  z\EMbr =x+iy,\qquad\EMbr  \Delta z\EMbr =\Delta x+i\Delta y}
\EMdisp{\begin{aligned}
f(z+\Delta z)-f(z)&=\big[u(x+\Delta x,y+\Delta y)+iv(x+\Delta x,y+\Delta y)\big]-\big[u(x,y)+iv(x,y)\big]\\
&=\big[u(x+\Delta x,y+\Delta y)-u(x,y)\big]+i\big[v(x+\Delta x,y+\Delta y)-v(x,y)\big]
\end{aligned}}
\EMdisp{\frac{f(z+\Delta z)-f(z)}{\Delta z}\EMbr =\frac{u(x+\Delta x,y+\Delta y)-u(x,y)}{\Delta x+i\Delta y}+i\,\frac{v(x+\Delta x,y+\Delta y)-v(x,y)}{\Delta x+i\Delta y}}
شرط لازم برای وجود مشتق \EMim{w=f(z)} آن است که اگر از هر جهت به نقطهٔ \EMim{z} نزدیک شویم (از جمله در راستای افقی یا عمودی) مقدار حد مربوط به \EMim{f'(z)} تغییر نکند.

\textbf{مسیر ۱: راستای قائم.} \EMim{\Delta z=\Delta x+i\Delta y} و \EMim{\Delta x=0}، یعنی \EMim{\Delta z=i\Delta y}:
\EMdisp{\begin{aligned}
f'(z)&=\lim_{|\Delta z|\to0}\frac{f(z+\Delta z)-f(z)}{\Delta z}\\
&=\lim_{\Delta y\to0}\left[\frac{u(x,y+\Delta y)-u(x,y)}{i\Delta y}+i\,\frac{v(x,y+\Delta y)-v(x,y)}{i\Delta y}\right]\\
&=\frac1i\,\frac{\partial u}{\partial y}+\frac{\partial v}{\partial y}=\frac{\partial v}{\partial y}-i\,\frac{\partial u}{\partial y}
\end{aligned}}

\textbf{مسیر ۲: راستای افقی.} \EMim{\Delta z=\Delta x+i\Delta y} و \EMim{\Delta y=0}، یعنی \EMim{\Delta z=\Delta x}:
\EMdisp{\begin{aligned}
f'(z)&=\lim_{|\Delta z|\to0}\frac{f(z+\Delta z)-f(z)}{\Delta z}\\
&=\lim_{\Delta x\to0}\left[\frac{u(x+\Delta x,y)-u(x,y)}{\Delta x}+i\,\frac{v(x+\Delta x,y)-v(x,y)}{\Delta x}\right]\\
&=\frac{\partial u}{\partial x}+i\,\frac{\partial v}{\partial x}
\end{aligned}}

بنابراین
\EMdisp{f'(z)\EMbr =\frac{\partial v}{\partial y}-i\,\frac{\partial u}{\partial y}\EMbr =\frac{\partial u}{\partial x}+i\,\frac{\partial v}{\partial x}}
و با مساوی قرار دادن بخش‌های حقیقی و موهومی:

\end{EMproof}

\EMfigure{m02-cr-paths}{دو مسیر نزدیک شدن به \EMim{z}: راستای قائم (\EMim{\Delta z=i\Delta y}) و
راستای افقی (\EMim{\Delta z=\Delta x})}

\begin{EMimportant}{شرایط کشی--ریمان}

\EMdisp{\frac{\partial u}{\partial x}\EMbr =\frac{\partial v}{\partial y},\qquad\EMbr  \frac{\partial u}{\partial y}\EMbr =-\frac{\partial v}{\partial x}}

\end{EMimportant}

\begin{EMremark}{نتیجه}

\begin{enumerate}
\def\labelenumi{\arabic{enumi}.}
\tightlist
\item
  اگر این شرایط برقرار نباشد مشتق در نقطهٔ \EMim{z} وجود ندارد.
\item
  روش‌های زیر جهت مشتق‌گیری پیشنهاد می‌شود:
  \EMdisp{f'(z)\EMbr =\lim_{|\Delta z|\to0}\frac{f(z+\Delta z)-f(z)}{\Delta z}\EMbr =\frac{\partial v}{\partial y}-i\,\frac{\partial u}{\partial y}\EMbr =\frac{\partial u}{\partial x}+i\,\frac{\partial v}{\partial x}\EMbr =\frac{\partial u}{\partial x}-i\,\frac{\partial u}{\partial y}\EMbr =\frac{\partial v}{\partial y}+i\,\frac{\partial v}{\partial x}}
\end{enumerate}

\end{EMremark}

\begin{EMtheorem}{شرط کافی مشتق‌پذیری}

شرط کافی برای آنکه تابع مختلط \EMim{w=f(z)} مشتق‌پذیر باشد آن‌است که شرایط کشی--ریمان برقرار باشند.

اثبات خارج از حوصلهٔ کلاس است.

\end{EMtheorem}

\begin{EMtheorem}{قضیهٔ کشی--ریمان}

شرط لازم و کافی برای آنکه تابع مختلط \EMim{w=f(z)} مشتق‌پذیر باشد آن‌است که شرایط کشی--ریمان (شرایط مقابل) برقرار باشند.
\EMdisp{\begin{cases}\dfrac{\partial u}{\partial x}=\dfrac{\partial v}{\partial y}\\\noalign{\vskip 2mm} \dfrac{\partial u}{\partial y}=-\dfrac{\partial v}{\partial x}\end{cases}\qquad\EMbr  w\EMbr =u+iv\EMbr =f(x+iy)\EMbr =f(z)}

\end{EMtheorem}

\EMsubsection{نتیجهٔ برقراری شرایط کشی--ریمان (تحلیلی بودن یک تابع)}{نتیجهٔ برقراری شرایط کشی–ریمان (تحلیلی بودن یک تابع)}

از شرایط کشی--ریمان نسبت به \EMim{x} و \EMim{y} مشتق می‌گیریم:

\EMdisp{\left.\begin{aligned}\frac{\partial^2u}{\partial x^2}&=\frac{\partial^2v}{\partial x\partial y}\\[1mm] \frac{\partial^2u}{\partial y^2}&=-\frac{\partial^2v}{\partial x\partial y}\end{aligned}\right\}\ \xrightarrow{\ \EMt{جمع}\ }\ \frac{\partial^2u}{\partial x^2}+\frac{\partial^2u}{\partial y^2}\EMbr =0}

\EMdisp{\left.\begin{aligned}\frac{\partial^2u}{\partial x\partial y}&=\frac{\partial^2v}{\partial y^2}\\[1mm] \frac{\partial^2u}{\partial x\partial y}&=-\frac{\partial^2v}{\partial x^2}\end{aligned}\right\}\ \xrightarrow{\ \EMt{تفریق}\ }\ \frac{\partial^2v}{\partial x^2}+\frac{\partial^2v}{\partial y^2}\EMbr =0}

\begin{EMremark}{نتیجه}

قسمت حقیقی و موهومی یک تابع تحلیلی در معادلهٔ لاپلاس دوبعدی صدق می‌کند.

\end{EMremark}

\EMhold

\EMsection{مشتق‌پذیری توابع مختلط: مثال‌ها}{مشتق‌پذیری توابع مختلط: مثال‌ها}

\begin{EMexample}{}

تابع \EMim{w=z^2} را در نظر بگیرید؛ آیا این تابع تحلیلی است؟

\EMdisp{w\EMbr =f(z)\EMbr =z^2\EMbr =(x+iy)^2\EMbr =(x^2-y^2)+i\,2xy,\qquad\EMbr  u(x,y)\EMbr =x^2-y^2,\qquad\EMbr  v(x,y)\EMbr =2xy}
\EMdisp{\frac{\partial u}{\partial x}\EMbr =2x\EMbr =\frac{\partial v}{\partial y}\ \EMcheck \qquad\EMbr \qquad\EMbr  \frac{\partial u}{\partial y}\EMbr =-2y\EMbr =-\frac{\partial v}{\partial x}\ \EMcheck }
\EMdisp{f'(z)\EMbr =\frac{\partial u}{\partial x}+i\,\frac{\partial v}{\partial x}\EMbr =2x+i\,2y\EMbr =2(x+iy)\EMbr =2z}
\EMdisp{\nabla^2u\EMbr =2-2\EMbr =0,\qquad\EMbr  \nabla^2v\EMbr =0+0\EMbr =0}

\end{EMexample}

\begin{EMexample}{}

تابع \EMim{w=|z|^2} را در نظر بگیرید؛ آیا این تابع تحلیلی است؟

\EMdisp{w\EMbr =f(z)\EMbr =|z|^2\EMbr =(x+iy)(x-iy)\EMbr =x^2+y^2,\qquad\EMbr  u(x,y)\EMbr =x^2+y^2,\qquad\EMbr  v(x,y)\EMbr =0}
\EMdisp{\frac{\partial u}{\partial x}\EMbr =2x,\quad\EMbr  \frac{\partial v}{\partial y}\EMbr =0,\qquad\EMbr  \frac{\partial u}{\partial y}\EMbr =2y,\quad\EMbr  -\frac{\partial v}{\partial x}\EMbr =0}
\EMdisp{\nabla^2u\EMbr =2+2\EMbr =4,\qquad\EMbr  \nabla^2v\EMbr =0+0\EMbr =0}
تحلیلی در هیچ‌جا به جز مبدأ.

\end{EMexample}

\begin{EMexercise}{}

بررسی کنید آیا هر کدام از توابع زیر تحلیلی‌اند. در صورت تحلیلی بودن مشتق آن توابع را به دست آورید.

\begin{enumerate}
\def\labelenumi{\arabic{enumi}.}
\tightlist
\item
  \EMim{f(z)=\bar z}
\item
  \EMim{f(z)=\operatorname{Re}\{z\}}
\item
  \EMim{f(z)=z^3+2z^2+3i}
\end{enumerate}

\end{EMexercise}

\EMhold

\EMsection{توابع مختلط تحلیلی}{توابع مختلط تحلیلی}

\begin{EMexample}{}

\EMim{u=x^2-y^2} قسمت حقیقی یک تابع تحلیلی \EMim{f(z)} است. مطلوب است محاسبهٔ \EMim{f(z)}.

از آنجائی که \EMim{f(z)} تحلیلی است، قسمت‌های حقیقی و موهومی تابع در شرایط کوشی--ریمان صدق می‌کنند و در نتیجه این قسمت‌ها در معادلهٔ لاپلاس دو بعدی نیز باید صدق کنند.

۱) اول چک کنیم که آیا قسمت حقیقی تابع در معادلهٔ لاپلاس ۲ بعدی صدق می‌کند؟
\EMdisp{\nabla^2u\EMbr =\frac{\partial^2}{\partial x^2}(x^2-y^2)+\frac{\partial^2}{\partial y^2}(x^2-y^2)\EMbr =2-2\EMbr =0\ \EMcheck }

۲) از روابط کشی--ریمان استفاده می‌کنیم:
\EMdisp{\frac{\partial u}{\partial x}\EMbr =2x\EMbr =\frac{\partial v}{\partial y}\quad\EMbr \EMbr \Longrightarrow\quad\EMbr  v(x,y)\EMbr =2xy+\varphi(x)}
\EMdisp{-\frac{\partial v}{\partial x}\EMbr =-2y-\varphi'(x)\EMbr =\frac{\partial u}{\partial y}\EMbr =-2y\quad\EMbr \EMbr \Longrightarrow\quad\EMbr \varphi'(x)\EMbr =0\quad\EMbr \EMbr \Longrightarrow\quad\EMbr \varphi(x)\EMbr =c,\ \ c\in\mathbb{R}}
\EMdisp{f(z)\EMbr =(x^2-y^2)+i(2xy+c)\EMbr =z^2+C}

\end{EMexample}

\begin{EMexample}{}

اگر \EMim{u=\cos x\cosh y}،
الف) نشان دهید این تابع در صفحه هارمونیک است (\EMim{\nabla^2u=0}).
ب) تابع \EMim{v(x,y)} را چنان بیابید که \EMim{u+iv} تحلیلی شود.

الف) چک کنیم که آیا قسمت حقیقی تابع در معادلهٔ لاپلاس ۲ بعدی صدق می‌کند؟
\EMdisp{\nabla^2u\EMbr =\frac{\partial^2}{\partial x^2}(\cos x\cosh y)+\frac{\partial^2}{\partial y^2}(\cos x\cosh y)\EMbr =-\cos x\cosh y+\cos x\cosh y\EMbr =0\ \EMcheck }

ب) از روابط کشی--ریمان استفاده می‌کنیم:
\EMdisp{\frac{\partial u}{\partial x}\EMbr =-\sin x\cosh y\EMbr =\frac{\partial v}{\partial y}\quad\EMbr \EMbr \Longrightarrow\quad\EMbr  v(x,y)\EMbr =-\sin x\sinh y+\varphi(x)}
\EMdisp{-\frac{\partial v}{\partial x}\EMbr =\cos x\sinh y-\varphi'(x)\EMbr =\frac{\partial u}{\partial y}\EMbr =\cos x\sinh y\quad\EMbr \EMbr \Longrightarrow\quad\EMbr \varphi'(x)\EMbr =0}
\EMdisp{\varphi(x)\EMbr =c,\ \ c\in\mathbb{R}\quad\EMbr \EMbr \Longrightarrow\quad\EMbr  v(x,y)\EMbr =-\sin x\sinh y+c}

\end{EMexample}

\EMsection{شرایط کشی--ریمان در مختصات قطبی}{شرایط کشی–ریمان در مختصات قطبی}

\EMdisp{\begin{cases}x=r\cos\theta\\ y=r\sin\theta\end{cases}\qquad\EMbr \begin{cases}r=\sqrt{x^2+y^2}\\ \theta=\tan^{-1}\dfrac{y}{x}\end{cases}\qquad\EMbr  r_x\EMbr =\frac{x}{r},\quad\EMbr \theta_x\EMbr =\frac{-y}{r^2},\quad\EMbr  r_y\EMbr =\frac{y}{r},\quad\EMbr \theta_y\EMbr =\frac{x}{r^2}}

\EMdisp{\frac{\partial u}{\partial x}\EMbr =\frac{\partial u}{\partial r}\frac{\partial r}{\partial x}+\frac{\partial u}{\partial\theta}\frac{\partial\theta}{\partial x}\EMbr =\frac xr\frac{\partial u}{\partial r}-\frac{y}{r^2}\frac{\partial u}{\partial\theta}\EMbr =\cos\theta\frac{\partial u}{\partial r}-\frac1r\sin\theta\frac{\partial u}{\partial\theta}}

\EMdisp{\frac{\partial v}{\partial y}\EMbr =\frac{\partial v}{\partial r}\frac{\partial r}{\partial y}+\frac{\partial v}{\partial\theta}\frac{\partial\theta}{\partial y}\EMbr =\frac yr\frac{\partial v}{\partial r}+\frac{x}{r^2}\frac{\partial v}{\partial\theta}\EMbr =\sin\theta\frac{\partial v}{\partial r}+\frac1r\cos\theta\frac{\partial v}{\partial\theta}}

\EMdisp{\frac{\partial u}{\partial x}\EMbr =\frac{\partial v}{\partial y}\;\EMbr \Longrightarrow\;\cos\theta\frac{\partial u}{\partial r}-\frac1r\sin\theta\frac{\partial u}{\partial\theta}\EMbr =\sin\theta\frac{\partial v}{\partial r}+\frac1r\cos\theta\frac{\partial v}{\partial\theta}}

مشابهاً

\EMdisp{\frac{\partial u}{\partial y}\EMbr =-\frac{\partial v}{\partial x}\;\EMbr \Longrightarrow\;\sin\theta\frac{\partial u}{\partial r}+\frac1r\cos\theta\frac{\partial u}{\partial\theta}\EMbr =-\cos\theta\frac{\partial v}{\partial r}+\frac1r\sin\theta\frac{\partial v}{\partial\theta}}

\begin{EMimportant}{شرایط کشی--ریمان در فرم قطبی}

\EMdisp{\frac{\partial u}{\partial r}\EMbr =\frac1r\frac{\partial v}{\partial\theta},\qquad\EMbr  -\frac1r\frac{\partial u}{\partial\theta}\EMbr =\frac{\partial v}{\partial r}}

\end{EMimportant}

\begin{EMexercise}{}

با استفاده از شرایط کشی--ریمان در فرم قطبی ثابت کنید که \EMim{u} و \EMim{v} در معادلهٔ لاپلاس در فرم قطبی صدق می‌کنند، یعنی:
\EMdisp{\nabla^2u\EMbr =u_{rr}+\frac1{r^2}u_{\theta\theta}+\frac1ru_r\EMbr =0,\qquad\EMbr  \nabla^2v\EMbr =v_{rr}+\frac1{r^2}v_{\theta\theta}+\frac1rv_r\EMbr =0}

\end{EMexercise}

\EMsetmodule{3}{توابع مختلط مقدماتی}{Elementary Complex Functions}
\EMchapterpage{3}{3}{توابع مختلط مقدماتی}{Elementary Complex Functions}
\begin{EMminitoc}
\EMtocitem{1}{تابع نمایی}
\EMtocitem{2}{توابع مثلثاتی}
\EMtocitem{3}{توابع هذلولی‌گون}
\EMtocitem{4}{تابع لگاریتم}
\end{EMminitoc}
\EMchapterend
توابع مقدماتی مختلط تعمیم توابع مقدماتی حقیقی‌اند (\lr{Elementary Complex Functions}).

\EMsection{تابع نمایی}{تابع نمایی}

تابع نمایی حقیقی \EMim{f(x)=e^x} دارای خواص زیر است:

\EMdisp{\begin{aligned}
&1)\quad f(x_1+x_2)=f(x_1)f(x_2)\qquad &&e^{x_1}e^{x_2}=e^{x_1+x_2}\\
&2)\quad \frac{d}{dx}f(x)=f(x)\qquad &&\frac{d}{dx}e^x=e^x\\
&3)\quad e^x=1+x+\frac{x^2}{2!}+\cdots+\frac{x^n}{n!}+\cdots
\end{aligned}}

تابع نمایی مختلط را با الهام از خاصیت سوم تعریف می‌کنیم و سایر خواص را در مورد این تابع تحقیق می‌نمائیم.

\begin{EMdefinition}{تابع نمایی مختلط}

\EMdisp{w\EMbr =e^z\triangleq1+z+\frac{z^2}{2!}+\cdots+\frac{z^n}{n!}+\cdots}

\end{EMdefinition}

اگر در تعریف تابع نمایی مختلط به جای \EMim{z} مقدار \EMim{z=iy} قرار دهیم:

\EMdisp{\begin{aligned}
e^{iy}&=1+iy+\frac1{2!}(iy)^2+\frac1{3!}(iy)^3+\cdots\\
&=1+iy-\frac1{2!}y^2-\frac{i}{3!}y^3+\frac1{4!}y^4+\frac{i}{5!}y^5+\cdots\\
&=\underbrace{\left(1-\frac1{2!}y^2+\frac1{4!}y^4+\cdots\right)}_{\cos y}+i\underbrace{\left(y-\frac1{3!}y^3+\frac1{5!}y^5+\cdots\right)}_{\sin y}=\cos y+i\sin y
\end{aligned}}

\begin{EMimportant}{}

\EMdisp{w\EMbr =e^z\EMbr =e^{x+iy}\EMbr =e^xe^{iy}\EMbr =e^x(\cos y+i\sin y)}

\end{EMimportant}

\begin{EMexample}{}

آیا تابع \EMim{w=e^z} تحلیلی است؟

\EMdisp{u(x,y)\EMbr =e^x\cos y,\qquad\EMbr  v(x,y)\EMbr =e^x\sin y}
\EMdisp{\frac{\partial u}{\partial x}\EMbr =e^x\cos y\EMbr =\frac{\partial v}{\partial y}\ \EMcheck \qquad\EMbr \qquad\EMbr  \frac{\partial u}{\partial y}\EMbr =-e^x\sin y\EMbr =-\frac{\partial v}{\partial x}\ \EMcheck }
چون شرایط کشی--ریمان در تمام صفحهٔ مختلط برقرار است، تابع \EMim{w=e^z} در تمام صفحهٔ مختلط مشتق‌پذیر است.

\end{EMexample}

\begin{EMexample}{}

مشتق تابع \EMim{w=e^z} را به‌دست آورید.

\EMdisp{\frac{d}{dz}e^z\EMbr =\frac{\partial u}{\partial x}+i\,\frac{\partial v}{\partial x}\EMbr =e^x\cos y+ie^x\sin y\EMbr =e^x(\cos y+i\sin y)\EMbr =e^z}

\end{EMexample}

\begin{EMimportant}{نکتهٔ مهم}

از این پس می‌توانیم هر عدد مختلط را به صورت نمائی نمایش دهیم:
\EMdisp{z\EMbr =x+iy\EMbr =|z|(\cos\theta+i\sin\theta)\EMbr =|z|\,e^{i\theta}}

\end{EMimportant}

مشاهده می‌شود که:
\EMdisp{e^{z+i2k\pi}\EMbr =e^ze^{i2k\pi}\EMbr =e^z\big[\cos(2k\pi)+i\sin(2k\pi)\big]\EMbr =e^z}
پس تابع \EMim{w=e^z} در صفحهٔ اعداد مختلط تابعی متناوب با دورهٔ تناوب \EMim{i2\pi} است.

\EMfigure{m03-exp-strip}{باند اصلی تابع نمایی: \EMim{-\pi<y\le\pi}}

\EMsection{توابع مثلثاتی}{توابع مثلثاتی}

\EMdisp{e^{i\theta}\EMbr =\cos\theta+i\sin\theta\quad\EMbr \EMbr \Longrightarrow\quad\EMbr  e^{-i\theta}\EMbr =\cos\theta-i\sin\theta}
\EMdisp{\Longrightarrow\quad\EMbr \cos\theta\EMbr =\frac{e^{i\theta}+e^{-i\theta}}{2},\qquad\EMbr  \sin\theta\EMbr =\frac{e^{i\theta}-e^{-i\theta}}{2i}}

\begin{EMdefinition}{کسینوس و سینوس مختلط}

\EMdisp{\cos z\triangleq\frac{e^{iz}+e^{-iz}}{2},\qquad\EMbr  \sin z\triangleq\frac{e^{iz}-e^{-iz}}{2i}}

\end{EMdefinition}

\EMdisp{\begin{aligned}
\cos z&=\frac12\Big[e^{i(x+iy)}+e^{-i(x+iy)}\Big]=\frac12\Big[e^{ix}e^{-y}+e^{-ix}e^{y}\Big]\\
&=\frac12\Big[e^{-y}(\cos x+i\sin x)+e^{y}(\cos x-i\sin x)\Big]\\
&=\cos x\left(\frac{e^y+e^{-y}}{2}\right)-i\sin x\left(\frac{e^y-e^{-y}}{2}\right)\\
&=\cos x\cosh y-i\sin x\sinh y
\end{aligned}}

\EMdisp{\begin{aligned}
\sin z&=\frac1{2i}\Big[e^{i(x+iy)}-e^{-i(x+iy)}\Big]=\frac1{2i}\Big[e^{ix}e^{-y}-e^{-ix}e^{y}\Big]\\
&=\frac1{2i}\Big[e^{-y}(\cos x+i\sin x)-e^{y}(\cos x-i\sin x)\Big]\\
&=\sin x\left(\frac{e^y+e^{-y}}{2}\right)+i\cos x\left(\frac{e^y-e^{-y}}{2}\right)\\
&=\sin x\cosh y+i\cos x\sinh y
\end{aligned}}

\EMsubsection{تحلیلی بودن و مشتق‌پذیری \EMim{\cos z}}{تحلیلی بودن و مشتق‌پذیری cos z}

\EMdisp{\cos z\EMbr =u(x,y)+iv(x,y)\EMbr =\cos x\cosh y-i\sin x\sinh y,\qquad\EMbr  u\EMbr =\cos x\cosh y,\quad\EMbr  v\EMbr =-\sin x\sinh y}
\EMdisp{\frac{\partial u}{\partial x}\EMbr =-\sin x\cosh y\EMbr =\frac{\partial v}{\partial y}\ \EMcheck \qquad\EMbr \qquad\EMbr  \frac{\partial u}{\partial y}\EMbr =\cos x\sinh y\EMbr =-\frac{\partial v}{\partial x}\ \EMcheck }
پس این تابع در همهٔ صفحهٔ مختلط تحلیلی است. مشتق این تابع را توسط یکی از روش‌های ذکر شده به‌دست می‌آوریم:
\EMdisp{f'(z)\EMbr =\frac{d}{dz}\cos z\EMbr =\frac{\partial u}{\partial x}+i\,\frac{\partial v}{\partial x}\EMbr =-\sin x\cosh y-i\cos x\sinh y\EMbr =-\sin z}

\EMsubsection{تحلیلی بودن و مشتق‌پذیری \EMim{\sin z}}{تحلیلی بودن و مشتق‌پذیری sin z}

\EMdisp{\sin z\EMbr =u(x,y)+iv(x,y)\EMbr =\sin x\cosh y+i\cos x\sinh y,\qquad\EMbr  u\EMbr =\sin x\cosh y,\quad\EMbr  v\EMbr =\cos x\sinh y}
\EMdisp{\frac{\partial u}{\partial x}\EMbr =\cos x\cosh y\EMbr =\frac{\partial v}{\partial y}\ \EMcheck \qquad\EMbr \qquad\EMbr  \frac{\partial u}{\partial y}\EMbr =\sin x\sinh y\EMbr =-\frac{\partial v}{\partial x}\ \EMcheck }
پس این تابع در همهٔ صفحهٔ مختلط تحلیلی است و
\EMdisp{f'(z)\EMbr =\frac{d}{dz}\sin z\EMbr =\frac{\partial u}{\partial x}+i\,\frac{\partial v}{\partial x}\EMbr =\cos x\cosh y-i\sin x\sinh y\EMbr =\cos z}

\begin{EMimportant}{}

\EMdisp{(\sin z)'\EMbr =\cos z,\qquad\EMbr  (\cos z)'\EMbr =-\sin z}

\end{EMimportant}

\EMsubsection{بی‌کرانی \EMim{\sin z} و \EMim{\cos z}}{بی‌کرانی sin z و cos z}

\EMdisp{|\cos z|\EMbr =\sqrt{\cos^2x\cosh^2y+\sin^2x\sinh^2y}\EMbr =\sqrt{\cos^2x+\sinh^2y},\qquad\EMbr (\cosh^2y\EMbr =1+\sinh^2y,\ \ \sin^2x\EMbr =1-\cos^2x)}
\EMdisp{|\sin z|\EMbr =\sqrt{\sin^2x\cosh^2y+\cos^2x\sinh^2y}\EMbr =\sqrt{\sin^2x+\sinh^2y},\qquad\EMbr (\cosh^2y\EMbr =1+\sinh^2y,\ \ \cos^2x\EMbr =1-\sin^2x)}

از آنجائی که \EMim{\sinh y} تابعی بی‌کران است، پس توابع \EMim{\cos z} و \EMim{\sin z} توابعی \textbf{بی‌کران} خواهند بود.

\EMkeepfig{m03-trig-strip}{5}

\EMsubsection{تناوب}{تناوب}

\EMdisp{\begin{aligned}
\cos(z+2\pi)&=\cos(x+iy+2\pi)=\cos(x+2\pi)\cosh y-i\sin(x+2\pi)\sinh y\\
&=\cos x\cosh y-i\sin x\sinh y=\cos z
\end{aligned}}
به همین ترتیب \EMim{\sin(z+2\pi)=\sin z}.

\EMfigure{m03-trig-strip}{باند اصلی توابع \EMim{\sin z} و \EMim{\cos z}: \EMim{-\pi<x\le\pi}}

\EMsubsection{نکته}{نکته}

\EMdisp{\cos(iy)\EMbr =\cosh y,\qquad\EMbr  \sin(iy)\EMbr =i\sinh y}
\EMdisp{\cos(x+iy)\EMbr =\cos x\cos(iy)-\sin x\sin(iy)\EMbr =\cos x\cosh y-i\sin x\sinh y}
\EMdisp{\sin(x+iy)\EMbr =\sin x\cos(iy)+\cos x\sin(iy)\EMbr =\sin x\cosh y+i\cos x\sinh y}

سایر توابع مثلثاتی:
\EMdisp{\tan z\triangleq\frac{\sin z}{\cos z},\qquad\EMbr  \cot z\triangleq\frac{\cos z}{\sin z},\qquad\EMbr  \sec z\triangleq\frac1{\cos z},\qquad\EMbr  \csc z\triangleq\frac1{\sin z}}
\EMdisp{(\tan z)'\EMbr =\sec^2z,\quad\EMbr \dots}

\EMsection{توابع هذلولی‌گون}{توابع هذلولی‌گون}

\EMdisp{\cosh z\triangleq\frac{e^z+e^{-z}}{2},\qquad\EMbr  \sinh z\triangleq\frac{e^z-e^{-z}}{2},\qquad\EMbr  e^z\EMbr =e^x(\cos y+i\sin y)}

\begin{EMexercise}{}

تحقیق کنید:
\EMdisp{\cosh z\EMbr =\cosh x\cos y+i\sinh x\sin y,\qquad\EMbr  \sinh z\EMbr =\sinh x\cos y+i\cosh x\sin y}

\end{EMexercise}

\begin{EMexercise}{}

تحقیق کنید توابع \EMim{\cosh z} و \EMim{\sinh z} در تمام صفحهٔ مختلط تحلیلی‌اند و مشتق آنها از طریق روابط زیر به دست می‌آید:
\EMdisp{(\cosh z)'\EMbr =\sinh z,\qquad\EMbr  (\sinh z)'\EMbr =\cosh z}

\end{EMexercise}

\EMdisp{\tanh z\EMbr =\frac{\sinh z}{\cosh z},\quad\EMbr  \coth z\EMbr =\frac{\cosh z}{\sinh z},\quad\EMbr  \operatorname{sech}z\EMbr =\frac1{\cosh z},\quad\EMbr  \operatorname{csch}z\EMbr =\frac1{\sinh z}}

\EMdisp{\cosh(z+i2\pi)\EMbr =\cosh(z),\qquad\EMbr  \sinh(z+i2\pi)\EMbr =\sinh(z)}

\EMfigure{m03-hyp-strip}{باند اصلی توابع هذلولی‌گون: \EMim{-\pi<y\le\pi}}

\EMsection{تابع لگاریتم}{تابع لگاریتم}

\EMdisp{w\EMbr =\ln z,\qquad\EMbr  z\EMbr =|z|e^{i\arg z}\EMbr =e^w\EMbr =e^{\ln|z|}e^{i\arg z}}
\EMdisp{w\EMbr =\ln z\EMbr =\ln|z|+i\arg z\EMbr =\ln\!\big(x^2+y^2\big)^{1/2}+i\tan^{-1}\frac yx}
\EMdisp{u(x,y)\EMbr =\frac12\ln\!\big(x^2+y^2\big),\qquad\EMbr  v(x,y)\EMbr =\tan^{-1}\frac yx}
\EMdisp{\frac{\partial u}{\partial x}\EMbr =\frac{x}{x^2+y^2}\EMbr =\frac{\partial v}{\partial y}\ \EMcheck \qquad\EMbr \qquad\EMbr  \frac{\partial u}{\partial y}\EMbr =\frac{y}{x^2+y^2}\EMbr =-\frac{\partial v}{\partial x}\ \EMcheck }

بنابراین تابع لگاریتم در تمام صفحهٔ مختلط بجز مبدأ تحلیلی است.

\EMdisp{\frac{d}{dz}\ln z\EMbr =\frac{\partial u}{\partial x}+i\,\frac{\partial v}{\partial x}\EMbr =\frac{x}{x^2+y^2}+i\,\frac{-y}{x^2+y^2}\EMbr =\frac{\bar z}{z\bar z}\EMbr =\frac1z}

\EMdisp{z\EMbr =re^{i\theta}\EMbr =re^{i(\theta+2k\pi)}}

برای آنکه تعریف تابع نقض نشود برد این تابع ناحیهٔ نشان داده شده در شکل است.

\EMfigure{m03-log-range}{برد تابع لگاریتم: نوار \EMim{-\pi<\operatorname{Im}w\le\pi}}

خواص زیر برای تابع لگاریتم برقرارند:
\EMdisp{\ln(z_1\times z_2)\EMbr =\ln z_1+\ln z_2,\qquad\EMbr  \ln\!\left(\frac{z_1}{z_2}\right)\EMbr =\ln z_1-\ln z_2,\qquad\EMbr  \ln(z^c)\EMbr =c\ln z}

\EMsetmodule{4}{نگاشت‌های خطی}{Linear Mappings}
\EMchapterpage{4}{4}{نگاشت‌های خطی}{Linear Mappings}
\begin{EMminitoc}
\EMtocitem{1}{انتقال}
\EMtocitem{2}{دوران}
\EMtocitem{3}{انبساط یا انقباض}
\EMtocitem{4}{نگاشت خطی کلی}
\end{EMminitoc}
\EMchapterend
در این فصل برخی نگاشت‌های ساده از صفحهٔ دامنه (\EMim{z}) به صفحهٔ برد (\EMim{w}) بررسی می‌شوند.

\EMkeepfig{m04-translation}{5}

\EMsection{انتقال}{انتقال}

\EMdisp{\begin{aligned}
w=z+b&=(x+iy)+\operatorname{Re}\{b\}+i\operatorname{Im}\{b\}\\
&=\big(x+\operatorname{Re}\{b\}\big)+i\big(y+\operatorname{Im}\{b\}\big)
\end{aligned}}

\EMfigure{m04-translation}{انتقال \EMim{w=z+b}: ناحیه‌ای از صفحهٔ \EMim{z} بدون تغییر شکل به اندازهٔ
\EMim{b} جابه‌جا می‌شود}

\EMkeepfig{m04-rotation}{5}

\EMsection{دوران}{دوران}

\EMdisp{w\EMbr =az,\qquad\EMbr  |a|\EMbr =1,\qquad\EMbr  a\EMbr =e^{i\varphi},\qquad\EMbr  \varphi\EMbr =\arg\{a\}}
\EMdisp{z\EMbr =|z|e^{i\theta}\;\EMbr \Longrightarrow\;w\EMbr =|z|e^{i(\theta+\varphi)}}

\EMfigure{m04-rotation}{دوران \EMim{w=az} با \EMim{|a|=1}: ناحیه‌ای از صفحهٔ \EMim{z} به اندازهٔ
\EMim{\varphi=\arg a} حول مبدأ می‌چرخد}

\EMkeepfig{m04-dilation}{5}

\EMsection{انبساط یا انقباض}{انبساط یا انقباض}

\EMdisp{w\EMbr =az,\qquad\EMbr  a\in\mathbb{R}^+}
\EMdisp{z\EMbr =|z|e^{i\theta}\;\EMbr \Longrightarrow\;w\EMbr =a|z|e^{i\theta}}

\EMfigure{m04-dilation}{انقباض \EMim{w=az} با \EMim{0<a<1}؛ برای \EMim{a>1} نگاشت یک انبساط است}

\EMsection{نگاشت خطی کلی}{نگاشت خطی کلی}

\EMdisp{w\EMbr =az+b,\qquad\EMbr  a,b\in\mathbb{C}}

الف) دوران به اندازهٔ \EMim{\arg a}
ب) انبساط یا انقباض به اندازهٔ \EMim{|a|}
ج) انتقال به اندازهٔ \EMim{b}

\begin{EMexample}{}

اگر \EMim{w=(1+i)z+(1-i)} مطلوب است تصویر مربع واحد.

\EMdisp{w(0)\EMbr =1-i,\qquad\EMbr  w(1)\EMbr =2,\qquad\EMbr  w(i)\EMbr =0,\qquad\EMbr  w(1+i)\EMbr =1+i}

\end{EMexample}

\EMfigure{m04-linear-example}{تصویر مربع واحد تحت نگاشت \EMim{w=(1+i)z+(1-i)}}

\EMsetmodule{5}{انتگرال توابع مختلط}{Integration of Complex Functions}
\EMchapterpage{5}{5}{انتگرال توابع مختلط}{Integration of Complex Functions}
\begin{EMminitoc}
\EMtocitem{1}{انتگرال معین و مسیر انتگرال‌گیری}
\EMtocitem{2}{انتگرال معین توابع مختلط}
\EMtocitem{3}{منحنی ساده بسته و ناحیهٔ همبند ساده}
\EMtocitem{4}{انتگرال با تابع اولیه}
\EMtocitem{5}{انتگرال توابع مختلط روی مسیر بسته}
\EMtocitem{6}{فرمول انتگرال کشی}
\EMtocitem{7}{کاربرد: محاسبهٔ انتگرال‌های حقیقی}
\EMtocitem{8}{قضیهٔ انتگرال کشی برای چند نقطهٔ غیرتحلیلی}
\EMtocitem{9}{مشتق‌های تابع تحلیلی}
\end{EMminitoc}
\EMchapterend
\EMsection{انتگرال معین و مسیر انتگرال‌گیری}{انتگرال معین و مسیر انتگرال‌گیری}

انتگرال توابع مختلط همچون توابع حقیقی به دو صورت مطرح می‌شود:

\begin{itemize}
\tightlist
\item
  انتگرال نامعین، که مثل محاسبهٔ تابع اولیه است؛
\item
  انتگرال معین، که به عنوان انتگرال روی مسیر شناخته می‌شود.
\end{itemize}

انتگرال معین را در حوزهٔ توابع مختلط انتگرال روی مسیر یا \lr{Line Integral} نیز می‌نامیم و با نماد زیر نشان می‌دهیم:
\EMdisp{\int_C f(z)\,dz}
به \EMim{C} مسیر انتگرال‌گیری گوئیم.

\EMfigure{m05-path}{مسیر انتگرال‌گیری \EMim{C} از نقطهٔ \EMim{A} تا نقطهٔ \EMim{B} در صفحهٔ
\EMim{z}}

\EMsubsection{مسیر انتگرال‌گیری}{مسیر انتگرال‌گیری}

مسیر انتگرال‌گیری در صفحهٔ مختلط از نقطهٔ شروع، مسیر میانی، و نقطهٔ پایانی تشکیل می‌شود. در بسیاری موارد بهترین نحوهٔ بیان مسیر، استفاده از یک متغیر کمکی است که با تغییرات صعودی خود مسیر از ابتدا تا پایان در جهت مطلوب طی می‌شود.

\EMdisp{z(t)\EMbr =x(t)+iy(t),\qquad\EMbr  a\le t\le b,\qquad\EMbr  t,a,b\in\mathbb{R}}
\EMdisp{z(a)\EMbr =A,\qquad\EMbr  z(b)\EMbr =B,\qquad\EMbr  z,A,B\in\mathbb{C}}

\EMfigure{m05-path-param}{مسیر \EMim{C} با جهت از \EMim{A=z(a)} به \EMim{B=z(b)}}

\begin{EMexample}{}

\EMdisp{z(t)\EMbr =t+i2t,\qquad\EMbr  1\le t\le2}
این مسیر پاره‌خطی است از نقطهٔ \EMim{1+2i} تا نقطهٔ \EMim{2+4i}.

\end{EMexample}

\EMfigure{m05-line-path}{مسیر \EMim{z(t)=t+i2t}، \EMim{1\le t\le2}}

\begin{EMexample}{}

\EMdisp{z(t)\EMbr =\cos t+i\sin t,\qquad\EMbr  0\le t\le2\pi}
\EMdisp{z(t)\EMbr =e^{it},\qquad\EMbr  0\le t\le2\pi}
این مسیر دایرهٔ واحد است که در جهت مثلثاتی پیموده می‌شود.

\end{EMexample}

\EMfigure{m05-circle-path}{مسیر \EMim{z(t)=e^{it}}، \EMim{0\le t\le2\pi} (دایرهٔ واحد)}

\EMkeepfig{m05-tangent}{5}

\EMsubsection{مشتق مسیر انتگرال‌گیری}{مشتق مسیر انتگرال‌گیری}

\EMdisp{\dot z(t)\EMbr =\frac{d}{dt}z(t)\EMbr =\lim_{\Delta t\to0}\frac{z(t+\Delta t)-z(t)}{\Delta t}\EMbr =\frac{d}{dt}x(t)+i\frac{d}{dt}y(t)\EMbr =\dot x(t)+i\dot y(t)}

\EMfigure{m05-tangent}{مشتق مسیر \EMim{\dot z(t)} به عنوان حد
\EMim{\frac{z(t+\Delta t)-z(t)}{\Delta t}}}

\EMsection{انتگرال معین توابع مختلط}{انتگرال معین توابع مختلط}

\EMdisp{\int_Cf(z)\,dz\EMbr =\int_C(u+iv)(dx+i\,dy)\EMbr =\int_C(u\,dx-v\,dy)+i\int_C(v\,dx+u\,dy)}
\EMdisp{\int_Cf(z)\,dz\EMbr =\int_tf\big(z(t)\big)\,\dot z(t)\,dt}

\begin{EMdefinition}{مراحل محاسبهٔ انتگرال معین مختلط روی مسیر}

\begin{enumerate}
\def\labelenumi{\arabic{enumi}.}
\tightlist
\item
  بیان مسیر به صورت پارامتریک: \EMim{z(t)=x(t)+iy(t)}، \EMim{a\le t\le b}.
\item
  محاسبهٔ مشتق \EMim{z(t)} نسبت به \EMim{t}: \EMim{\dot z(t)=\dot x(t)+i\,\dot y(t)}.
\item
  جایگذاری \EMim{z(t)} در \EMim{f(z)}.
\item
  انتگرال‌گیری از \EMim{f\big(z(t)\big)\dot z(t)} نسبت به \EMim{t} برای \EMim{a\le t\le b}.
\end{enumerate}

\end{EMdefinition}

\begin{EMexample}{}

مطلوب است انتگرال روی مسیر شکل زیر از تابع \EMim{f(z)=\bar z}.

\EMdisp{z(t)\EMbr =t+it^2,\quad\EMbr  -1\le t\le1,\qquad\EMbr  \dot z(t)\EMbr =1+i2t}
\EMdisp{f(z)\EMbr =\bar z\EMbr =x-iy,\qquad\EMbr  f\big(z(t)\big)\EMbr =\bar z(t)\EMbr =t-it^2}
\EMdisp{\begin{aligned}
\int_{-1+i}^{1+i}\bar z\,dz&=\int_{t=-1}^{1}(t-it^2)(1+i2t)\,dt\\
&=\int_{t=-1}^{1}\Big[(t+2t^3)+i(2t^2-t^2)\Big]dt=\int_{t=-1}^{1}\Big[(t+2t^3)+it^2\Big]dt\\
&=\left.\left(\frac12t^2+\frac24t^4\right)+i\times\frac13t^3\right|_{-1}^{1}=0+i\,\frac23
\end{aligned}}

\end{EMexample}

\EMfigure{m05-parabola-path}{مسیر \EMim{z(t)=t+it^2}، \EMim{-1\le t\le1} از \EMim{-1+i} تا \EMim{1+i}}

\begin{EMexample}{}

مطلوب است انتگرال روی مسیر شکل زیر از تابع \EMim{f(z)=\bar z}.

\EMdisp{z(t)\EMbr =t+i,\quad\EMbr  -1\le t\le1,\qquad\EMbr  \dot z(t)\EMbr =1}
\EMdisp{f(z)\EMbr =\bar z\EMbr =x-iy,\qquad\EMbr  f\big(z(t)\big)\EMbr =\bar z(t)\EMbr =t-i}
\EMdisp{\int_{-1+i}^{1+i}\bar z\,dz\EMbr =\int_{t=-1}^{1}(t-i)(1)\,dt\EMbr =\left.\left(\frac{t^2}{2}-it\right)\right|_{-1}^{1}\EMbr =0-i2}
بنابراین برای این تابع مختلط مقدار انتگرال به مسیر انتخابی بستگی دارد.

\end{EMexample}

\EMfigure{m05-line-path-2}{مسیر \EMim{z(t)=t+i}، \EMim{-1\le t\le1} از \EMim{-1+i} تا \EMim{1+i}}

\EMkeepfig{m05-split-path}{5}

\EMsubsection{بعضی خواص انتگرال معین}{بعضی خواص انتگرال معین}

\EMdisp{\int_C\big[k_1f_1(z)+k_2f_2(z)\big]dz\EMbr =k_1\int_Cf_1(z)\,dz+k_2\int_Cf_2(z)\,dz\qquad\EMbr \EMt{(خطی بودن)}}
\EMdisp{\int_{z_0}^{Z}f(z)\,dz\EMbr =-\int_{Z}^{z_0}f(z)\,dz\qquad\EMbr \EMt{(معکوس کردن جهت انتگرال‌گیری)}}
\EMdisp{\int_Cf(z)\,dz\EMbr =\int_{C_1}f(z)\,dz+\int_{C_2}f(z)\,dz\qquad\EMbr \EMt{(چند تکه کردن مسیر انتگرال‌گیری)}}

\EMfigure{m05-split-path}{مسیر \EMim{C} از \EMim{z_0} تا \EMim{Z} که به دو تکهٔ \EMim{C_1} و
\EMim{C_2} تقسیم شده است}

\EMhold

\EMsection{منحنی ساده بسته و ناحیهٔ همبند ساده}{منحنی ساده بسته و ناحیهٔ همبند ساده}

\begin{EMdefinition}{منحنی ساده بسته}

منحنی ساده بسته، منحنی بسته‌ای است که خود را قطع نکند.

\end{EMdefinition}

\EMfigure{m05-curve-types}{منحنی‌های ساده بسته و غیر ساده بسته}

\begin{EMdefinition}{ناحیهٔ همبند ساده}

ناحیهٔ \EMim{D} که اگر هر منحنی بستهٔ ساده‌ای به طور کامل در آن قرار گیرد، نقاط داخل منحنی همه مربوط به این ناحیه باشند.

\end{EMdefinition}

\EMfigure{m05-connected-types}{ناحیه‌های همبند ساده و غیر همبند}

\EMhold

\EMsection{انتگرال با تابع اولیه}{انتگرال با تابع اولیه}

\begin{EMtheorem}{}

فرض کنید تابع مختلط \EMim{f(z)} تابعی تحلیلی در ناحیهٔ همبند سادهٔ \EMim{D} باشد. اگر تابع \EMim{F(z)} تابع اولیهٔ \EMim{f(z)} باشد، یعنی \EMim{F'(z)=f(z)}، آنگاه داریم:
\EMdisp{\int_{z_0}^{z_1}f(z)\,dz\EMbr =F(z_1)-F(z_0)}

\end{EMtheorem}

\begin{EMexample}{}

\EMdisp{\int_0^{1+i}z^2\,dz\EMbr =\left.\frac13z^3\right|_0^{1+i}\EMbr =\frac13(1+i)^3\EMbr =\frac23(-1+i)}

\end{EMexample}

\begin{EMexample}{}

\EMdisp{\int_{-\pi i}^{\pi i}\cos z\,dz\EMbr =\sin z\Big|_{-\pi i}^{\pi i}\EMbr =2\sin(\pi i)\EMbr =2i\sinh\pi\EMbr =23.097\,i}

\end{EMexample}

\begin{EMexample}{}

\EMdisp{\int_{8+\pi i}^{8-3\pi i}e^{z/2}\,dz\EMbr =2e^{z/2}\Big|_{8+\pi i}^{8-3\pi i}\EMbr =2\left(e^{4-3\pi i/2}-e^{4+\pi i/2}\right)\EMbr =0}

\end{EMexample}

\begin{EMexample}{}

\EMdisp{\int_{-i}^{i}\frac{dz}{z}\EMbr =\ln i-\ln(-i)\EMbr =\frac{i\pi}{2}-\left(-\frac{i\pi}{2}\right)\EMbr =i\pi}

\end{EMexample}

\EMsection{انتگرال توابع مختلط روی مسیر بسته}{انتگرال توابع مختلط روی مسیر بسته}

انتگرال روی مسیر بستهٔ سادهٔ \EMim{C} که با نماد زیر بیان می‌شود از اهمیت زیادی برخوردار است.
\EMdisp{\oint_Cf(z)\,dz}

\begin{EMtheorem}{قضیهٔ انتگرال کشی}

اگر تابع \EMim{f(z)} در ناحیهٔ همبند سادهٔ \EMim{D} تحلیلی باشد، آنگاه برای هر مسیر سادهٔ بستهٔ \EMim{C} در این ناحیه،
\EMdisp{\oint_Cf(z)\,dz\EMbr =0}

\end{EMtheorem}

\EMfigure{m05-cauchy-region}{مسیر ساده بستهٔ \EMim{C} در ناحیهٔ همبند سادهٔ \EMim{D}}

\begin{EMproof}

طبق آنچه قبلاً دیدیم، از آنجا که تابع مختلط تحلیلی است، داریم:
\EMdisp{\frac{\partial u}{\partial x}\EMbr =\frac{\partial v}{\partial y},\qquad\EMbr  \frac{\partial u}{\partial y}\EMbr =-\frac{\partial v}{\partial x}}
\EMdisp{\oint_Cf(z)\,dz\EMbr =\oint_C(u+iv)(dx+i\,dy)\EMbr =\oint_C(u\,dx-v\,dy)+i\oint_C(v\,dx+u\,dy)}
طبق قضیهٔ گرین (که در ریاضی ۲ آموخته‌اید) داریم:
\EMdisp{\oint_C(L\,dx+M\,dy)\EMbr =\iint_R\left(\frac{\partial M}{\partial x}-\frac{\partial L}{\partial y}\right)dx\,dy}
بنابراین
\EMdisp{\oint_C(u\,dx-v\,dy)+i\oint_C(v\,dx+u\,dy)\EMbr =\iint_R\underbrace{\left(-\frac{\partial v}{\partial x}-\frac{\partial u}{\partial y}\right)}_{0}dx\,dy+i\iint_R\underbrace{\left(\frac{\partial u}{\partial x}-\frac{\partial v}{\partial y}\right)}_{0}dx\,dy\EMbr =0}
که در آن عبارت اول به دلیل شرط دوم کشی--ریمان و عبارت دوم به دلیل شرط اول کشی--ریمان صفر می‌شوند.

\end{EMproof}

\begin{EMexample}{}

اگر تابع در همهٔ صفحهٔ مختلط تحلیلی باشد، آنگاه انتگرال روی هر مسیر بسته‌ای صفر خواهد بود:
\EMdisp{\oint_C\cos z\,dz\EMbr =0,\qquad\EMbr  \oint_Ce^z\,dz\EMbr =0,\qquad\EMbr  \oint_Cz^n\,dz\EMbr =0,\quad\EMbr  n\EMbr =1,2,3,\dots}

\end{EMexample}

\begin{EMexample}{}

اگر تابع نقاط تحلیلی نباشد، ولی آن نقاط بر روی مسیر و داخل \EMim{C} قرار نداشته باشند، انتگرال باز هم صفر است:
\EMdisp{\oint_{|z|=1}\sec z\,dz,\qquad\EMbr  \oint_{|z|=1}\frac{dz}{z^2+4}}
(مسیر دایرهٔ واحد است.)

\end{EMexample}

\begin{EMexample}{}

اگر تابع غیرتحلیلی باشد، انتگرال روی مسیر بسته \textbf{ممکن است} صفر نباشد:
\EMdisp{\oint_{|z|=1}\bar z\,dz,\qquad\EMbr  z(t)\EMbr =\cos t+i\sin t\EMbr =e^{it},\quad\EMbr \dot z(t)\EMbr =ie^{it}}
\EMdisp{\oint_{|z|=1}\bar z\,dz\EMbr =\int_0^{2\pi}e^{-it}\,ie^{it}\,dt\EMbr =2\pi i}

\end{EMexample}

\begin{EMexample}{}

اگر تابع روی مسیر یا داخل آن غیرتحلیلی باشد، انتگرال روی مسیر بسته \textbf{ممکن است} صفر باشد:
\EMdisp{\oint_C\frac1{z^2}\,dz\EMbr =0}

\end{EMexample}

\begin{EMexercise}{}

ثابت کنید \EMim{\displaystyle\oint_C\frac1{z^2}\,dz=0}.

\end{EMexercise}

\begin{EMremark}{نتیجه ۱}

همبند بودن ناحیه در قضیهٔ انتگرال کشی اهمیت اساسی دارد.

\end{EMremark}

\begin{EMremark}{نتیجه ۲}

اگر تابع \EMim{f(z)} در ناحیهٔ همبند سادهٔ \EMim{D} تحلیلی باشد، آنگاه انتگرال معین به انتخاب مسیر مربوط نیست.

\end{EMremark}

\EMdisp{\int_{C_1}f(z)\,dz+\int_{C_2^*}f(z)\,dz\EMbr =0\;\EMbr \Longrightarrow\;\int_{C_1}f(z)\,dz\EMbr =-\int_{C_2^*}f(z)\,dz\;\EMbr \Longrightarrow\;\int_{C_1}f(z)\,dz\EMbr =\int_{C_2}f(z)\,dz}

بنابراین انتگرال از مبدأ تا مقصد به مسیر انتخابی وابسته نیست.

\EMfigure{m05-path-independence}{مسیرهای \EMim{C_1}، \EMim{C_2} و \EMim{C_2^*} (مسیر \EMim{C_2} با جهت
معکوس) میان \EMim{z_1} و \EMim{z_2}}

\EMkeepfig{m05-circle-integral}{5}

\EMsubsection{یک مثال بسیار مهم}{یک مثال بسیار مهم}

\EMdisp{\oint_C\frac{dz}{z-z_0}\EMbr =\;?\qquad\EMbr  C:\ z(t)\EMbr =z_0+re^{it}}

\EMfigure{m05-circle-integral}{دایرهٔ \EMim{C} به مرکز \EMim{z_0} و شعاع \EMim{r}}

\EMdisp{z(t)\EMbr =z_0+re^{it},\qquad\EMbr  \dot z(t)\EMbr =rie^{it}}
\EMdisp{\oint_C\frac{dz}{z-z_0}\EMbr =\int_0^{2\pi}\frac{rie^{it}}{re^{it}}\,dt\EMbr =\int_0^{2\pi}i\,dt\EMbr =2\pi i}

\begin{EMimportant}{}

\EMdisp{\oint_C\frac{dz}{z-z_0}\EMbr =2\pi i}

\end{EMimportant}

\EMhold

\EMsection{فرمول انتگرال کشی}{فرمول انتگرال کشی}

\begin{EMtheorem}{فرمول انتگرال کشی}

اگر تابع \EMim{f(z)} در ناحیهٔ همبند سادهٔ \EMim{D} تحلیلی بوده و \EMim{z_0} نقطه‌ای در داخل \EMim{C} باشد، آنگاه
\EMdisp{\frac1{2\pi i}\oint_C\frac{f(z)}{z-z_0}\,dz\EMbr =f(z_0)}

\end{EMtheorem}

\begin{EMproof}

طبق آنچه قبلاً دیدیم، از آنجا که تابع \EMim{\dfrac{f(z)}{z-z_0}} داخل ناحیهٔ همبند \EMim{\Gamma} تحلیلی است، داریم:
\EMdisp{\oint_\Gamma\frac{f(z)}{z-z_0}\,dz\EMbr =0\EMbr =\int_{ACA'}\frac{f(z)}{z-z_0}\,dz+\int_{A'B'}\frac{f(z)}{z-z_0}\,dz+\int_{B'C'B}\frac{f(z)}{z-z_0}\,dz+\int_{BA}\frac{f(z)}{z-z_0}\,dz}
\EMdisp{\lim_{\substack{A\to A'\\ B\to B'}}\int_{BA}\frac{f(z)}{z-z_0}\,dz\EMbr =-\int_{A'B'}\frac{f(z)}{z-z_0}\,dz}
در نتیجه (دو انتگرال روی پیوندها یکدیگر را خنثی می‌کنند)
\EMdisp{\lim_{\substack{A\to A'\\ B\to B'}}\int_{ACA'}\frac{f(z)}{z-z_0}\,dz+\int_{B'C'B}\frac{f(z)}{z-z_0}\,dz\EMbr =0\;\EMbr \Longrightarrow\;\oint_C\frac{f(z)}{z-z_0}\,dz\EMbr =\oint_{C'}\frac{f(z)}{z-z_0}\,dz}
(در اینجا \EMim{C'} در جهت مثلثاتی پیموده می‌شود.) ادامه:
\EMdisp{\begin{aligned}
\oint_C\frac{f(z)}{z-z_0}\,dz&=\oint_{C'}\frac{f(z)}{z-z_0}\,dz=\oint_{C'}\frac{f(z)-f(z_0)+f(z_0)}{z-z_0}\,dz\\
&=\oint_{C'}\frac{f(z)-f(z_0)}{z-z_0}\,dz+\oint_{C'}\frac{f(z_0)}{z-z_0}\,dz\\
&=\oint_{C'}g(z)\,dz+f(z_0)\oint_{C'}\frac1{z-z_0}\,dz=0+2\pi i\,f(z_0)=2\pi i\,f(z_0)
\end{aligned}}
که در آن
\EMdisp{g(z)\triangleq\begin{cases}\dfrac{f(z)-f(z_0)}{z-z_0}&z\neq z_0\\\noalign{\vskip 2mm} f'(z_0)&z=z_0\end{cases}}
تابعی تحلیلی در داخل \EMim{C'} است. بنابراین
\EMdisp{\frac1{2\pi i}\oint_C\frac{f(z)}{z-z_0}\,dz\EMbr =f(z_0)}

\end{EMproof}

\EMfigure{m05-cauchy-proof}{مسیر \EMim{\Gamma} در اثبات فرمول انتگرال کشی: منحنی \EMim{C}، دایرهٔ
کوچک \EMim{C'} به مرکز \EMim{z_0} و دو پیوند \EMim{AA'} و \EMim{BB'}}

\EMhold

\EMsection{کاربرد: محاسبهٔ انتگرال‌های حقیقی}{کاربرد: محاسبهٔ انتگرال‌های حقیقی}

\begin{EMexample}{}

مطلوب است محاسبهٔ انتگرال \EMim{\displaystyle\int_{-\infty}^{\infty}\frac{dx}{1+x^2}} با استفاده از فرمول انتگرال کشی.

تابع \EMim{g(z)} را به صورت \EMim{g(z)=\dfrac{1}{1+z^2}} در نظر می‌گیریم. با استفاده از مسیر مشخص شده در شکل انتگرال روی مسیر بسته را به دست می‌آوریم.
\EMdisp{\oint_C\frac{dz}{1+z^2}\EMbr =\oint_C\frac{\left[\dfrac1{z+i}\right]}{z-i}\,dz\EMbr =2\pi i\,\frac1{i+i}\EMbr =\pi,\qquad\EMbr  f(z)\EMbr =\frac1{z+i},\quad\EMbr  z_0\EMbr =i}
از طرف دیگر
\EMdisp{\oint_C\frac{dz}{1+z^2}\EMbr =\int_{-R}^{R}\frac{dx}{1+x^2}+\int_0^\pi\frac{iRe^{it}\,dt}{1+R^2e^{i2t}}}
در عبارت بالا جملهٔ اول وقتی \EMim{R\to\infty} مقدار مطلوب خواهد بود.
\EMdisp{\oint_C\frac{dz}{1+z^2}\EMbr =\pi\EMbr =\int_{-\infty}^{\infty}\frac{dx}{1+x^2}+I,\qquad\EMbr  I\EMbr =\lim_{R\to\infty}\int_0^\pi\frac{iRe^{it}\,dt}{1+R^2e^{i2t}}\EMbr =0}
بنابراین
\EMdisp{\int_{-\infty}^{\infty}\frac{dx}{1+x^2}\EMbr =\pi}

\end{EMexample}

\EMfigure{m05-semicircle-1}{مسیر نیم‌دایره در نیم‌صفحهٔ بالایی؛ قطب‌ها در \EMim{z=\pm i}}

\begin{EMexample}{}

مطلوب است محاسبهٔ انتگرال \EMim{\displaystyle\int_{-\infty}^{\infty}\frac{\cos\omega x}{1+x^2}\,dx} با استفاده از فرمول انتگرال کشی (\EMim{\omega>0}).

تابع \EMim{g(z)} را به صورت \EMim{g(z)=\dfrac{e^{i\omega z}}{1+z^2}} در نظر می‌گیریم. با استفاده از مسیر مشخص شده در شکل (همان نیم‌دایرهٔ شکل قبل) انتگرال روی مسیر بسته را به دست می‌آوریم.
\EMdisp{\oint_C\frac{e^{i\omega z}}{1+z^2}\,dz\EMbr =\oint_C\frac{\left[\dfrac{e^{i\omega z}}{z+i}\right]}{z-i}\,dz\EMbr =2\pi i\,\frac{e^{i\omega i}}{i+i}\EMbr =\pi e^{-\omega},\qquad\EMbr  f(z)\EMbr =\frac{e^{i\omega z}}{z+i},\quad\EMbr  z_0\EMbr =i}
از طرف دیگر
\EMdisp{\oint_C\frac{e^{i\omega z}\,dz}{1+z^2}\EMbr =\int_{-R}^{R}\frac{e^{i\omega x}\,dx}{1+x^2}+\int_0^\pi\frac{e^{i\omega Re^{it}}\,iRe^{it}\,dt}{1+R^2e^{i2t}}}
اما:
\EMdisp{I_1\EMbr =\lim_{R\to\infty}\int_0^\pi\frac{e^{i\omega Re^{it}}\,iRe^{it}\,dt}{1+R^2e^{i2t}}\to0}
\textbf{اثبات:}
\EMdisp{e^{i\omega Re^{it}}\EMbr =e^{i\omega R(\cos t+i\sin t)}\EMbr =e^{i\omega R\cos t}\,e^{-\omega R\sin t}}
\EMdisp{\lim_{R\to\infty}\left|e^{i\omega Re^{it}}\right|\EMbr =\lim_{R\to\infty}\left|e^{i\omega R\cos t}\right|\times\left|e^{-\omega R\sin t}\right|\EMbr =1\times0\EMbr =0,\qquad\EMbr  t\in[0,\pi]}
از طرف دیگر در انتگرال \EMim{I_1} بقیهٔ انتگرند زمانی که \EMim{R} بی‌نهایت می‌شود به سمت صفر میل می‌کند. پس حد \EMim{I_1} وقتی \EMim{R\to\infty} صفر می‌شود.

بنابراین
\EMdisp{\pi e^{-\omega}\EMbr =\oint_C\frac{e^{i\omega z}\,dz}{1+z^2}\EMbr =\int_{-R}^{R}\frac{e^{i\omega x}\,dx}{1+x^2}\EMbr =\int_{-R}^{R}\frac{(\cos x+i\sin x)\,dx}{1+x^2}\EMbr =\int_{-R}^{R}\frac{\cos x\,dx}{1+x^2}+i\int_{-R}^{R}\frac{\sin x\,dx}{1+x^2}}
انتگرال دوم صفر است. چرا؟
\EMdisp{\int_{-\infty}^{\infty}\frac{\cos\omega x}{1+x^2}\,dx\EMbr =\pi e^{-\omega}}

\end{EMexample}

\begin{EMexample}{}

مطلوب است محاسبهٔ انتگرال \EMim{\displaystyle I=\int_{-\infty}^{\infty}\frac{dx}{(x^2+1)(x^2+4)}} با استفاده از فرمول انتگرال کشی.

\EMdisp{\frac{1}{(x^2+1)(x^2+4)}\EMbr =\frac{A}{x^2+1}+\frac{B}{x^2+4}\EMbr =\frac{1/3}{x^2+1}+\frac{-1/3}{x^2+4}}
\EMdisp{\int_{-\infty}^{\infty}\frac{dx}{(x^2+1)(x^2+4)}\EMbr =\frac13\int_{-\infty}^{\infty}\frac{dx}{x^2+1}-\frac13\int_{-\infty}^{\infty}\frac{dx}{x^2+4}}

\textbf{انتگرال اول:} \EMim{I_1=\displaystyle\int_{-\infty}^{\infty}\frac{dx}{x^2+1}} و \EMim{g_1(z)=\dfrac1{z^2+1}}:
\EMdisp{\oint_Cg_1(z)\,dz\EMbr =\oint_C\frac1{z^2+1}\,dz\EMbr =\oint_C\frac{\left[\dfrac1{z+i}\right]}{z-i}\,dz\EMbr =2\pi i\,\frac1{i+i}\EMbr =\pi}
\EMdisp{\oint_C\frac1{z^2+1}\,dz\EMbr =\pi\EMbr =\int_{-R}^{R}\frac{dx}{1+x^2}+\int_0^\pi\frac{iRe^{it}\,dt}{1+R^2e^{i2t}}}
\EMdisp{\lim_{R\to\infty}\oint_C\frac1{z^2+1}\,dz\EMbr =\int_{-R}^{R}\frac{dx}{1+x^2}+0\EMbr =\pi\;\EMbr \Longrightarrow\;I_1\EMbr =\pi}

\textbf{انتگرال دوم:} \EMim{I_2=\displaystyle\int_{-\infty}^{\infty}\frac{dx}{x^2+4}} و \EMim{g_2(z)=\dfrac1{z^2+4}}:
\EMdisp{\oint_Cg_2(z)\,dz\EMbr =\oint_C\frac1{z^2+4}\,dz\EMbr =\oint_C\frac{\left[\dfrac1{z+2i}\right]}{z-2i}\,dz\EMbr =2\pi i\,\frac1{2i+2i}\EMbr =\frac\pi2}
\EMdisp{\oint_C\frac1{z^2+4}\,dz\EMbr =\frac\pi2\EMbr =\int_{-R}^{R}\frac{dx}{x^2+4}+\int_0^\pi\frac{iRe^{it}\,dt}{R^2e^{i2t}+4}}
\EMdisp{\lim_{R\to\infty}\oint_C\frac1{z^2+4}\,dz\EMbr =\int_{-R}^{R}\frac{dx}{x^2+4}+0\EMbr =\frac\pi2\;\EMbr \Longrightarrow\;I_2\EMbr =\frac\pi2}

در نتیجه
\EMdisp{I\EMbr =\frac13I_1-\frac13I_2\EMbr =\frac13\pi-\frac13\times\frac\pi2\EMbr =\frac\pi6}

\end{EMexample}

\EMfigure{m05-semicircle-2}{مسیر نیم‌دایره؛ قطب‌ها در \EMim{z=\pm i} و \EMim{z=\pm2i}}

\begin{EMexample}{}

مطلوب است محاسبهٔ انتگرال \EMim{\displaystyle I=\int_0^\infty\frac{dx}{x^4+1}} با استفاده از فرمول انتگرال کشی.

\EMdisp{g(z)\EMbr =\frac1{z^4+1},\qquad\EMbr  z^4+1\EMbr =0\;\EMbr \Longrightarrow\;z_1\EMbr =e^{i\pi/4},\quad\EMbr  z_2\EMbr =e^{i3\pi/4},\quad\EMbr  z_3\EMbr =e^{i5\pi/4},\quad\EMbr  z_4\EMbr =e^{i7\pi/4}}
در داخل مسیر \EMim{C} فقط دو قطب \EMim{z_1} و \EMim{z_2} قرار دارند:
\EMdisp{\begin{aligned}
\oint_Cg(z)\,dz&=\oint_C\frac1{z^4+1}\,dz=\oint_{C_1}\frac1{z^4+1}\,dz+\oint_{C_2}\frac1{z^4+1}\,dz\\
&=\oint_{C_1}\frac{\dfrac1{(z-z_2)(z-z_3)(z-z_4)}}{z-z_1}\,dz+\oint_{C_2}\frac{\dfrac1{(z-z_1)(z-z_3)(z-z_4)}}{z-z_2}\,dz\\
&=2\pi i\left.\frac1{(z-z_2)(z-z_3)(z-z_4)}\right|_{z=z_1}+2\pi i\left.\frac1{(z-z_1)(z-z_3)(z-z_4)}\right|_{z=z_2}\\
&=2\pi i\,\frac1{\underbrace{(e^{i\pi/4}-e^{i3\pi/4})}_{\sqrt2}\underbrace{(e^{i\pi/4}-e^{i5\pi/4})}_{\sqrt2+i\sqrt2}\underbrace{(e^{i\pi/4}-e^{i7\pi/4})}_{i\sqrt2}}\\
&\quad+2\pi i\,\frac1{\underbrace{(e^{i3\pi/4}-e^{i\pi/4})}_{-\sqrt2}\underbrace{(e^{i3\pi/4}-e^{i5\pi/4})}_{i\sqrt2}\underbrace{(e^{i3\pi/4}-e^{i7\pi/4})}_{-\sqrt2+i\sqrt2}}=\frac\pi{\sqrt2}
\end{aligned}}
از طرف دیگر
\EMdisp{\oint_C\frac1{z^4+1}\,dz\EMbr =\frac\pi{\sqrt2}\EMbr =\int_{-R}^{R}\frac{dx}{x^4+1}+\int_0^\pi\frac{iRe^{it}\,dt}{R^4e^{i4t}+1}}
\EMdisp{\lim_{R\to\infty}\oint_C\frac1{z^4+1}\,dz\EMbr =\int_{-\infty}^{\infty}\frac{dx}{x^4+1}+\underbrace{\lim_{R\to\infty}\int_0^\pi\frac{iRe^{it}\,dt}{R^4e^{i4t}+1}}_{0}\EMbr =\frac\pi{\sqrt2}}
\EMdisp{\int_{-\infty}^{\infty}\frac{dx}{x^4+1}\EMbr =\frac\pi{\sqrt2}\;\EMbr \Longrightarrow\;\int_0^\infty\frac{dx}{x^4+1}\EMbr =\frac\pi{2\sqrt2}}

\end{EMexample}

\EMfigure{m05-semicircle-3}{مسیر نیم‌دایره و دایره‌های \EMim{C_1} و \EMim{C_2} حول قطب‌های \EMim{z_1} و
\EMim{z_2}؛ قطب‌های \EMim{z_3} و \EMim{z_4} در نیم‌صفحهٔ پایینی}

\EMhold

\EMsection{قضیهٔ انتگرال کشی برای چند نقطهٔ غیرتحلیلی}{قضیهٔ انتگرال کشی برای چند نقطهٔ غیرتحلیلی}

\begin{EMtheorem}{}

اگر تابع \EMim{f(z)} در همهٔ نقاط ناحیهٔ \EMim{D} بجز نقاط \EMim{z_1,z_2,z_3,\dots,z_n} در داخل \EMim{C} تحلیلی باشد، آنگاه
\EMdisp{\oint_Cf(z)\,dz\EMbr =\oint_{C_1}f(z)\,dz+\oint_{C_2}f(z)\,dz+\cdots+\oint_{C_n}f(z)\,dz}
که در این رابطه، دایرهٔ \EMim{C_k} دایره‌ای به مرکز \EMim{z_k} است که داخل \EMim{C} قرار داشته سایر نقاط \EMim{z_1} تا \EMim{z_n} را شامل نشود.

\end{EMtheorem}

\EMfigure{m05-multi-points}{مسیر \EMim{C} و نقاط غیرتحلیلی \EMim{z_1,\dots,z_n} در داخل آن}

\begin{EMproof}

مسیر \EMim{\Gamma} را از \EMim{C} (در جهت مثلثاتی)، دایره‌های کوچک \EMim{C_1,\dots,C_n} (در جهت ساعت‌گرد) و پیوندهایی که آن‌ها را به \EMim{C} وصل می‌کنند می‌سازیم؛ انتگرال روی هر پیوند دو بار و در دو جهت مخالف محاسبه می‌شود و حذف می‌گردد. تابع در داخل \EMim{\Gamma} تحلیلی است، پس
\EMdisp{\oint_\Gamma f(z)\,dz\EMbr =0\;\EMbr \Longrightarrow\;\oint_Cf(z)\,dz+\oint_{C_1^{\circlearrowright}}f(z)\,dz+\cdots+\oint_{C_n^{\circlearrowright}}f(z)\,dz\EMbr =0}
\EMdisp{\oint_Cf(z)\,dz\EMbr =-\oint_{C_1^{\circlearrowright}}f(z)\,dz-\cdots-\oint_{C_n^{\circlearrowright}}f(z)\,dz\EMbr =\oint_{C_1}f(z)\,dz+\cdots+\oint_{C_n}f(z)\,dz}
که در آن \EMim{C_k^{\circlearrowright}} دایرهٔ \EMim{C_k} در جهت ساعت‌گرد و \EMim{C_k} همان دایره در جهت مثلثاتی است.

\end{EMproof}

\EMfigure{m05-multi-proof}{مسیر \EMim{\Gamma} برای اثبات: \EMim{C}، دایره‌های \EMim{C_1,\dots,C_n} و
پیوندها}

\EMhold

\EMsection{مشتق‌های تابع تحلیلی}{مشتق‌های تابع تحلیلی}

\begin{EMtheorem}{فرمول انتگرال کشی برای مشتق}

اگر تابع \EMim{f(z)} در ناحیهٔ همبند سادهٔ \EMim{D} تحلیلی بوده و \EMim{z_0} نقطه‌ای در داخل \EMim{C} باشد، آنگاه
\EMdisp{\frac1{2\pi i}\oint_C\frac{f(z)}{(z-z_0)^2}\,dz\EMbr =f'(z_0)}

\end{EMtheorem}

\begin{EMproof}

\EMdisp{f'(z_0)\EMbr =\lim_{\Delta t\to0}\frac{f(z_0+\Delta t)-f(z_0)}{\Delta t},\qquad\EMbr  f(z_0)\EMbr =\frac1{2\pi i}\oint_C\frac{f(z)}{z-z_0}\,dz,\qquad\EMbr  f(z_0+\Delta t)\EMbr =\frac1{2\pi i}\oint_C\frac{f(z)}{z-z_0-\Delta t}\,dz}
\EMdisp{\begin{aligned}
\frac{f(z_0+\Delta t)-f(z_0)}{\Delta t}&=\frac1{2\pi i\,\Delta t}\left\{\oint_C\frac{f(z)}{z-z_0-\Delta t}\,dz-\oint_C\frac{f(z)}{z-z_0}\,dz\right\}\\
&=\frac1{2\pi i\,\Delta t}\oint_Cf(z)\left[\frac1{z-z_0-\Delta t}-\frac1{z-z_0}\right]dz\\
&=\frac1{2\pi i\,\Delta t}\oint_Cf(z)\left[\frac{\Delta t}{(z-z_0-\Delta t)(z-z_0)}\right]dz\\
&=\frac1{2\pi i}\oint_Cf(z)\left[\frac1{(z-z_0-\Delta t)(z-z_0)}\right]dz
\end{aligned}}
\EMdisp{f'(z_0)\EMbr =\lim_{\Delta t\to0}\frac{f(z_0+\Delta t)-f(z_0)}{\Delta t}\EMbr =\lim_{\Delta t\to0}\frac1{2\pi i}\oint_Cf(z)\left[\frac1{(z-z_0-\Delta t)(z-z_0)}\right]dz\EMbr =\frac1{2\pi i}\oint_C\frac{f(z)}{(z-z_0)^2}\,dz}

\end{EMproof}

\begin{EMtheorem}{تعمیم فرمول انتگرال کشی}

اگر تابع \EMim{f(z)} در ناحیهٔ همبند سادهٔ \EMim{D} تحلیلی بوده و \EMim{z_0} نقطه‌ای در داخل \EMim{C} باشد، آنگاه
\EMdisp{\frac{n!}{2\pi i}\oint_C\frac{f(z)}{(z-z_0)^{n+1}}\,dz\EMbr =f^{(n)}(z_0)}

\end{EMtheorem}

\begin{EMexercise}{}

از طریق استقراء ریاضی ثابت کنید.

\end{EMexercise}

\begin{EMexample}{}

مطلوب است محاسبهٔ انتگرال زیر. (\EMim{C} دایرهٔ واحد است.)
\EMdisp{I\EMbr =\oint_C\frac{z^2}{(2z-1)^2}\,dz}

\EMdisp{I\EMbr =\oint_C\frac{z^2}{4(z-1/2)^2}\,dz\EMbr =\frac14\oint_C\frac{z^2}{(z-1/2)^2}\,dz\EMbr =\frac14\times2\pi i\times\left.\frac{d}{dz}\big(z^2\big)\right|_{z=1/2}\EMbr =\frac14\times2\pi i\times(2z)\Big|_{z=1/2}\EMbr =\frac\pi2\,i}
(\EMim{z=1/2} داخل دایرهٔ واحد است.)

\end{EMexample}

\begin{EMexample}{}

مطلوب است محاسبهٔ انتگرال زیر. (\EMim{C} دایرهٔ واحد است.)
\EMdisp{I\EMbr =\oint_C\frac{e^{-z}\sin z}{z^2}\,dz}

\EMdisp{f(z)\EMbr =e^{-z}\sin z,\quad\EMbr  z_0\EMbr =0,\qquad\EMbr  f'(z)\EMbr =-e^{-z}\sin z+e^{-z}\cos z,\qquad\EMbr  f'(z_0)\EMbr =f'(0)\EMbr =1}
\EMdisp{I\EMbr =2\pi i\,f'(z_0)\EMbr =2\pi i}

\end{EMexample}

\begin{EMexercise}{مسئله}

هرگاه \EMim{f(z)} در تمام صفحهٔ \EMim{z} تحلیلی و در تمام صفحه کران‌دار باشد، یعنی \EMim{|f(z)|<M}، ثابت کنید: \EMim{f(z)=k}، \EMim{k\in\mathbb{C}}.

مسیر \EMim{C} را به صورت دایره‌ای به شعاع \EMim{R} و به مرکز \EMim{z_0} انتخاب می‌کنیم و \EMim{R} را به سمت بی‌نهایت میل می‌دهیم.
\EMdisp{f'(z_0)\EMbr =\frac1{2\pi i}\oint_C\frac{f(z)}{(z-z_0)^2}\,dz}
\EMdisp{|f'(z_0)|\EMbr =\left|\frac1{2\pi i}\oint_C\frac{f(z)}{(z-z_0)^2}\,dz\right|\le\frac1{2\pi}\oint_C\frac{|f(z)|}{|z-z_0|^2}\,|dz|\le\frac1{2\pi}\,\frac{M}{R^2}\oint_C|dz|\EMbr =\frac1{2\pi}\,\frac{M}{R^2}\cdot2\pi R\EMbr =\frac MR\to0\quad\EMbr (R\to\infty)}
\EMdisp{\Longrightarrow\;|f'(z_0)|\EMbr =0\;\EMbr \Longrightarrow\;f'(z)\EMbr =0\;\EMbr \Longrightarrow\;f(z)\EMbr =k}

\end{EMexercise}

\EMfigure{m05-liouville-circle}{دایرهٔ \EMim{C} به مرکز \EMim{z_0} و شعاع \EMim{R}}

\begin{EMexercise}{}

مطلوب است محاسبهٔ انتگرال‌های زیر:

\begin{enumerate}
\def\labelenumi{\arabic{enumi}.}
\tightlist
\item
  \EMim{\displaystyle\oint_C\frac{z^2}{(2z-1)^4}\,dz}، \EMim{C}: دایرهٔ واحد
\item
  \EMim{\displaystyle\int_{1+i}^{1+3i}e^z\sin z\,dz}
\item
  \EMim{\displaystyle\int_{1+i}^{2+2i}\operatorname{Re}\{z^2\}\,dz} روی مسیر خط راست
\item
  \EMim{\displaystyle\oint_C\frac{e^{\cos z}\sin z}{1+z^6}\,dz} روی مسیر زیر
\end{enumerate}

\end{EMexercise}

\EMfigure{m05-exercise-circle}{مسیر تمرین ۴: دایره‌ای به مرکز \EMim{3+3i} و شعاع \EMim{r=1}}

\EMsetmodule{6}{سری لورنت، نقاط تکین و مانده‌ها}{Laurent Series, Singularities and Residues}
\EMchapterpage{6}{6}{سری لورنت، نقاط تکین و مانده‌ها}{Laurent Series, Singularities and Residues}
\begin{EMminitoc}
\EMtocitem{1}{سری لورنت}
\EMtocitem{2}{مانده}
\EMtocitem{3}{قضیهٔ مانده‌ها}
\EMtocitem{4}{مثال‌هایی از بسط تیلور و لورنت}
\EMtocitem{5}{مثال‌هایی از نقاط تکین و مانده}
\end{EMminitoc}
\EMchapterend
\EMhold

\EMsection{سری لورنت}{سری لورنت}

\begin{EMtheorem}{سری لورنت}

اگر تابع \EMim{f(z)} تابعی تحلیلی در ناحیهٔ \EMim{D} باشد، سری لورنت این تابع حول نقطهٔ \EMim{z_0} به صورت زیر نوشته می‌شود:
\EMdisp{f(z)\EMbr =\sum_{n=-\infty}^{\infty}a_n\,(z-z_0)^n,\qquad\EMbr  a_n\EMbr =\frac1{2\pi i}\oint_C\frac{f(v)}{(v-z_0)^{n+1}}\,dv}
که در آن \EMim{C} دایره‌ای به مرکز \EMim{z_0} است و تابع در همهٔ نقاط داخل و روی این دایره، مگر استثنائاً در \EMim{z_0} تحلیلی است.

شعاع همگرائی، فاصلهٔ \EMim{z_0} تا نزدیک‌ترین نقطهٔ غیرتحلیلی تابع \EMim{f(z)} است.

\end{EMtheorem}

برای \EMim{n} های نامنفی داریم:
\EMdisp{a_n\EMbr =\frac1{2\pi i}\oint_C\frac{f(v)}{(v-z_0)^{n+1}}\,dv\EMbr =\frac{f^{(n)}(z_0)}{n!},\qquad\EMbr  n\EMbr =0,1,2,\dots}
بنابراین خواهیم داشت:
\EMdisp{\begin{aligned}
f(z)&=\sum_{n=0}^{\infty}\frac{f^{(n)}(z_0)}{n!}(z-z_0)^n+\sum_{n=-\infty}^{-1}a_n(z-z_0)^n\\
&=\sum_{n=0}^{\infty}\frac{f^{(n)}(z_0)}{n!}(z-z_0)^n+\sum_{n=1}^{\infty}a_{-n}(z-z_0)^{-n}\\
&=\sum_{n=0}^{\infty}\frac{f^{(n)}(z_0)}{n!}(z-z_0)^n+\sum_{n=1}^{\infty}\frac{A_n}{(z-z_0)^n},\qquad A_n=a_{-n}
\end{aligned}}
\EMdisp{A_n\EMbr =a_{-n}\EMbr =\frac1{2\pi i}\oint_C\frac{f(v)}{(v-z_0)^{-n+1}}\,dv\EMbr =\frac1{2\pi i}\oint_C(v-z_0)^{n-1}f(v)\,dv,\qquad\EMbr  n\EMbr =1,2,\dots}

\begin{EMimportant}{سری لورنت: بخش تیلور و بخش تکین}

\EMdisp{f(z)\EMbr =\underbrace{\sum_{n=0}^{\infty}\frac{f^{(n)}(z_0)}{n!}(z-z_0)^n}_{\EMt{بخش سری تیلور}}+\underbrace{\sum_{n=1}^{\infty}\frac{A_n}{(z-z_0)^n}}_{\EMt{بخش تکین}},\qquad\EMbr  A_n\EMbr =\frac1{2\pi i}\oint_C(v-z_0)^{n-1}f(v)\,dv}

\end{EMimportant}

در صورتی که \EMim{z_0} نقطه‌ای غیر عادی (تکین) باشد، به عبارت دیگر تابع در این نقطه تحلیلی نباشد، بخش سری تیلور برای بیان تابع به تنهائی کافی نیست.

\EMsubsection{حالت‌های مختلف نقطهٔ \EMim{z_0}}{حالت‌های مختلف نقطهٔ z0}

\begin{enumerate}
\def\labelenumi{\arabic{enumi}.}
\tightlist
\item
  اگر تمامی \EMim{A_n} ها صفر باشند، نقطهٔ \EMim{z_0} یک نقطهٔ \textbf{معمولی} (\lr{Regular}) نامیده می‌شود. سری لورنت در این حالت همان سری تیلور خواهد بود.
\item
  اگر \EMim{A_p\neq0} ولی \EMim{(A_n=0\ \ \forall n>p)}، آنگاه \EMim{z_0} را \textbf{قطب} (\lr{pole}) \textbf{از مرتبهٔ \EMim{p}} می‌نامند.
\item
  اگر \EMim{z_0} نقطهٔ عادی نبوده و ضمناً قطب نیز نباشد، (به عبارت دیگر مقداری مانند \EMim{p} یافت نشود به قسمی که \EMim{A_p\neq0} ولی \EMim{(A_n=0\ \ \forall n>p)})، آنگاه \EMim{z_0} را نقطهٔ \textbf{تکین اصلی} (\lr{Essential Singularity}) می‌نامند.
\end{enumerate}

\EMhold

\EMsection{مانده}{مانده}

\begin{EMdefinition}{مانده (\lr{Residue})}

مانده تابع \EMim{f(z)} در نقطهٔ \EMim{z_0}: \EMim{A_1}.
\EMdisp{\operatorname{Res}\{f(z)\}\Big|_{z=z_0}\EMbr =A_1\EMbr =\frac1{2\pi i}\oint_Cf(v)\,dv}

\end{EMdefinition}

\begin{EMtheorem}{}

اگر تابع \EMim{f(z)} در نقطهٔ \EMim{z_0} دارای قطب از مرتبهٔ \EMim{p} باشد داریم:
\EMdisp{A_k\EMbr =\frac1{(p-k)!}\left.\frac{d^{\,p-k}}{dz^{\,p-k}}\Big[(z-z_0)^p\,f(z)\Big]\right|_{z=z_0},\qquad\EMbr  k\EMbr =1,2,\dots,p}

\end{EMtheorem}

\begin{EMproof}

\EMdisp{f(z)\EMbr =\sum_{n=1}^{\infty}\frac{A_n}{(z-z_0)^n}+\sum_{n=0}^{\infty}\frac{f^{(n)}(z_0)}{n!}(z-z_0)^n}
اگر تابع \EMim{f(z)} در نقطهٔ \EMim{z_0} دارای قطب از مرتبهٔ \EMim{p} باشد داریم:
\EMdisp{f(z)\EMbr =\frac{A_p}{(z-z_0)^p}+\frac{A_{p-1}}{(z-z_0)^{p-1}}+\cdots+\frac{A_2}{(z-z_0)^2}+\frac{A_1}{z-z_0}+\sum_{n=0}^{\infty}\frac{f^{(n)}(z_0)}{n!}(z-z_0)^n}
اگر طرفین را در \EMim{(z-z_0)^p} ضرب کنیم داریم:
\EMdisp{(z-z_0)^pf(z)\EMbr =A_p+A_{p-1}(z-z_0)+\cdots+A_2(z-z_0)^{p-2}+A_1(z-z_0)^{p-1}+\sum_{n=0}^{\infty}\frac{f^{(n)}(z_0)}{n!}(z-z_0)^{n+p}}
حال اگر به \EMim{z} مقدار \EMim{z_0} بدهیم، داریم:
\EMdisp{(z-z_0)^pf(z)\Big|_{z=z_0}\EMbr =A_p}
اگر از طرفین \EMim{(z-z_0)^pf(z)} یک بار مشتق بگیریم و به \EMim{z} مقدار \EMim{z_0} بدهیم، داریم:
\EMdisp{\begin{aligned}
\frac{d}{dz}\Big[(z-z_0)^pf(z)\Big]\Big|_{z=z_0}&=\Bigg\{A_{p-1}+\cdots+A_2(p-2)(z-z_0)^{p-3}+A_1(p-1)(z-z_0)^{p-2}\\
&\qquad+\sum_{n=0}^{\infty}\frac{f^{(n)}(z_0)}{n!}(z-z_0)^{n+p-1}(n+p)\Bigg\}\Bigg|_{z=z_0}=A_{p-1}
\end{aligned}}
اگر از طرفین \EMim{(z-z_0)^pf(z)} دو بار مشتق بگیریم و به \EMim{z} مقدار \EMim{z_0} بدهیم، داریم:
\EMdisp{\begin{aligned}
\frac{d^2}{dz^2}\Big[(z-z_0)^pf(z)\Big]\Big|_{z=z_0}&=\Bigg\{2A_{p-2}+\cdots+A_2(p-2)(p-3)(z-z_0)^{p-4}+A_1(p-1)(p-2)(z-z_0)^{p-3}\\
&\qquad+\sum_{n=0}^{\infty}\frac{f^{(n)}(z_0)}{n!}(z-z_0)^{n+p-2}(n+p)(n+p-1)\Bigg\}\Bigg|_{z=z_0}=2A_{p-2}
\end{aligned}}
\EMdisp{\Longrightarrow\quad\EMbr  A_{p-2}\EMbr =\frac1{2!}\left.\frac{d^2}{dz^2}\Big[(z-z_0)^pf(z)\Big]\right|_{z=z_0}}

\end{EMproof}

\begin{EMexercise}{}

به همین ترتیب مشتق‌گیری را ادامه دهید و اثبات قضیه را کامل کنید.

\end{EMexercise}

\begin{EMremark}{نتیجه}

در حالت خاص، در صورتی که تابع \EMim{f(z)} در نقطهٔ \EMim{z_0} دارای قطب از مرتبهٔ \EMim{p} باشد، مقدار مانده از رابطهٔ زیر به دست می‌آید:
\EMdisp{\operatorname{Res}\{f(z)\}\EMbr =A_1\EMbr =\frac1{(p-1)!}\left.\frac{d^{\,p-1}}{dz^{\,p-1}}\Big[(z-z_0)^p\,f(z)\Big]\right|_{z=z_0}}

\end{EMremark}

\EMhold

\EMsection{قضیهٔ مانده‌ها}{قضیهٔ مانده‌ها}

\begin{EMtheorem}{قضیهٔ مانده‌ها}

اگر تابع \EMim{f(z)} در همه‌جای داخل مسیر \EMim{C} به‌جز نقاط \EMim{z_1,z_2,\dots,z_n} تحلیلی باشد و مانده تابع در این نقاط به ترتیب \EMim{A_1^1,A_1^2,\dots,A_1^n} باشد، آنگاه:
\EMdisp{\oint_Cf(z)\,dz\EMbr =2\pi i\sum_{k=1}^{n}A_1^k}

\end{EMtheorem}

\EMfigure{m06-residue-theorem}{مسیر \EMim{C} و نقاط تکین \EMim{z_1,\dots,z_n} در داخل آن؛ حول هر نقطه
دایره‌ای کوچک رسم شده است}

اثبات به عنوان تمرین بر عهدهٔ خواننده است.

\EMhold

\EMsection{مثال‌هایی از بسط تیلور و لورنت}{مثال‌هایی از بسط تیلور و لورنت}

\begin{EMexample}{}

مطلوب است محاسبهٔ سری لورنت تابع \EMim{f(z)=\dfrac1{z^4+1}} در همسایگی نقطهٔ \EMim{z_0=0}.

دیده می‌شود که تابع \EMim{f(z)} در نقطهٔ \EMim{z_0=0} تحلیلی است و لذا این نقطه یک نقطهٔ معمولی است و سری لورنت همان سری تیلور (مک‌لورن) خواهد بود. لذا داریم:
\EMdisp{f(z)\EMbr =f(0)+\frac1{1!}f^{(1)}(0)z+\frac1{2!}f^{(2)}(0)z^2+\frac1{3!}f^{(3)}(0)z^3+\cdots}
سری هندسی:
\EMdisp{1+q+q^2+q^3+q^4+\cdots\EMbr =\frac1{1-q},\qquad\EMbr  |q|<1}
\EMdisp{\frac1{1+z^4}\EMbr =1-z^4+z^8-z^{12}+-\cdots,\qquad\EMbr  |z^4|<1\ \EMbr \Longrightarrow\ |z|<1}
(سری تیلور/مک‌لورن؛ ناحیهٔ همگرائی با شعاع ۱.)
\EMdisp{z^4+1\EMbr =0\;\EMbr \Longrightarrow\;z_1\EMbr =e^{i\pi/4},\quad\EMbr  z_2\EMbr =e^{i3\pi/4},\quad\EMbr  z_3\EMbr =e^{i5\pi/4},\quad\EMbr  z_4\EMbr =e^{i7\pi/4}}
مشاهده می‌شود که فاصلهٔ \EMim{z_0=0} تا نزدیک‌ترین نقطهٔ غیرتحلیلی تابع (شعاع همگرائی) عدد ۱ می‌باشد.

\end{EMexample}

\EMfigure{m06-radius-1}{چهار قطب \EMim{z^4+1=0} روی دایرهٔ واحد؛ شعاع همگرائی سری حول مبدأ برابر
۱ است}

\begin{EMexample}{}

مطلوب است محاسبهٔ سری لورنت تابع \EMim{f(z)=\dfrac{4z^2+30z+68}{(z+4)^2(z-2)}} در همسایگی نقطهٔ \EMim{z_0=0}.

دیده می‌شود که تابع \EMim{f(z)} در نقطهٔ \EMim{z_0=0} تحلیلی است و لذا این نقطه یک نقطهٔ معمولی است و سری لورنت همان سری تیلور (مک‌لورن) خواهد بود. لذا داریم:
\EMdisp{f(z)\EMbr =f(0)+\frac1{1!}f^{(1)}(0)z+\frac1{2!}f^{(2)}(0)z^2+\frac1{3!}f^{(3)}(0)z^3+\cdots}
\EMdisp{\frac1{z-2}\EMbr =\frac{-1}{2-z}\EMbr =\frac{-1}2\,\frac1{1-z/2}\EMbr =\frac{-1}2\left(1+\frac z2+\frac{z^2}4+\frac{z^3}8+\cdots\right),\qquad\EMbr  \left|\frac z2\right|<1\ \EMbr \Longrightarrow\ |z|<2}
\EMdisp{\frac1{4+z}\EMbr =\frac14\,\frac1{1+z/4}\EMbr =\frac14\left(1-\frac z4+\frac{z^2}{16}-\frac{z^3}{64}+\cdots\right),\qquad\EMbr  \left|-\frac z4\right|<1\ \EMbr \Longrightarrow\ |z|<4}
\EMdisp{\frac d{dz}\left(\frac1{4+z}\right)\EMbr =\frac{-1}{(4+z)^2}\EMbr =\frac14\left(-\frac14+\frac{2z}{16}-\frac{3z^2}{64}+\cdots\right)\;\EMbr \Longrightarrow\;\frac1{(4+z)^2}\EMbr =\frac14\left(\frac14-\frac{2z}{16}+\frac{3z^2}{64}-\cdots\right)}
\EMdisp{f(z)\EMbr =(4z^2+30z+68)\times\frac1{(z+4)^2}\times\frac1{z-2}}
\EMdisp{f(z)\EMbr =(4z^2+30z+68)\times\left(\frac{-1}2\right)\left(1+\frac z2+\frac{z^2}4+\frac{z^3}8+\cdots\right)\times\frac14\left(\frac14-\frac z8+\frac{3z^2}{64}-\cdots\right)}
\EMdisp{f(z)\EMbr =\frac{-1}{16}\left(34+15z+\frac{67}8z^2+\cdots\right),\qquad\EMbr  |z|<2}
(ناحیهٔ همگرائی با شعاع ۲.)

\end{EMexample}

\EMfigure{m06-radius-2}{قطب‌های \EMim{z=-4} و \EMim{z=2}؛ دایرهٔ همگرائی حول مبدأ شعاع ۲ دارد}

\begin{EMexample}{}

مطلوب است محاسبهٔ سری لورنت تابع \EMim{f(z)=\cos z^2} در همسایگی نقطهٔ \EMim{z_0=0}.

با توجه به تعریف تابع کسینوس، داریم:
\EMdisp{\cos z\EMbr =1-\frac{z^2}{2!}+\frac{z^4}{4!}-\frac{z^6}{6!}+-\cdots}
\EMdisp{f(z)\EMbr =\cos z^2\EMbr =1-\frac{z^4}{2!}+\frac{z^8}{4!}-\frac{z^{12}}{6!}+-\cdots}
ناحیهٔ همگرائی تمام صفحهٔ مختلط با شعاع همگرائی بی‌نهایت.

\end{EMexample}

\begin{EMexample}{}

مطلوب است محاسبهٔ سری لورنت تابع \EMim{f(z)=\dfrac{\cos z}{1-z^2}} در همسایگی نقطهٔ \EMim{z_0=0}.

\EMdisp{\left\{\begin{aligned}
\frac1{1-z^2}&=1+z^2+z^4+z^6+\cdots&&|z^2|<1\ \ \ |z|<1\\
\cos z&=1-\frac{z^2}{2!}+\frac{z^4}{4!}-\frac{z^6}{6!}+-\cdots&&|z|<\infty
\end{aligned}\right.}
\EMdisp{\begin{aligned}
f(z)=\frac{\cos z}{1-z^2}&=\big(1+z^2+z^4+z^6+\cdots\big)\left(1-\frac{z^2}{2!}+\frac{z^4}{4!}-\frac{z^6}{6!}+-\cdots\right)\\
&=1+\left(1-\frac1{2!}\right)z^2+\left(1+\frac1{4!}-\frac1{2!}\right)z^4+\cdots,\qquad |z|<1
\end{aligned}}

\end{EMexample}

\begin{EMexample}{}

مطلوب است محاسبهٔ سری لورنت تابع \EMim{f(z)=\dfrac1{z^4}} در همسایگی نقطهٔ \EMim{z_0=1}.

\EMdisp{g(z)\EMbr =\frac1z\EMbr =\frac1{z-1+1}\EMbr =\frac1{1+(z-1)}\EMbr =1-(z-1)+(z-1)^2-(z-1)^3+-\cdots,\qquad\EMbr  |z-1|<1}
\EMdisp{\frac d{dz}g(z)\EMbr =\frac{-1}{z^2}\EMbr =-1+2(z-1)-3(z-1)^2+4(z-1)^3-+\cdots,\qquad\EMbr  |z-1|<1}
\EMdisp{\frac{d^2}{dz^2}g(z)\EMbr =\frac2{z^3}\EMbr =2-6(z-1)+12(z-1)^2-20(z-1)^3+-\cdots,\qquad\EMbr  |z-1|<1}
\EMdisp{\frac{d^3}{dz^3}g(z)\EMbr =\frac{-6}{z^4}\EMbr =-6+24(z-1)-60(z-1)^2-+\cdots,\qquad\EMbr  |z-1|<1}
\EMdisp{f(z)\EMbr =\frac1{z^4}\EMbr =1-4(z-1)+10(z-1)^2-+\cdots}

\end{EMexample}

\EMfigure{m06-radius-3}{ناحیهٔ همگرائی سری حول \EMim{z_0=1} دایره‌ای به شعاع ۱ است (تا قطب
\EMim{z=0})}

\begin{EMexercise}{}

مطلوب است بسط لورنت توابع زیر در همسایگی \EMim{z_0=0}:

\begin{enumerate}
\def\labelenumi{\arabic{enumi}.}
\tightlist
\item
  \EMim{f_1(z)=\dfrac1{\sqrt{1-z^2}}}
\item
  \EMim{f_2(z)=\sin^{-1}z}
\item
  \EMim{f_3(z)=\dfrac{e^z}{z^2(z^2+1)}}
\item
  تابع \EMim{\sinh z} را بر حسب توانهائی از \EMim{(z-\pi i)} بسط داده، ثابت کنید: \EMim{\displaystyle\lim_{z\to\pi i}\frac{\sinh z}{z-\pi i}=-1}
\item
  ثابت کنید: \EMim{\displaystyle z\cosh z^2=z+\sum_{n=1}^{\infty}\frac1{(2n)!}\,z^{4n+1}}
\end{enumerate}

\end{EMexercise}

\EMhold

\EMsection{مثال‌هایی از نقاط تکین و مانده}{مثال‌هایی از نقاط تکین و مانده}

\begin{EMexample}{}

مطلوب است محاسبهٔ سری لورنت تابع \EMim{f(z)=\dfrac{z}{(z+1)(z+4)^3}} در همسایگی نقاط غیرتحلیلی آن و محاسبهٔ مانده در آن نقاط.

تابع در \EMim{z=-1} و \EMim{z=-4} تحلیلی نیست.

\textbf{الف)} در همسایگی \EMim{z=-1} دایره‌ای به مرکز این نقطه و به شعاع کمتر از ۳ (\EMim{|z+1|<3}) نقطهٔ تکین (ایزولهٔ) \EMim{z=-4} را در بر ندارد.
\EMdisp{\frac{z}{(z+4)^3}\EMbr =\frac{z+4-4}{(z+4)^3}\EMbr =\frac{z+4}{(z+4)^3}+\frac{-4}{(z+4)^3}\EMbr =\frac1{(z+4)^2}-\frac4{(z+4)^3}}
\EMdisp{\frac1{z+4}\EMbr =\frac1{(z+1)+3}\EMbr =\frac13\,\frac1{1+\frac{z+1}3}\EMbr =\frac13\left[1-\frac{z+1}3+\left(\frac{z+1}3\right)^2-\left(\frac{z+1}3\right)^3+-\cdots\right]}
\EMdisp{\frac{-1}{(z+4)^2}\EMbr =\frac13\left[-\frac13+\frac2{3^2}(z+1)-\frac3{3^3}(z+1)^2+-\cdots\right]\;\EMbr \Longrightarrow\;\frac1{(z+4)^2}\EMbr =\frac19-\frac2{27}(z+1)+\frac1{27}(z+1)^2-+\cdots}
\EMdisp{\frac1{(z+4)^3}\EMbr =\frac16\left[\frac2{3^2}-\frac{2\times3}{3^3}(z+1)+\frac{3\times4}{3^4}(z+1)^2-\frac{4\times5}{3^5}(z+1)^3+-\cdots\right]\EMbr =\frac1{27}-\frac1{27}(z+1)+\frac2{81}(z+1)^2-+\cdots}
\EMdisp{\frac z{(z+4)^3}\EMbr =\frac1{(z+4)^2}-\frac4{(z+4)^3}\EMbr =-\frac1{27}+\frac2{27}(z+1)-\frac5{81}(z+1)^2+\cdots}
\EMdisp{f(z)\EMbr =\frac1{z+1}\,\frac z{(z+4)^3}\EMbr =-\frac1{27}\,\frac1{z+1}+\frac2{27}-\frac5{81}(z+1)+\cdots}
نوع نقطهٔ تکین: قطب ساده (از مرتبهٔ ۱). مانده: \EMim{\operatorname{Res}\{f(z)\}\big|_{z=-1}=-\dfrac1{27}}.

\textbf{ب)} در همسایگی \EMim{z=-4} دایره‌ای به مرکز این نقطه و به شعاع کمتر از ۳ (\EMim{|z+4|<3}) نقطهٔ تکین (ایزولهٔ) \EMim{z=-1} را در بر ندارد.
\EMdisp{\frac z{z+1}\EMbr =\frac{z+1-1}{z+1}\EMbr =\frac{z+1}{z+1}-\frac1{z+1}\EMbr =1-\frac1{z+1}}
\EMdisp{\frac1{z+1}\EMbr =\frac1{z+4-3}\EMbr =\frac1{-3+(z+4)}\EMbr =\frac{-1}3\,\frac1{1-\frac{z+4}3}\EMbr =\frac{-1}3\left[1+\frac{z+4}3+\frac{(z+4)^2}{3^2}+\cdots\right],\qquad\EMbr  |z+4|<3}
\EMdisp{\frac z{z+1}\EMbr =1+\frac13+\frac1{3^2}(z+4)+\frac1{3^3}(z+4)^2+\frac1{3^4}(z+4)^3+\cdots}
\EMdisp{f(z)\EMbr =\frac z{(z+1)(z+4)^3}\EMbr =\frac1{(z+4)^3}\,\frac z{z+1}\EMbr =\frac{4/3}{(z+4)^3}+\frac1{3^2}\frac1{(z+4)^2}+\frac1{3^3}\frac1{z+4}+\frac1{3^4}+\frac1{3^5}(z+4)+\frac1{3^6}(z+4)^2+\cdots}
نوع نقطهٔ تکین: قطب از مرتبهٔ ۳. مانده: \EMim{\operatorname{Res}\{f(z)\}\big|_{z=-4}=\dfrac1{27}}.

\end{EMexample}

\EMfigure{m06-poles-1}{قطب‌های \EMim{z=-1} و \EMim{z=-4} تابع \EMim{f(z)=\frac{z}{(z+1)(z+4)^3}}}

\begin{EMexample}{}

مطلوب است محاسبهٔ سری لورنت تابع \EMim{f(z)=e^{1/z}} در همسایگی نقطهٔ غیرتحلیلی آن و محاسبهٔ مانده در آن نقطه.

\EMdisp{e^z\EMbr =1+z+\frac{z^2}{2!}+\frac{z^3}{3!}+\cdots\;\EMbr \Longrightarrow\;e^{1/z}\EMbr =1+\frac1z+\frac1{2!\,z^2}+\frac1{3!\,z^3}+\cdots}
نوع نقطهٔ تکین: تکین اصلی (بی‌نهایت جمله با توان‌های منفی). مانده: \EMim{\operatorname{Res}\{f(z)\}\big|_{z=0}=1}.

\end{EMexample}

\begin{EMexample}{}

مطلوب است محاسبهٔ سری لورنت تابع \EMim{f(z)=\cos\!\left(\dfrac z{z-1}\right)} در همسایگی نقطهٔ غیرتحلیلی آن و محاسبهٔ مانده در آن نقطه.

\EMdisp{\frac z{z-1}\EMbr =1+\frac1{z-1}\;\EMbr \Longrightarrow\;\cos\!\left(\frac z{z-1}\right)\EMbr =\cos\!\left(1+\frac1{z-1}\right)\EMbr =\cos1\cos\!\left(\frac1{z-1}\right)-\sin1\sin\!\left(\frac1{z-1}\right)}
\EMdisp{\cos z\EMbr =1-\frac{z^2}{2!}+\frac{z^4}{4!}-\frac{z^6}{6!}+-\cdots,\qquad\EMbr  \sin z\EMbr =z-\frac{z^3}{3!}+\frac{z^5}{5!}-+\cdots}
\EMdisp{\begin{aligned}
\cos\!\left(\frac z{z-1}\right)&=\cos1\left(1-\frac1{2!}\frac1{(z-1)^2}+\frac1{4!}\frac1{(z-1)^4}-+\cdots\right)\\
&\quad-\sin1\left(\frac1{z-1}-\frac1{3!}\frac1{(z-1)^3}+\frac1{5!}\frac1{(z-1)^5}-+\cdots\right)\\
&=\cos1-\sin1\,\frac1{z-1}-\cos1\,\frac1{2!}\frac1{(z-1)^2}+\sin1\,\frac1{3!}\frac1{(z-1)^3}+\cdots
\end{aligned}}
نوع نقطهٔ تکین: تکین اصلی. مانده: \EMim{\operatorname{Res}\{f(z)\}\big|_{z=1}=-\sin1}.

\end{EMexample}

\begin{EMtheorem}{مانده در قطب ساده}

ثابت کنید هرگاه تابع \EMim{f(z)=q(z)/g(z)} بوده و \EMim{g(z)} در نقطهٔ \EMim{z_0} ریشهٔ مرتبهٔ اول داشته باشد، آنگاه
\EMdisp{\operatorname{Res}\{f(z)\}\Big|_{z=z_0}\EMbr =\frac{q(z_0)}{g'(z_0)}}

\end{EMtheorem}

\begin{EMproof}

چون نقطهٔ \EMim{z_0} قطب مرتبهٔ ۱ است داریم (\EMim{p=1}):
\EMdisp{\begin{aligned}
\operatorname{Res}\{f(z)\}\Big|_{z=z_0}=A_1&=\frac1{(p-1)!}\left.\frac{d^{\,p-1}}{dz^{\,p-1}}\Big[(z-z_0)^pf(z)\Big]\right|_{z=z_0}=(z-z_0)\frac{q(z)}{g(z)}\Bigg|_{z=z_0}\\
&=(z-z_0)\frac{q(z)}{g(z)-g(z_0)}\Bigg|_{z=z_0}=\frac{q(z)}{\dfrac{g(z)-g(z_0)}{z-z_0}}\Bigg|_{z=z_0}=\frac{q(z_0)}{g'(z_0)}
\end{aligned}}

\end{EMproof}

\begin{EMexercise}{}

سری لورنت توابع زیر را حول \EMim{z=0} به دست آورده، نوع این نقطه را تعیین کنید.
\EMdisp{f_1(z)\EMbr =\frac z{e^z-1},\qquad\EMbr  f_2(z)\EMbr =\frac{e^{-z}}{z^3},\qquad\EMbr  f_3(z)\EMbr =\frac1{z^2(1-z^2)},\qquad\EMbr  f_4(z)\EMbr =\frac{\sin z^2}{z^3}}

\end{EMexercise}

\EMsetmodule{7}{محاسبهٔ انتگرال‌های معین به روش مانده‌ها}{Evaluating Definite Integrals by Residues}
\EMchapterpage{7}{7}{محاسبهٔ انتگرال‌های معین به روش مانده‌ها}{Evaluating Definite Integrals by Residues}
\begin{EMminitoc}
\EMtocitem{1}{انتگرال‌های دستهٔ اول}
\EMtocitem{2}{انتگرال‌های دستهٔ دوم}
\end{EMminitoc}
\EMchapterend
\begin{EMtheorem}{قضیهٔ مانده‌ها}

اگر تابع \EMim{f(z)} در همه‌جای داخل مسیر \EMim{C} به‌جز نقاط \EMim{z_1,z_2,\dots,z_n} تحلیلی باشد و مانده تابع در این نقاط به ترتیب \EMim{A_1^1,A_1^2,\dots,A_1^n} باشد، آنگاه:
\EMdisp{\oint_Cf(z)\,dz\EMbr =2\pi i\sum_{k=1}^{n}A_1^k}

\end{EMtheorem}

دسته انتگرال‌های مختلفی توسط روش مانده‌ها قابل حل هستند. در اینجا به دو دستهٔ آنها اشاره می‌شود:

\begin{EMdefinition}{انتگرال‌های دستهٔ اول}

در این نوع انتگرال‌ها، مخرج انتگرند به صورت چندجمله‌ای بوده، صورت به فرم تابعی از نوع چندجمله‌ای یا نمائی و یا توابع سینوس یا کسینوس (که آنها هم به صورت نمائی قابل بیان هستند) می‌باشد.

\end{EMdefinition}

\begin{EMdefinition}{انتگرال‌های دستهٔ دوم}

در این نوع انتگرال‌ها، انتگرند تابعی از سینوس یا کسینوس می‌باشد.

\end{EMdefinition}

\EMhold

\EMsection{انتگرال‌های دستهٔ اول}{انتگرال‌های دستهٔ اول}

\begin{EMexample}{}

مطلوب است محاسبهٔ انتگرال \EMim{\displaystyle\int_0^\infty\frac{\cos\omega x}{(1+x^2)^2}\,dx} (\EMim{\omega>0}).

تابع \EMim{f(z)=e^{i\omega z}\big/(1+z^2)^2} را در نظر می‌گیریم. این تابع در \EMim{z=i} و \EMim{z=-i} تحلیلی نیست. مسیر \EMim{C} نیم‌دایرهٔ بالایی به شعاع \EMim{R} است.
\EMdisp{\begin{aligned}
\oint_C\frac{e^{i\omega z}}{(1+z^2)^2}\,dz&=2\pi i\times\operatorname{Res}\left\{\frac{e^{i\omega z}}{(1+z^2)^2}\right\}\Bigg|_{z=i}\\
&=2\pi i\times\left[\frac d{dz}\left\{(z-i)^2\,\frac{e^{i\omega z}}{(z-i)^2(z+i)^2}\right\}\right]_{z=i}\\
&=2\pi i\times\left[\frac{i\omega e^{i\omega z}(z+i)-2e^{i\omega z}}{(z+i)^3}\right]_{z=i}\\
&=2\pi i\times e^{-\omega}\left(\frac{i\omega(2i)-2}{(2i)^3}\right)\\
&=2\pi i\times e^{-\omega}\,\frac{\omega+1}{4i}=\frac\pi2(\omega+1)e^{-\omega}
\end{aligned}}
از طرف دیگر
\EMdisp{\oint_C\frac{e^{i\omega z}}{(1+z^2)^2}\,dz\EMbr =\int_{-R}^{R}\frac{e^{i\omega x}}{(1+x^2)^2}\,dx+\int_0^\pi\frac{e^{i\omega Re^{it}}\,Rie^{it}}{(1+R^2e^{i2t})^2}\,dt}
و وقتی \EMim{R\to\infty} انتگرال روی نیم‌دایره صفر می‌شود و بخش موهومی انتگرال روی محور حقیقی (تابع فرد \EMim{\sin\omega x/(1+x^2)^2}) صفر است:
\EMdisp{\oint_C\frac{e^{i\omega z}}{(1+z^2)^2}\,dz\ \xrightarrow{\ R\to\infty\ }\ 2\int_0^\infty\frac{\cos\omega x}{(1+x^2)^2}\,dx}
\EMdisp{\int_0^\infty\frac{\cos\omega x}{(1+x^2)^2}\,dx\EMbr =\frac\pi4(\omega+1)e^{-\omega}}

\end{EMexample}

\EMfigure{m07-semicircle}{مسیر نیم‌دایرهٔ بالایی؛ قطب مرتبهٔ دوم \EMim{z=i} داخل مسیر و \EMim{z=-i}
بیرون آن است}

\EMsection{انتگرال‌های دستهٔ دوم}{انتگرال‌های دستهٔ دوم}

\EMdisp{\int_0^{2\pi}f(a\cos t,\,b\sin t)\,dt}
برای حل این‌گونه انتگرال‌ها فرض می‌کنیم \EMim{z=e^{it}}، در این صورت داریم:
\EMdisp{dz\EMbr =ie^{it}\,dt\;\EMbr \Longrightarrow\;dt\EMbr =\frac{dz}{ie^{it}}\EMbr =\frac{dz}{iz}}
\EMdisp{\sin t\EMbr =\frac{e^{it}-e^{-it}}{2i}\EMbr =\frac{z-z^{-1}}{2i},\qquad\EMbr  \cos t\EMbr =\frac{e^{it}+e^{-it}}2\EMbr =\frac{z+z^{-1}}2}

\begin{EMimportant}{}

\EMdisp{\int_0^{2\pi}f(a\cos t,\,b\sin t)\,dt\EMbr =\oint_Cf\!\left(a\,\frac{z+z^{-1}}{2},\;b\,\frac{z-z^{-1}}{2i}\right)\frac{dz}{iz}}

\end{EMimportant}

در اینجا \EMim{C} دایرهٔ واحد است و به سهولت با استفاده از روش مانده‌ها انتگرال حل می‌شود.

\EMfigure{m07-unit-circle}{دایرهٔ واحد \EMim{z=e^{it}}، \EMim{0\le t\le2\pi}}

\begin{EMexample}{}

مطلوب است محاسبهٔ انتگرال \EMim{\displaystyle\int_0^{2\pi}\frac1{2+\cos x}\,dx}.

\EMdisp{z\EMbr =e^{ix},\qquad\EMbr  dz\EMbr =ie^{ix}dx,\qquad\EMbr  dx\EMbr =\frac{dz}{iz},\qquad\EMbr  \cos x\EMbr =\frac{z+z^{-1}}2\EMbr =\frac{z^2+1}{2z}}
\EMdisp{2+\cos x\EMbr =2+\frac{z^2+1}{2z}\EMbr =\frac{z^2+4z+1}{2z}}
\EMdisp{\int_0^{2\pi}\frac1{2+\cos x}\,dx\EMbr =\oint_C\frac{2z}{z^2+4z+1}\,\frac{dz}{iz}\EMbr =\frac2i\oint_C\frac1{z^2+4z+1}\,dz\EMbr =\frac2i\times2\pi i\times\sum_k\operatorname{Res}\left\{\frac1{z^2+4z+1}\right\}\Bigg|_{z=z_k}}
\EMdisp{z^2+4z+1\EMbr =0\;\EMbr \Longrightarrow\;z_1\EMbr =-2+\sqrt3\EMbr \approx-0.27,\qquad\EMbr  z_2\EMbr =-2-\sqrt3\EMbr \approx-3.73}
فقط \EMim{z_1} داخل دایرهٔ واحد قرار دارد.
\EMdisp{\int_0^{2\pi}\frac1{2+\cos x}\,dx\EMbr =4\pi\times\operatorname{Res}\left\{\frac1{z^2+4z+1}\right\}\Bigg|_{z=z_1}\EMbr =4\pi\left[\big(z+2-\sqrt3\big)\frac1{\big(z+2-\sqrt3\big)\big(z+2+\sqrt3\big)}\right]_{z_1=-2+\sqrt3}\EMbr =\frac{2\pi}{\sqrt3}}

\end{EMexample}

\EMfigure{m07-poles-cos}{دایرهٔ واحد و قطب‌های \EMim{z_1\approx-0.27} (داخل) و
\EMim{z_2\approx-3.73} (بیرون)}

\begin{EMexercise}{کاربرد قضیهٔ مانده‌ها در حل انتگرال}

با استفاده از قضیهٔ مانده‌ها مقدار انتگرال‌های معین زیر را به دست آورید:
\EMdisp{I_1\EMbr =\int_0^{2\pi}\frac{dt}{1-2p\cos t+p^2},\qquad\EMbr  I_2\EMbr =\int_{-\infty}^{\infty}\frac{\sin2x}{x^2+x+1}\,dx}
\EMdisp{I_3\EMbr =\int_{-\infty}^{\infty}\frac{x}{(x^2-2x+2)^2}\,dx,\qquad\EMbr  I_4\EMbr =\int_{-\infty}^{\infty}\frac{\cos^4x}{(x^2+1)(x^2+4)}\,dx}

\end{EMexercise}

\EMpartpage{2}{۲}{آنالیز فوریه}{Fourier Analysis}{فصل‌های ۸ تا ۱۰}
\EMsetmodule{8}{سری فوریه}{Fourier Series}
\EMchapterpage{8}{8}{سری فوریه}{Fourier Series}
\begin{EMminitoc}
\EMtocitem{1}{توابع متناوب}
\EMtocitem{2}{یادآوری از مبحث مثلثات}
\EMtocitem{3}{یادآوری: انتگرال‌ها و تعامد}
\EMtocitem{4}{سری فوریه}
\EMtocitem{5}{یادآوری: انتگرال معین حاصل‌ضرب دو تابع}
\EMtocitem{6}{مثال‌های سری فوریه}
\EMtocitem{7}{سری فوریهٔ توابع متناوب با دورهٔ تناوب دلخواه}
\EMtocitem{8}{قضیهٔ دریشله}
\EMtocitem{9}{حل چند مسئله}
\end{EMminitoc}
\EMchapterend
\EMkeepfig{m08-periodic}{5}

\EMsection{توابع متناوب}{توابع متناوب}

\EMdisp{f(x+p)\EMbr =f(x)\quad\EMbr \EMbr \Longrightarrow\quad\EMbr  f(x+np)\EMbr =f(x)}

\EMfigure{m08-periodic}{تابعی متناوب با دورهٔ تناوب \EMim{p}}

\begin{itemize}
\tightlist
\item
  کوچک‌ترین مقدار مثبت \EMim{p} به قسمی که تابع با آن تناوبی باشد، دورهٔ تناوب اصلی تابع نامیده می‌شود.
\item
  اگر تابع با دورهٔ تناوب \EMim{p} تناوبی باشد آنگاه با دورهٔ \EMim{np}، که \EMim{n} یک عدد طبیعی است، نیز متناوب است.
\end{itemize}

\EMhold

\EMsubsection{توابع متناوب سینوس و کسینوس}{توابع متناوب سینوس و کسینوس}

\begin{EMexample}{}

تابع \EMim{\sin x} تابعی با دورهٔ تناوب \EMim{2\pi} تناوبی است و \EMim{\sin2x} با دورهٔ تناوب اصلی \EMim{\pi} تناوبی است هر چند \EMim{2\pi} نیز دورهٔ تناوب آن هست. همینطور تابع \EMim{\sin3x} با دورهٔ تناوب \EMim{2\pi/3} تناوبی است، هر چند \EMim{2\pi} نیز دورهٔ تناوب آن هست.

\end{EMexample}

\EMfigure{m08-sincos}{توابع \EMim{\sin nx} و \EMim{\cos nx} برای \EMim{n=1,2,3} در بازهٔ
\EMim{[0,2\pi]}}

\EMsection{یادآوری از مبحث مثلثات}{یادآوری از مبحث مثلثات}

\EMdisp{\begin{gathered}
\sin^2\alpha+\cos^2\alpha=1\\
\sin(\alpha+\beta)=\sin\alpha\cos\beta+\cos\alpha\sin\beta\\
\cos(\alpha+\beta)=\cos\alpha\cos\beta-\sin\alpha\sin\beta\\
\cos2\alpha=2\cos^2\alpha-1=1-2\sin^2\alpha=\cos^2\alpha-\sin^2\alpha\\
\sin2\alpha=2\sin\alpha\cos\alpha\\
\sin^2\alpha=\frac{1-\cos2\alpha}{2}\qquad\cos^2\alpha=\frac{1+\cos2\alpha}{2}\\
\sin\alpha\sin\beta=\tfrac12\big[\cos(\alpha-\beta)-\cos(\alpha+\beta)\big]\\
\cos\alpha\cos\beta=\tfrac12\big[\cos(\alpha-\beta)+\cos(\alpha+\beta)\big]\\
\sin\alpha\cos\beta=\tfrac12\big[\sin(\alpha-\beta)+\sin(\alpha+\beta)\big]
\end{gathered}}

\EMsection{یادآوری: انتگرال‌ها و تعامد}{یادآوری: انتگرال‌ها و تعامد}

\EMdisp{\int_{-\pi}^{\pi}\sin nt\,dt\EMbr =0,\qquad\EMbr  n\in\mathbb{Z}}
\EMdisp{\int_{-\pi}^{\pi}\cos nt\,dt\EMbr =0,\qquad\EMbr  n\in\mathbb{Z}\setminus\{0\}}

\begin{EMproof}

\EMdisp{\int_{-\pi}^{\pi}\sin nt\,dt\EMbr =-\frac1n\cos nt\Big|_{-\pi}^{\pi}\EMbr =0,\qquad\EMbr  \int_{-\pi}^{\pi}\cos nt\,dt\EMbr =\frac1n\sin nt\Big|_{-\pi}^{\pi}\EMbr =0\qquad\EMbr (n\neq0)}

\end{EMproof}

\EMsubsection{تعامد}{تعامد}

\EMdisp{\int_{-\pi}^{\pi}\sin nt\sin mt\,dt\EMbr =0,\qquad\EMbr  n,m\in\mathbb{Z},\ n\neq m}
\EMdisp{\int_{-\pi}^{\pi}\cos nt\cos mt\,dt\EMbr =0,\qquad\EMbr  n,m\in\mathbb{Z},\ n\neq m}
\EMdisp{\int_{-\pi}^{\pi}\sin nt\cos mt\,dt\EMbr =0,\qquad\EMbr  n,m\in\mathbb{Z}\ \ (n\neq m\ \EMt{یا}\ n\EMbr =m)}

\begin{EMproof}

\EMdisp{\int_{-\pi}^{\pi}\sin nx\sin mx\,dx\EMbr =\frac12\int_{-\pi}^{\pi}\cos(n-m)x\,dx-\frac12\int_{-\pi}^{\pi}\cos(n+m)x\,dx}
\EMdisp{\int_{-\pi}^{\pi}\cos nx\cos mx\,dx\EMbr =\frac12\int_{-\pi}^{\pi}\cos(n+m)x\,dx+\frac12\int_{-\pi}^{\pi}\cos(n-m)x\,dx}
\EMdisp{\int_{-\pi}^{\pi}\sin nx\cos mx\,dx\EMbr =\frac12\int_{-\pi}^{\pi}\sin(n+m)x\,dx+\frac12\int_{-\pi}^{\pi}\sin(n-m)x\,dx}

\end{EMproof}

\EMdisp{\int_{-\pi}^{\pi}\sin^2nt\,dt\EMbr =\pi,\qquad\EMbr  \int_{-\pi}^{\pi}\cos^2nt\,dt\EMbr =\pi,\qquad\EMbr  n\in\mathbb{Z}\setminus\{0\}}

\begin{EMproof}

\EMdisp{\int_{-\pi}^{\pi}\sin^2nt\,dt\EMbr =\int_{-\pi}^{\pi}\frac{1-\cos2nt}{2}\,dt\EMbr =\pi,\qquad\EMbr  \int_{-\pi}^{\pi}\cos^2nt\,dt\EMbr =\int_{-\pi}^{\pi}\frac{1+\cos2nt}{2}\,dt\EMbr =\pi}

\end{EMproof}

\EMhold

\EMsection{سری فوریه}{سری فوریه}

\begin{EMremark}{تاریخچه}

ژوزف فوریه (۱۷۶۸--۱۸۳۰) ایدهٔ بیان توابع متناوب به صورت یک سری از توابع سینوسی را ارائه کرد. در زمان ارائهٔ این ایدهٔ انقلابی، بعضی از ریاضیدانان، از جمله لاگرانژ با عمومیت این ایده مخالفت کردند تا زمانی که دریشله شرایط همگرایی سری فوریه را بیان کرد.

\end{EMremark}

\begin{EMtheorem}{سری فوریه}

اگر تابع \EMim{f(x)} با دورهٔ تناوب \EMim{2\pi} متناوب باشد، آنگاه می‌توان آن را بر حسب یک سری مثلثاتی به صورت زیر بیان کرد:
\EMdisp{f(x)\EMbr =a_0+\sum_{n=1}^{\infty}\big(a_n\cos nx+b_n\sin nx\big)}
که در آن
\EMdisp{a_0\EMbr =\frac1{2\pi}\int_{-\pi}^{\pi}f(x)\,dx,\qquad\EMbr  a_n\EMbr =\frac1\pi\int_{-\pi}^{\pi}f(x)\cos nx\,dx,\qquad\EMbr  b_n\EMbr =\frac1\pi\int_{-\pi}^{\pi}f(x)\sin nx\,dx}

\end{EMtheorem}

\begin{EMproof}

فرض کنید تابع متناوب را به صورت سری مثلثاتی بالا نوشته‌ایم؛ کافی است ضرایب \EMim{a_0}، \EMim{a_n} و \EMim{b_n} را پیدا کنیم.

\textbf{محاسبهٔ \EMim{a_0}.} اگر از طرفین تساوی بالا روی بازه‌ای به طول \EMim{2\pi} انتگرال بگیریم، خواهیم داشت:
\EMdisp{\int_{-\pi}^{\pi}f(x)\,dx\EMbr =\int_{-\pi}^{\pi}\left[a_0+\sum_{n=1}^{\infty}(a_n\cos nx+b_n\sin nx)\right]dx\EMbr =a_0\int_{-\pi}^{\pi}dx+\sum_{n=1}^{\infty}\left(a_n\underbrace{\int_{-\pi}^{\pi}\cos nx\,dx}_{0}+b_n\underbrace{\int_{-\pi}^{\pi}\sin nx\,dx}_{0}\right)\EMbr =2\pi a_0}
\EMdisp{\Longrightarrow\quad\EMbr  a_0\EMbr =\frac1{2\pi}\int_{-\pi}^{\pi}f(x)\,dx}
\EMim{a_0} مقدار متوسط یا سطح \lr{DC} تابع است.

\textbf{محاسبهٔ \EMim{a_n}.} اگر از طرفین تساوی بالا در \EMim{\cos mx} ضرب کنیم و روی بازه‌ای به طول \EMim{2\pi} انتگرال بگیریم، خواهیم داشت:
\EMdisp{\int_{-\pi}^{\pi}f(x)\cos mx\,dx\EMbr =\int_{-\pi}^{\pi}a_0\cos mx\,dx+\sum_{n=1}^{\infty}\int_{-\pi}^{\pi}a_n\cos nx\cos mx\,dx+\sum_{n=1}^{\infty}\int_{-\pi}^{\pi}b_n\sin nx\cos mx\,dx\EMbr =\pi a_m}
\EMdisp{\Longrightarrow\quad\EMbr  a_n\EMbr =\frac1\pi\int_{-\pi}^{\pi}f(x)\cos nx\,dx}

\textbf{محاسبهٔ \EMim{b_n}.} اگر از طرفین تساوی بالا در \EMim{\sin mx} ضرب کنیم و روی بازه‌ای به طول \EMim{2\pi} انتگرال بگیریم، خواهیم داشت:
\EMdisp{\int_{-\pi}^{\pi}f(x)\sin mx\,dx\EMbr =\int_{-\pi}^{\pi}a_0\sin mx\,dx+\sum_{n=1}^{\infty}\int_{-\pi}^{\pi}a_n\cos nx\sin mx\,dx+\sum_{n=1}^{\infty}\int_{-\pi}^{\pi}b_n\sin nx\sin mx\,dx\EMbr =\pi b_m}
\EMdisp{\Longrightarrow\quad\EMbr  b_n\EMbr =\frac1\pi\int_{-\pi}^{\pi}f(x)\sin nx\,dx}

\end{EMproof}

\begin{EMimportant}{سری فوریه و روابط اویلر}

\EMdisp{f(x)\EMbr =a_0+\sum_{n=1}^{\infty}\big(a_n\cos nx+b_n\sin nx\big),\qquad\EMbr  f(x+2\pi)\EMbr =f(x)}
\EMdisp{a_0\EMbr =\frac1{2\pi}\int_{2\pi}f(x)\,dx,\qquad\EMbr  a_n\EMbr =\frac1\pi\int_{2\pi}f(x)\cos nx\,dx,\qquad\EMbr  b_n\EMbr =\frac1\pi\int_{2\pi}f(x)\sin nx\,dx}
(انتگرال روی یک دورهٔ تناوب \EMim{2\pi} گرفته می‌شود.)

\end{EMimportant}

\begin{EMremark}{}

مقدار سری فوریه در نقاط ناپیوستگی تابع برابر است با متوسط حد بالا و پایین در آن نقطه.

\end{EMremark}

\EMsection{یادآوری: انتگرال معین حاصل‌ضرب دو تابع}{یادآوری: انتگرال معین حاصل‌ضرب دو تابع}

با استفاده از روش جزء به جزء:
\EMdisp{\int_\alpha^\beta f(x)g(x)\,dx\EMbr =\Big[f(x)G_1(x)-f^{(1)}(x)G_2(x)+f^{(2)}(x)G_3(x)-+\cdots\Big]_\alpha^\beta}
که در آن \EMim{f^{(k)}} مشتق \EMim{k}-ام \EMim{f} و \EMim{G_k} انتگرال \EMim{k}-ام \EMim{g} است:

{\def\LTcaptype{none} % do not increment counter
\Needspace{13\baselineskip}
\begin{longtable}[c]{@{}ccc@{}}
\toprule\noalign{}
\rowcolor{EMindigoL}مشتق \EMim{f} & & انتگرال \EMim{g} \\
\midrule\noalign{}
\endhead
\bottomrule\noalign{}
\endlastfoot
\EMim{f(x)} & \EMim{\searrow} & \EMim{g(x)} \\
\EMim{f^{(1)}(x)} & \EMim{\searrow} & \EMim{G_1(x)} \\
\EMim{f^{(2)}(x)} & \EMim{\searrow} & \EMim{G_2(x)} \\
\EMim{f^{(3)}(x)} & \EMim{\searrow} & \EMim{G_3(x)} \\
\EMim{\vdots} & & \EMim{\vdots} \\
\end{longtable}
}

\EMhold

\EMsection{مثال‌های سری فوریه}{مثال‌های سری فوریه}

\begin{EMexample}{}

سری فوریهٔ تابع زیر را به دست آورید:
\EMdisp{f(x)\EMbr =\begin{cases}-k&-\pi<x<0\\ \ \ k&0<x<\pi\end{cases}\qquad\EMbr  f(x+2\pi)\EMbr =f(x)}

\EMdisp{a_0\EMbr =\frac1{2\pi}\int_{-\pi}^{\pi}f(x)\,dx\EMbr =0}
\EMdisp{a_n\EMbr =\frac1\pi\int_{-\pi}^{\pi}f(x)\cos nx\,dx\EMbr =\frac1\pi\left[\int_{-\pi}^{0}(-k)\cos nx\,dx+\int_0^\pi k\cos nx\,dx\right]\EMbr =\frac1\pi\left[-k\,\frac{\sin nx}{n}\Big|_{-\pi}^{0}+k\,\frac{\sin nx}{n}\Big|_0^\pi\right]\EMbr =0}
\EMdisp{\begin{aligned}
b_n&=\frac1\pi\int_{-\pi}^{\pi}f(x)\sin nx\,dx=\frac1\pi\left[\int_{-\pi}^{0}(-k)\sin nx\,dx+\int_0^\pi k\sin nx\,dx\right]\\
&=\frac1\pi\left[k\,\frac{\cos nx}{n}\Big|_{-\pi}^{0}-k\,\frac{\cos nx}{n}\Big|_0^\pi\right]=\frac k{n\pi}\big[\cos0-\cos(-n\pi)-\cos n\pi+\cos0\big]=\frac{2k}{n\pi}(1-\cos n\pi)
\end{aligned}}
\EMdisp{\cos n\pi\EMbr =\begin{cases}-1&n\ \EMt{فرد}\\ \ \ 1&n\ \EMt{زوج}\end{cases}\quad\EMbr \EMbr \Longrightarrow\quad\EMbr 1-\cos n\pi\EMbr =\begin{cases}2&n\ \EMt{فرد}\\ 0&n\ \EMt{زوج}\end{cases}}
\EMdisp{b_1\EMbr =\frac{4k}{\pi},\quad\EMbr  b_2\EMbr =0,\quad\EMbr  b_3\EMbr =\frac{4k}{3\pi},\quad\EMbr  b_4\EMbr =0,\quad\EMbr  b_5\EMbr =\frac{4k}{5\pi},\ \dots}
\EMdisp{f(x)\EMbr =\frac{4k}\pi\left(\sin x+\frac13\sin3x+\frac15\sin5x+\cdots\right)}

\end{EMexample}

\EMfigure{m08-square-wave}{موج مربعی \EMim{f(x)} با دورهٔ تناوب \EMim{2\pi} و مقادیر \EMim{\pm k}}

مجموع‌های جزئی سری:

\EMdisp{S_1\EMbr =\frac{4k}\pi\sin x,\qquad\EMbr  S_2\EMbr =\frac{4k}\pi\left(\sin x+\frac13\sin3x\right),\qquad\EMbr  S_3\EMbr =\frac{4k}\pi\left(\sin x+\frac13\sin3x+\frac15\sin5x\right)}

\EMfigure{m08-partial-sums}{مجموع‌های جزئی \EMim{S_1}، \EMim{S_2} و \EMim{S_3} سری فوریهٔ موج مربعی
(خط‌چین خود موج مربعی است)}

با قرار دادن \EMim{x=\pi/2}:
\EMdisp{f(\pi/2)\EMbr =\frac{4k}\pi\left(1-\frac13+\frac15-\frac17+-\cdots\right)\EMbr =k\quad\EMbr \EMbr \Longrightarrow\quad\EMbr \pi\EMbr =4\left(1-\frac13+\frac15-\frac17+-\cdots\right)}

\begin{EMexample}{}

سری فوریهٔ تابع زیر را به دست آورید:
\EMdisp{f(x)\EMbr =x,\qquad\EMbr  -\pi\le x\le\pi,\qquad\EMbr  f(x+2\pi)\EMbr =f(x)}

\EMdisp{a_0\EMbr =\frac1{2\pi}\int_{-\pi}^{\pi}f(x)\,dx\EMbr =0}
\EMdisp{a_n\EMbr =\frac1\pi\int_{-\pi}^{\pi}x\cos nx\,dx\EMbr =\frac1\pi\left[x\left(\frac1n\sin nx\right)-(1)\left(-\frac1{n^2}\cos nx\right)\right]_{-\pi}^{\pi}\EMbr =\frac1\pi\left[\frac1{n^2}\big(\cos n\pi-\cos(-n\pi)\big)\right]\EMbr =0}
\EMdisp{b_n\EMbr =\frac1\pi\int_{-\pi}^{\pi}x\sin nx\,dx\EMbr =\frac2\pi\int_0^\pi x\sin nx\,dx\EMbr =\frac2\pi\left[x\left(\frac{-1}n\cos nx\right)-(1)\left(\frac{-1}{n^2}\sin nx\right)\right]_0^\pi\EMbr =\frac2\pi\left[\frac{-\pi}n\cos n\pi\right]\EMbr =\frac2n(-1)^{n+1}}
\EMdisp{f(x)\EMbr =\sum_{n=1}^{\infty}\frac2n(-1)^{n+1}\sin nx\EMbr =2\left(\sin x-\frac12\sin2x+\frac13\sin3x-+\cdots\right)}

\end{EMexample}

\EMfigure{m08-sawtooth}{موج دندانه‌اره‌ای \EMim{f(x)=x} روی \EMim{[-\pi,\pi]} با ادامهٔ متناوب}

با قرار دادن \EMim{x=\pi/2}:
\EMdisp{f\!\left(\frac\pi2\right)\EMbr =\frac\pi2\EMbr =\sum_{n=1}^{\infty}\frac2n(-1)^{n+1}\sin\!\left(n\frac\pi2\right)\EMbr =2\left(1-\frac13+\frac15-\frac17+-\cdots\right)\EMbr =\sum_{n=1}^{\infty}\frac2{2n-1}(-1)^{n+1}\ \EMbr \Longrightarrow\ \pi\EMbr =\sum_{n=1}^{\infty}\frac4{2n-1}(-1)^{n+1}}

مقدار سری در نقطهٔ ناپیوستگی \EMim{x=\pi}:
\EMdisp{f(\pi)\EMbr =\;?\quad\EMbr \EMbr \Longrightarrow\quad\EMbr \sum_{n=1}^{\infty}\frac2n(-1)^{n+1}\sin(n\pi)\EMbr =0}

\begin{EMimportant}{نکات مهم}

\begin{enumerate}
\def\labelenumi{\arabic{enumi}.}
\tightlist
\item
  اگر تابع تناوبی فرد باشد، ضرایب \EMim{a} همگی صفرند.
\item
  مقدار سری فوریه در نقطهٔ ناپیوستگی تابع برابر است با متوسط حد بالا و پایین در آن نقطه.
\end{enumerate}

\end{EMimportant}

\begin{EMexample}{}

سری فوریهٔ تابع زیر را به دست آورید:
\EMdisp{f(x)\EMbr =x^2,\qquad\EMbr  -\pi\le x\le\pi,\qquad\EMbr  f(x+2\pi)\EMbr =f(x)}

\EMdisp{a_0\EMbr =\frac1{2\pi}\int_{-\pi}^{\pi}x^2\,dx\EMbr =\frac{2}{2\pi}\int_0^\pi x^2\,dx\EMbr =\frac1\pi\left[\frac13x^3\right]_0^\pi\EMbr =\frac{\pi^2}{3}}
\EMdisp{a_n\EMbr =\frac1\pi\int_{-\pi}^{\pi}x^2\cos nx\,dx\EMbr =\frac2\pi\int_0^\pi x^2\cos nx\,dx\EMbr =\frac2\pi\left[x^2\left(\frac1n\sin nx\right)-2x\left(\frac{-1}{n^2}\cos nx\right)+2\left(\frac{-1}{n^3}\sin nx\right)\right]_0^\pi\EMbr =\frac2\pi\left[\frac{2\pi}{n^2}\cos n\pi\right]\EMbr =\frac4{n^2}(-1)^n}
\EMdisp{b_n\EMbr =\frac1\pi\int_{-\pi}^{\pi}x^2\sin nx\,dx\EMbr =0}
\EMdisp{f(x)\EMbr =\frac{\pi^2}3+\sum_{n=1}^{\infty}\frac4{n^2}(-1)^n\cos nx}

\end{EMexample}

\EMfigure{m08-parabola}{تابع \EMim{f(x)=x^2} روی \EMim{[-\pi,\pi]} با ادامهٔ متناوب}

\begin{EMimportant}{نکتهٔ مهم}

اگر تابع تناوبی زوج باشد، ضرایب \EMim{b} همگی صفرند.

\end{EMimportant}

با قرار دادن \EMim{x=0}:
\EMdisp{f(0)\EMbr =0\EMbr =\frac{\pi^2}3+\sum_{n=1}^{\infty}\frac4{n^2}(-1)^n\cos(n\cdot0)\quad\EMbr \EMbr \Longrightarrow\quad\EMbr \frac{\pi^2}3\EMbr =\sum_{n=1}^{\infty}\frac4{n^2}(-1)^{n+1}}
\EMdisp{\pi^2\EMbr =\sum_{n=1}^{\infty}\frac{12}{n^2}(-1)^{n+1}\EMbr =12\left(1-\frac14+\frac19-\frac1{16}+-\cdots\right)}

\begin{EMexample}{}

سری فوریهٔ تابع زیر را به دست آورید:
\EMdisp{f(x)\EMbr =x^2,\qquad\EMbr  0\le x\le2\pi,\qquad\EMbr  f(x+2\pi)\EMbr =f(x)}

\EMdisp{a_0\EMbr =\frac1{2\pi}\int_0^{2\pi}x^2\,dx\EMbr =\frac1{2\pi}\left[\frac13x^3\right]_0^{2\pi}\EMbr =\frac1{2\pi}\cdot\frac138\pi^3\EMbr =\frac{4\pi^2}3}
\EMdisp{a_n\EMbr =\frac1\pi\int_0^{2\pi}x^2\cos nx\,dx\EMbr =\frac1\pi\left[x^2\left(\frac1n\sin nx\right)-2x\left(\frac{-1}{n^2}\cos nx\right)+2\left(\frac{-1}{n^3}\sin nx\right)\right]_0^{2\pi}\EMbr =\frac1\pi(2\times2\pi)\left(\frac1{n^2}\cos2\pi n\right)\EMbr =\frac4{n^2}}
\EMdisp{b_n\EMbr =\frac1\pi\int_0^{2\pi}x^2\sin nx\,dx\EMbr =\frac1\pi\left[x^2\left(\frac{-1}n\cos nx\right)-2x\left(\frac{-1}{n^2}\sin nx\right)+2\left(\frac1{n^3}\cos nx\right)\right]_0^{2\pi}\EMbr =\frac1\pi\left[(4\pi^2)\left(\frac{-1}n\right)+2\cdot\frac1{n^3}-2\cdot\frac1{n^3}\right]\EMbr =-\frac{4\pi}n}
\EMdisp{f(x)\EMbr =\frac{4\pi^2}3+\sum_{n=1}^{\infty}\left(\frac4{n^2}\cos nx-\frac{4\pi}n\sin nx\right)}

\end{EMexample}

\EMfigure{m08-shifted-parabola}{تابع \EMim{f(x)=x^2} روی \EMim{[0,2\pi]} با ادامهٔ متناوب}

\EMsection{سری فوریهٔ توابع متناوب با دورهٔ تناوب دلخواه}{سری فوریهٔ توابع متناوب با دورهٔ تناوب دلخواه}

\EMdisp{f(x+T)\EMbr =f(x)}
اگر از متغیر کمکی \EMim{t=\dfrac{2\pi}{T}x} استفاده کنیم (\EMim{x:0\to T} و \EMim{t:0\to2\pi}):
\EMdisp{f\!\left(\frac T{2\pi}t\right)\EMbr =a_0+\sum_{n=1}^{\infty}a_n\cos nt+b_n\sin nt}
\EMdisp{a_0\EMbr =\frac1{2\pi}\int_{2\pi}f\!\left(\frac T{2\pi}t\right)dt,\qquad\EMbr  a_n\EMbr =\frac1\pi\int_{2\pi}f\!\left(\frac T{2\pi}t\right)\cos nt\,dt,\qquad\EMbr  b_n\EMbr =\frac1\pi\int_{2\pi}f\!\left(\frac T{2\pi}t\right)\sin nt\,dt}
با بازگشت به \EMim{x=\dfrac T{2\pi}t}:
\EMdisp{f(x)\EMbr =a_0+\sum_{n=1}^{\infty}a_n\cos\frac{2\pi}Tnx+b_n\sin\frac{2\pi}Tnx}
\EMdisp{a_0\EMbr =\frac1{2\pi}\int_Tf(x)\frac{2\pi}T\,dx\EMbr =\frac1T\int_Tf(x)\,dx}
\EMdisp{a_n\EMbr =\frac1\pi\int_Tf(x)\cos\!\left(\frac{2\pi}Tnx\right)\frac{2\pi}T\,dx\EMbr =\frac2T\int_Tf(x)\cos\!\left(\frac{2\pi}Tnx\right)dx}
\EMdisp{b_n\EMbr =\frac1\pi\int_Tf(x)\sin\!\left(\frac{2\pi}Tnx\right)\frac{2\pi}T\,dx\EMbr =\frac2T\int_Tf(x)\sin\!\left(\frac{2\pi}Tnx\right)dx}

\begin{EMtheorem}{سری فوریه (دورهٔ تناوب \EMim{T})}

اگر تابع \EMim{f(x)} با دورهٔ تناوب \EMim{T} متناوب باشد، آنگاه می‌توان آن را بر حسب یک سری مثلثاتی به صورت زیر بیان کرد:
\EMdisp{f(x)\EMbr =a_0+\sum_{n=1}^{\infty}a_n\cos\frac{2\pi}Tnx+b_n\sin\frac{2\pi}Tnx,\qquad\EMbr  f(x+T)\EMbr =f(x)}
\EMdisp{a_0\EMbr =\frac1T\int_Tf(x)\,dx,\qquad\EMbr  a_n\EMbr =\frac2T\int_Tf(x)\cos\frac{2\pi}Tnx\,dx,\qquad\EMbr  b_n\EMbr =\frac2T\int_Tf(x)\sin\frac{2\pi}Tnx\,dx}

\end{EMtheorem}

\begin{EMremark}{}

مقدار سری فوریه در نقاط ناپیوستگی تابع برابر است با متوسط حد بالا و پایین در آن نقطه؛ مقدار سری فوریه در نقطهٔ \EMim{x}:
\EMdisp{\frac{f(x^+)+f(x^-)}{2}}

\end{EMremark}

\EMhold

\EMsection{قضیهٔ دریشله}{قضیهٔ دریشله}

\begin{EMtheorem}{قضیهٔ دریشله (\lr{Dirichlet Theorem})}

اگر تابع \EMim{f(x)} متناوب و \textbf{محدود} بوده، دارای تعداد \textbf{محدود ماکزیمم و مینیمم} و نیز تعداد \textbf{محدود نقاط ناپیوستگی} در هر دورهٔ تناوب باشد، آنگاه این تابع دارای سری فوریه بوده، مقدار سری فوریه در نقاط ناپیوستگی برابر مقدار متوسط حد راست و چپ در این نقاط است.

\end{EMtheorem}

مثال‌هایی از توابعی که شرایط دریشله را ندارند:

\EMdisp{f(x)\EMbr =1/x,\quad\EMbr  0<x\le1,\qquad\EMbr  f(x+1)\EMbr =f(x)}
\EMdisp{f(x)\EMbr =\sin(1/x),\quad\EMbr  0<x\le1,\qquad\EMbr  f(x+1)\EMbr =f(x)}

\EMfigure{m08-dirichlet}{توابعی که شرایط دریشله را ندارند: \EMim{1/x} (نامحدود)، \EMim{\sin(1/x)}
(بی‌نهایت ماکزیمم و مینیمم) و تابع پله‌ای با بی‌نهایت پله}

\EMhold

\EMsection{حل چند مسئله}{حل چند مسئله}

\begin{EMexample}{مسئلهٔ ۱}

سری فوریهٔ تابع زیر را به دست آورید:
\EMdisp{f(x)\EMbr =\sin\alpha x,\qquad\EMbr  \alpha\in\mathbb{R},\qquad\EMbr  x\in[-\pi,\pi),\qquad\EMbr  f(x+2\pi)\EMbr =f(x)}

از آنجا که این تابع فرد است، داریم \EMim{a_0=0} و \EMim{a_n=0}.
\EMdisp{\begin{aligned}
b_n&=\frac1\pi\int_{-\pi}^{\pi}\sin\alpha x\sin nx\,dx=\frac2\pi\int_0^\pi\sin\alpha x\sin nx\,dx=\frac1\pi\int_0^\pi\big[\cos(\alpha-n)x-\cos(\alpha+n)x\big]dx\\
&=\frac1\pi\left[\frac{\sin(\alpha-n)x}{\alpha-n}-\frac{\sin(\alpha+n)x}{\alpha+n}\right]_0^\pi=\cdots=\frac{2n(-1)^n\sin\alpha\pi}{\pi(\alpha^2-n^2)}
\end{aligned}}
\EMdisp{f(x)\EMbr =\sum_{n=1}^{\infty}\frac{2n(-1)^n\sin\alpha\pi}{\pi(\alpha^2-n^2)}\sin nx}

سری را وقتی \EMim{\alpha=1} بررسی می‌کنیم:
\EMdisp{\lim_{\alpha\to1}b_1\EMbr =\lim_{\alpha\to1}\frac{2(-1)^1\sin\alpha\pi}{\pi(\alpha^2-1^2)}\EMbr =\lim_{\alpha\to1}\frac{-2(\pi\cos\alpha\pi)}{\pi(2\alpha)}\EMbr =1,\qquad\EMbr  b_n\Big|_{n\neq1}\EMbr =0}
بنابراین سری فوریه همان \EMim{f(x)=\sin x} است.

\end{EMexample}

\begin{EMexercise}{}

سری فوریه را وقتی \EMim{\alpha=m}، \EMim{m\in\mathbb{N}} بررسی کنید.

\end{EMexercise}

\begin{EMexample}{مسئلهٔ ۲}

سری فوریهٔ تابع زیر را به دست آورید و \EMim{\csc\alpha\pi} را از روی آن بیابید.
\EMdisp{f(x)\EMbr =\cos\alpha x,\qquad\EMbr  \alpha\in\mathbb{R},\qquad\EMbr  x\in[-\pi,\pi),\qquad\EMbr  f(x+2\pi)\EMbr =f(x)}

از آنجا که این تابع زوج است، داریم \EMim{b_n=0} و \EMim{f(x)=a_0+\sum_{n=1}^{\infty}a_n\cos nx}.
\EMdisp{a_0\EMbr =\frac1{2\pi}\int_{-\pi}^{\pi}\cos\alpha x\,dx\EMbr =\frac1\pi\int_0^\pi\cos\alpha x\,dx\EMbr =\frac1\pi\left[\frac{\sin\alpha x}{\alpha}\right]_0^\pi\EMbr =\frac1{\alpha\pi}\sin\alpha\pi}
\EMdisp{a_n\EMbr =\frac2\pi\int_0^\pi\cos\alpha x\cos nx\,dx\EMbr =\frac1\pi\int_0^\pi\big[\cos(\alpha-n)x+\cos(\alpha+n)x\big]dx\EMbr =\frac1\pi\left[\frac{\sin(\alpha-n)x}{\alpha-n}+\frac{\sin(\alpha+n)x}{\alpha+n}\right]_0^\pi\EMbr =\cdots\EMbr =\frac{(-1)^n2\alpha\sin\alpha\pi}{\pi(\alpha^2-n^2)}}
\EMdisp{f(x)\EMbr =\frac1{\alpha\pi}\sin\alpha\pi+\frac{2\alpha\sin\alpha\pi}\pi\sum_{n=1}^{\infty}\frac{(-1)^n}{(\alpha^2-n^2)}\cos nx}
تابع تناوبی در \EMim{x=0} پیوسته است، بنابراین داریم:
\EMdisp{f(0)\EMbr =1\EMbr =\frac1{\alpha\pi}\sin\alpha\pi+\frac{2\alpha\sin\alpha\pi}\pi\sum_{n=1}^{\infty}\frac{(-1)^n}{(\alpha^2-n^2)}\quad\EMbr \EMbr \Longrightarrow\quad\EMbr \frac1{\sin\alpha\pi}\EMbr =\csc\alpha\pi\EMbr =\frac1{\alpha\pi}+\frac{2\alpha}\pi\sum_{n=1}^{\infty}\frac{(-1)^n}{(\alpha^2-n^2)}}

\end{EMexample}

\begin{EMexercise}{}

ثابت کنید: اگر \EMim{\alpha=m}، \EMim{m\in\mathbb{N}} خواهیم داشت:
\EMdisp{a_0\EMbr =0,\qquad\EMbr  a_n\EMbr =\begin{cases}1&n=m\\ 0&n\neq m\end{cases}}

\end{EMexercise}

\begin{EMexample}{مسئلهٔ ۳}

سری فوریهٔ تابع زیر را به دست آورید و \EMim{\coth\alpha\pi} را از روی آن بیابید.
\EMdisp{f(x)\EMbr =e^{-\alpha x},\qquad\EMbr  \alpha\in\mathbb{R},\qquad\EMbr  x\in[-\pi,\pi),\qquad\EMbr  f(x+2\pi)\EMbr =f(x)}

\textbf{لم ۱:}
\EMdisp{\int e^{\alpha x}\cos nx\,dx\EMbr =\frac{e^{\alpha x}}{\alpha^2+n^2}\big(\alpha\cos nx+n\sin nx\big)}
اثبات: با انتگرال‌گیری جزء به جزء دو بار، با \EMim{I=\int e^{\alpha x}\cos nx\,dx}:
\EMdisp{I\EMbr =\frac1\alpha e^{\alpha x}\cos nx-\int\frac1\alpha e^{\alpha x}(-n\sin nx)\,dx\EMbr =\frac1\alpha e^{\alpha x}\cos nx+\frac n\alpha\int e^{\alpha x}\sin nx\,dx}
\EMdisp{=\frac1\alpha e^{\alpha x}\cos nx+\frac n\alpha\left[\frac1\alpha e^{\alpha x}\sin nx-\int\frac1\alpha e^{\alpha x}(n\cos nx)\,dx\right]\EMbr =\frac1\alpha e^{\alpha x}\cos nx+\frac n{\alpha^2}e^{\alpha x}\sin nx-\frac{n^2}{\alpha^2}I}
\EMdisp{\Longrightarrow\quad\EMbr  I\EMbr =\frac{e^{\alpha x}}{\alpha^2+n^2}\big(\alpha\cos nx+n\sin nx\big)}

\textbf{لم ۲:}
\EMdisp{\int e^{\alpha x}\sin nx\,dx\EMbr =\frac{e^{\alpha x}}{\alpha^2+n^2}\big(\alpha\sin nx-n\cos nx\big)}
(اثبات به عنوان تمرین.)

تابع نه زوج است و نه فرد؛ لذا باید همهٔ ضرایب سری فوریه را حساب کنیم.
\EMdisp{a_0\EMbr =\frac1{2\pi}\int_{-\pi}^{\pi}e^{-\alpha x}\,dx\EMbr =\frac1{-2\pi\alpha}\Big[e^{-\alpha x}\Big]_{-\pi}^{\pi}\EMbr =\frac1{2\pi\alpha}\big[e^{\alpha\pi}-e^{-\alpha\pi}\big]\EMbr =\frac{\sinh\alpha\pi}{\alpha\pi}}
\EMdisp{a_n\EMbr =\frac1\pi\int_{-\pi}^{\pi}e^{-\alpha x}\cos nx\,dx\EMbr =\left\{\frac{e^{-\alpha x}}{\pi(\alpha^2+n^2)}\big(-\alpha\cos nx+n\sin nx\big)\right\}_{-\pi}^{\pi}\EMbr =\frac{\alpha(-1)^n}{\pi(\alpha^2+n^2)}\big[e^{\alpha\pi}-e^{-\alpha\pi}\big]\EMbr =\frac{2\alpha(-1)^n}{\pi(\alpha^2+n^2)}\sinh\alpha\pi}
\EMdisp{b_n\EMbr =\frac1\pi\int_{-\pi}^{\pi}e^{-\alpha x}\sin nx\,dx\EMbr =\left\{\frac{e^{-\alpha x}}{\pi(\alpha^2+n^2)}\big(-\alpha\sin nx-n\cos nx\big)\right\}_{-\pi}^{\pi}\EMbr =\frac{n(-1)^n}{\pi(\alpha^2+n^2)}\big[e^{\alpha\pi}-e^{-\alpha\pi}\big]\EMbr =\frac{2n(-1)^n}{\pi(\alpha^2+n^2)}\sinh\alpha\pi}
\EMdisp{\begin{aligned}
f(x)&=\frac{\sinh\alpha\pi}{\alpha\pi}+\sum_{n=1}^{\infty}\left(\frac{2\alpha(-1)^n}{\pi(\alpha^2+n^2)}\sinh\alpha\pi\cos nx+\frac{2n(-1)^n}{\pi(\alpha^2+n^2)}\sinh\alpha\pi\sin nx\right)\\
&=\sinh\alpha\pi\left(\frac1{\alpha\pi}+\sum_{n=1}^{\infty}\left(\frac{2\alpha(-1)^n}{\pi(\alpha^2+n^2)}\cos nx+\frac{2n(-1)^n}{\pi(\alpha^2+n^2)}\sin nx\right)\right)
\end{aligned}}
از آنجا که نقطهٔ \EMim{x=\pi} نقطهٔ ناپیوستگی است مقدار سری فوریه در این نقطه برابر است با متوسط حد بالا و پایین تابع در این نقطه، به عبارت دیگر:
\EMdisp{\frac{f(\pi^+)+f(\pi^-)}2\EMbr =\frac{e^{\alpha\pi}+e^{-\alpha\pi}}2\EMbr =\cosh\alpha\pi}
\EMdisp{\cosh\alpha\pi\EMbr =\sinh\alpha\pi\left(\frac1{\alpha\pi}+\sum_{n=1}^{\infty}\left(\frac{2\alpha(-1)^n}{\pi(\alpha^2+n^2)}\cos n\pi+\frac{2n(-1)^n}{\pi(\alpha^2+n^2)}\sin n\pi\right)\right)\EMbr =\sinh\alpha\pi\left(\frac1{\alpha\pi}+\sum_{n=1}^{\infty}\frac{2\alpha}{\pi(\alpha^2+n^2)}\right)}
\EMdisp{\coth\alpha\pi\EMbr =\frac{\cosh\alpha\pi}{\sinh\alpha\pi}\EMbr =\frac1{\alpha\pi}+\sum_{n=1}^{\infty}\frac{2\alpha}{\pi(\alpha^2+n^2)}}

\end{EMexample}

\EMsetmodule{9}{خواص سری فوریه، قضیهٔ پارسوال و صورت مختلط}{Properties of Fourier Series, Parseval's Theorem and the Complex Form}
\EMchapterpage{9}{9}{خواص سری فوریه، قضیهٔ پارسوال و صورت مختلط}{Properties of Fourier Series, Parseval's Theorem and the Complex Form}
\begin{EMminitoc}
\EMtocitem{1}{تقریب یک تابع متناوب توسط یک سری مثلثاتی با تعداد جملات محدود}
\EMtocitem{2}{مشتق و انتگرال سری فوریه}
\EMtocitem{3}{قضیهٔ رایلی}
\EMtocitem{4}{قضیهٔ پارسوال}
\EMtocitem{5}{صورت مختلط سری فوریه}
\end{EMminitoc}
\EMchapterend
\EMhold

\EMsection{تقریب یک تابع متناوب توسط یک سری مثلثاتی با تعداد جملات محدود}{تقریب یک تابع متناوب توسط یک سری مثلثاتی با تعداد جملات محدود}

\begin{EMtheorem}{}

برای تقریب \EMim{N} جمله‌ای یک تابع متناوب \EMim{f(x)} توسط یک سری سینوسی و کسینوسی، بهترین ضرایب (با معیار \lr{MMSE}) ضرایب سری فوریه می‌باشند.

\end{EMtheorem}

تابع متناوب با دورهٔ تناوب \EMim{2\pi} را در نظر بگیرید:
\EMdisp{f(x)\EMbr =a_0+\sum_{n=1}^{\infty}a_n\cos nx+b_n\sin nx,\qquad\EMbr  f(x)\EMbr =f(x+2\pi)}
تقریب \EMim{N} جمله‌ای مثلثاتی:
\EMdisp{f(x)\EMbr \approx a_0+\sum_{n=1}^{N}a_n\cos nx+b_n\sin nx}
آیا تقریب مثلثاتی مقابل \textbf{بهتر} نیست؟
\EMdisp{F(x)\triangleq\alpha_0+\sum_{n=1}^{N}\alpha_n\cos nx+\beta_n\sin nx}
\textbf{معیار} بهتر بودن تقریب چیست؟ یک معیار: تابع خطا:
\EMdisp{E(x)\EMbr =\big|f(x)-F(x)\big|\EMbr =\left|f(x)-\alpha_0-\sum_{n=1}^{N}\big(\alpha_n\cos nx+\beta_n\sin nx\big)\right|}
یک معیار خوب: متوسط مجذور خطا:
\EMdisp{D\triangleq\int_{-\pi}^{\pi}E^2(x)\,dx\EMbr =\int_{-\pi}^{\pi}\left|f(x)-\alpha_0-\sum_{n=1}^{N}\big(\alpha_n\cos nx+\beta_n\sin nx\big)\right|^2dx}

\begin{EMproof}

برای بهترین تقریب متوسط مجذور خطا را کمینه می‌کنیم. این کار با مشتق‌گیری از \EMim{D} بر حسب ضرایب سری مثلثاتی و برابر قرار دادن این مشتق با صفر به دست می‌آید. به عبارت دیگر:
\EMdisp{\frac{\partial D}{\partial\alpha_0}\EMbr =-2\int_{-\pi}^{\pi}\left[f(x)-\alpha_0-\sum_{n=1}^{N}\big(\alpha_n\cos nx+\beta_n\sin nx\big)\right]dx\EMbr =0}
\EMdisp{\int_{-\pi}^{\pi}f(x)\,dx\EMbr =\int_{-\pi}^{\pi}\left[\alpha_0+\sum_{n=1}^{N}\big(\alpha_n\cos nx+\beta_n\sin nx\big)\right]dx\EMbr =2\pi\alpha_0\quad\EMbr \EMbr \Longrightarrow\quad\EMbr \alpha_0\EMbr =\frac1{2\pi}\int_{-\pi}^{\pi}f(x)\,dx\EMbr =a_0}
\EMdisp{\frac{\partial D}{\partial\alpha_m}\EMbr =0\quad\EMbr \EMbr \Longrightarrow\quad\EMbr \alpha_m\EMbr =\frac1\pi\int_{-\pi}^{\pi}f(x)\cos mx\,dx\EMbr =a_m}
\EMdisp{\frac{\partial D}{\partial\beta_m}\EMbr =0\quad\EMbr \EMbr \Longrightarrow\quad\EMbr \beta_m\EMbr =\frac1\pi\int_{-\pi}^{\pi}f(x)\sin mx\,dx\EMbr =b_m}
بهترین تقریب مثلثاتی با تعداد جملات محدود از یک تابع تناوبی سری فوریه با جملات محدود است.

\end{EMproof}

\EMhold

\EMsection{مشتق و انتگرال سری فوریه}{مشتق و انتگرال سری فوریه}

\begin{EMtheorem}{مشتق سری فوریه}

اگر تابع \EMim{f(x)} با دورهٔ تناوب \EMim{T} تناوبی بوده، پیوسته و دارای سری فوریه باشد، آنگاه مشتق این تابع نیز تناوبی بوده سری فوریه آن با استفاده از مشتق‌گیری از سری فوریه تابع به دست می‌آید.
\EMdisp{f(x)\EMbr =a_0+\sum_{n=1}^{\infty}a_n\cos\frac{2\pi}Tnx+b_n\sin\frac{2\pi}Tnx}
\EMdisp{\Longrightarrow\quad\EMbr  f'(x)\EMbr =\sum_{n=1}^{\infty}\left(\frac{2\pi}Tnb_n\right)\cos\frac{2\pi}Tnx+\left(-\frac{2\pi}Tna_n\right)\sin\frac{2\pi}Tnx}

\end{EMtheorem}

\begin{EMtheorem}{انتگرال سری فوریه}

اگر تابع \EMim{f(x)} با دورهٔ تناوب \EMim{T} تناوبی بوده و دارای سری فوریه باشد، آنگاه انتگرال معین این تابع با استفاده از انتگرال‌گیری از سری فوریه تابع قابل محاسبه است.
\EMdisp{f(x)\EMbr =a_0+\sum_{n=1}^{\infty}a_n\cos\frac{2\pi}Tnx+b_n\sin\frac{2\pi}Tnx}
\EMdisp{\Longrightarrow\quad\EMbr \int_0^xf(t)\,dt\EMbr =a_0x+\sum_{n=1}^{\infty}\left(\frac T{2\pi n}a_n\sin\frac{2\pi}Tnx-\frac T{2\pi n}b_n\cos\frac{2\pi}Tnx+\frac T{2\pi n}b_n\right)}

\end{EMtheorem}

\EMhold

\EMsection{قضیهٔ رایلی}{قضیهٔ رایلی}

\begin{EMtheorem}{قضیهٔ رایلی (\lr{Rayleigh’s Theorem})}

اگر دو تابع \EMim{f(x)} و \EMim{g(x)} با دورهٔ تناوب \EMim{T} تناوبی بوده و سری فوریه این دو تابع به صورت زیر وجود داشته باشند:
\EMdisp{f(x)\EMbr =a_0+\sum_{n=1}^{\infty}a_n\cos\frac{2\pi}Tnx+b_n\sin\frac{2\pi}Tnx}
\EMdisp{g(x)\EMbr =\alpha_0+\sum_{n=1}^{\infty}\alpha_n\cos\frac{2\pi}Tnx+\beta_n\sin\frac{2\pi}Tnx}
آنگاه خواهیم داشت:
\EMdisp{\frac2T\int_Tf(x)\,\overline{g(x)}\,dx\EMbr =2a_0\overline{\alpha_0}+\sum_{n=1}^{\infty}a_n\overline{\alpha_n}+b_n\overline{\beta_n}}

\end{EMtheorem}

\begin{EMproof}

اثبات با فرض حقیقی بودن \EMim{g(x)}:
\EMdisp{\begin{aligned}
\frac2T\int_Tf(x)g(x)\,dx&=\frac2T\int_Tf(x)\left(\alpha_0+\sum_{n=1}^{\infty}\alpha_n\cos\frac{2\pi}Tnx+\beta_n\sin\frac{2\pi}Tnx\right)dx\\
&=\frac2T\int_Tf(x)\alpha_0\,dx+\sum_{n=1}^{\infty}\frac2T\int_Tf(x)\left(\alpha_n\cos\frac{2\pi}Tnx+\beta_n\sin\frac{2\pi}Tnx\right)dx\\
&=\underbrace{\frac2T\int_Tf(x)\,dx}_{2a_0}\,\alpha_0+\sum_{n=1}^{\infty}\underbrace{\frac2T\int_Tf(x)\cos\frac{2\pi}Tnx\,dx}_{a_n}\,\alpha_n+\underbrace{\frac2T\int_Tf(x)\sin\frac{2\pi}Tnx\,dx}_{b_n}\,\beta_n\\
&=2a_0\alpha_0+\sum_{n=1}^{\infty}a_n\alpha_n+b_n\beta_n
\end{aligned}}

\end{EMproof}

\EMhold

\EMsection{قضیهٔ پارسوال}{قضیهٔ پارسوال}

\begin{EMtheorem}{قضیهٔ پارسوال (\lr{Parseval’s Theorem})}

اگر تابع \EMim{f(x)} با دورهٔ تناوب \EMim{T} تناوبی بوده و سری فوریه آن به صورت زیر وجود داشته باشد:
\EMdisp{f(x)\EMbr =a_0+\sum_{n=1}^{\infty}a_n\cos\frac{2\pi}Tnx+b_n\sin\frac{2\pi}Tnx}
آنگاه خواهیم داشت:
\EMdisp{\frac2T\int_T|f(x)|^2\,dx\EMbr =2|a_0|^2+\sum_{n=1}^{\infty}\big(|a_n|^2+|b_n|^2\big)}

\end{EMtheorem}

\begin{EMproof}

در تساوی رایلی اگر به جای \EMim{g(x)} تابع \EMim{f(x)} را قرار دهیم به تساوی پارسوال می‌رسیم.

\end{EMproof}

با توجه به تساوی پارسوال مقدار توان تابع \EMim{f(x)} را می‌شود با استفاده از ضرایب سری فوریه به دست آورد. اگر تابع توسط تقریب \EMim{N} جمله‌ای از سینوس‌ها و کسینوس‌ها بیان شود طبق رابطهٔ بالا توان تابع تقریبی کمتر از توان تابع \EMim{f(x)} خواهد بود و داریم:

\begin{EMimportant}{نامساوی بسل}

\EMdisp{\frac2T\int_T|f(x)|^2\,dx\ \ge\ 2|a_0|^2+\sum_{n=1}^{N}\big(|a_n|^2+|b_n|^2\big)}

\end{EMimportant}

\begin{EMexample}{}

برای تابع \EMim{f(x)} زیر، چند جمله از سری فوریه نیاز است تا حداقل ۹۴ درصد توان در تابع تقریبی باقی بماند.

\EMdisp{f(x)\EMbr =\sum_{n=1}^{\infty}\frac{2k}{n\pi}(1-\cos n\pi)\sin nx}
\EMdisp{\frac2T\int_T|f(x)|^2\,dx\EMbr =\frac1\pi\int_{-\pi}^{\pi}k^2\,dx\EMbr =\frac2\pi\big[k^2x\big]_0^\pi\EMbr =2k^2}
\EMdisp{\frac{\displaystyle\sum_{n=1}^{N}\left(\frac{2k}{n\pi}(1-\cos n\pi)\right)^2}{2k^2}\ge0.94}
\EMdisp{\frac{\displaystyle\sum_{n=1}^{7}\left(\frac{2k}{n\pi}(1-\cos n\pi)\right)^2}{2k^2}\EMbr \cong0.9496,\qquad\EMbr  \frac{\displaystyle\sum_{n=1}^{5}\left(\frac{2k}{n\pi}(1-\cos n\pi)\right)^2}{2k^2}\EMbr \cong0.9331\quad\EMbr \EMbr \Longrightarrow\quad\EMbr  N\EMbr =7}

\end{EMexample}

\EMfigure{m09-power-example}{موج مربعی \EMim{\pm k} با دورهٔ تناوب \EMim{2\pi} در مثال توان}

\EMsection{صورت مختلط سری فوریه}{صورت مختلط سری فوریه}

فرض کنید تابع \EMim{f(x)} تناوبی و صرفاً جهت سادگی دارای دورهٔ تناوب \EMim{2\pi} می‌باشد. لذا داریم:
\EMdisp{f(x)\EMbr =a_0+\sum_{n=1}^{\infty}a_n\cos nx+b_n\sin nx}
با استفاده از رابطهٔ اویلر \EMim{e^{ix}=\cos x+i\sin x}:
\EMdisp{\cos nx\EMbr =\frac{e^{inx}+e^{-inx}}2,\qquad\EMbr  \sin nx\EMbr =\frac{e^{inx}-e^{-inx}}{2i},\qquad\EMbr  \frac1i\EMbr =-i}
\EMdisp{f(x)\EMbr =a_0+\sum_{n=1}^{\infty}a_n\left(\frac{e^{inx}+e^{-inx}}2\right)+b_n\left(\frac{e^{inx}-e^{-inx}}{2i}\right)\EMbr =\underbrace{a_0}_{c_0}+\sum_{n=1}^{\infty}\underbrace{\left(\frac{a_n-ib_n}2\right)}_{c_n}e^{inx}+\underbrace{\left(\frac{a_n+ib_n}2\right)}_{c_{-n}}e^{-inx}}
\EMdisp{c_n\EMbr =\frac{a_n-ib_n}2\EMbr =\frac12\left\{\frac1\pi\int_{2\pi}f(x)\cos nx\,dx-i\frac1\pi\int_{2\pi}f(x)\sin nx\,dx\right\}\EMbr =\frac1{2\pi}\int_{2\pi}f(x)\big(\cos nx-i\sin nx\big)\,dx\EMbr =\frac1{2\pi}\int_{2\pi}f(x)e^{-inx}\,dx}
\EMdisp{c_{-n}\EMbr =\frac{a_n+ib_n}2\EMbr =\frac1{2\pi}\int_{2\pi}f(x)\big(\cos nx+i\sin nx\big)\,dx\EMbr =\frac1{2\pi}\int_{2\pi}f(x)e^{inx}\,dx}

\begin{EMimportant}{صورت مختلط سری فوریه}

\EMdisp{f(x)\EMbr =\sum_{n=-\infty}^{\infty}c_ne^{inx},\qquad\EMbr  c_n\EMbr =\frac1{2\pi}\int_{2\pi}f(x)e^{-inx}\,dx}
برای دورهٔ تناوب دلخواه:
\EMdisp{f(x)\EMbr =\sum_{n=-\infty}^{\infty}c_ne^{i\frac{2\pi}Tnx},\qquad\EMbr  c_n\EMbr =\frac1T\int_Tf(x)e^{-i\frac{2\pi}Tnx}\,dx}
\EMdisp{\frac1T\int_T|f(x)|^2\,dx\EMbr =\sum_{n=-\infty}^{\infty}|c_n|^2}

\end{EMimportant}

\EMsetmodule{10}{انتگرال فوریه}{The Fourier Integral}
\EMchapterpage{10}{10}{انتگرال فوریه}{The Fourier Integral}
\begin{EMminitoc}
\EMtocitem{1}{تحلیل فوریه برای توابع غیرتناوبی}
\EMtocitem{2}{انتگرال فوریه}
\EMtocitem{3}{صورت مختلط انتگرال فوریه (تبدیل فوریه)}
\EMtocitem{4}{معرفی تابع \EMim{\operatorname{Si}(u)}}
\EMtocitem{5}{مثال‌های انتگرال فوریه}
\EMtocitem{6}{بعضی خواص انتگرال فوریه}
\end{EMminitoc}
\EMchapterend
\EMsection{تحلیل فوریه برای توابع غیرتناوبی}{تحلیل فوریه برای توابع غیرتناوبی}

برای آنکه بتوانیم از ابزار تحلیل فوریه در توابع غیرتناوبی استفاده کنیم، ابتدا تابع را تناوبی فرض می‌کنیم و سپس دورهٔ تناوب را به سمت بی‌نهایت میل می‌دهیم.

\EMfigure{m10-periodic-extension}{تابع غیرتناوبی \EMim{f(x)} و ادامهٔ متناوب آن \EMim{\tilde f(x)} با دورهٔ
تناوب \EMim{T}}

\EMdisp{\tilde f(x)\EMbr =a_0+\sum_{n=1}^{\infty}a_n\cos\frac{2\pi}Tnx+b_n\sin\frac{2\pi}Tnx}
\EMdisp{a_0\EMbr =\frac1T\int_{-T/2}^{T/2}f(x)\,dx,\qquad\EMbr  a_n\EMbr =\frac2T\int_{-T/2}^{T/2}f(x)\cos\frac{2\pi}Tnx\,dx,\qquad\EMbr  b_n\EMbr =\frac2T\int_{-T/2}^{T/2}f(x)\sin\frac{2\pi}Tnx\,dx}
\EMdisp{\begin{aligned}
\tilde f(x)&=\left(\frac1T\int_{-T/2}^{T/2}f(t)\,dt\right)+\sum_{n=1}^{\infty}\Bigg\{\left(\frac2T\int_{-T/2}^{T/2}f(t)\cos\frac{2\pi}Tnt\,dt\right)\cos\frac{2\pi}Tnx\\
&\qquad+\left(\frac2T\int_{-T/2}^{T/2}f(t)\sin\frac{2\pi}Tnt\,dt\right)\sin\frac{2\pi}Tnx\Bigg\}
\end{aligned}}
\EMdisp{\tilde f(x)\EMbr =\left(\frac1T\int_{-T/2}^{T/2}f(t)\,dt\right)+\frac2T\sum_{n=1}^{\infty}\left\{\int_{-T/2}^{T/2}f(t)\left(\cos\frac{2\pi}Tnt\cos\frac{2\pi}Tnx+\sin\frac{2\pi}Tnt\sin\frac{2\pi}Tnx\right)dt\right\}}
\EMdisp{\tilde f(x)\EMbr =\left(\frac1T\int_{-T/2}^{T/2}f(t)\,dt\right)+\frac2T\sum_{n=1}^{\infty}\left\{\int_{-T/2}^{T/2}f(t)\cos\frac{2\pi}Tn(t-x)\,dt\right\}}

\begin{itemize}
\tightlist
\item
  جهت به دست آوردن تحلیل فوریه برای تابع \EMim{f(x)} لازم است \EMim{T} را به سمت \textbf{بی‌نهایت} میل دهیم.
\item
  فرض کنید تابع \EMim{f(x)} یک تابع انتگرال‌پذیر است و داریم:
  \EMdisp{\lim_{T\to\infty}\frac1T\int_{-T/2}^{T/2}f(t)\,dt\EMbr =0}
\end{itemize}

در این صورت خواهیم داشت:
\EMdisp{f(x)\EMbr =\lim_{T\to\infty}\frac2T\sum_{n=1}^{\infty}\left\{\int_{-T/2}^{T/2}f(t)\cos\frac{2\pi}Tn(t-x)\,dt\right\}}
\EMdisp{\omega_n\triangleq\frac{2\pi}Tn\quad\EMbr \EMbr \Longrightarrow\quad\EMbr \Delta\omega\triangleq\omega_n-\omega_{n-1}\EMbr =\frac{2\pi}Tn-\frac{2\pi}T(n-1)\EMbr =\frac{2\pi}T\quad\EMbr \EMbr \Longrightarrow\quad\EMbr \frac2T\EMbr =\frac{\Delta\omega}\pi}
\EMdisp{f(x)\EMbr =\lim_{T\to\infty}\frac1\pi\sum_{n=1}^{\infty}\left\{\int_{-T/2}^{T/2}f(t)\cos\omega_n(t-x)\,dt\right\}\Delta\omega}
اگر \EMim{T\to\infty}، آنگاه \EMim{\sum_{n=1}^{\infty}\to\int_0^\infty}، \EMim{\Delta\omega\to d\omega} و \EMim{\omega_n\to\omega}:
\EMdisp{\begin{aligned}
f(x)&=\frac1\pi\int_0^\infty\left\{\int_{-\infty}^{\infty}f(t)\cos\omega(t-x)\,dt\right\}d\omega\\
&=\frac1\pi\int_0^\infty\left\{\int_{-\infty}^{\infty}f(t)\big(\cos\omega t\cos\omega x+\sin\omega t\sin\omega x\big)\,dt\right\}d\omega\\
&=\frac1\pi\int_0^\infty\left\{\cos\omega x\int_{-\infty}^{\infty}\big(f(t)\cos\omega t\big)\,dt+\sin\omega x\int_{-\infty}^{\infty}\big(f(t)\sin\omega t\big)\,dt\right\}d\omega
\end{aligned}}
\EMdisp{A(\omega)\triangleq\int_{-\infty}^{\infty}f(t)\cos\omega t\,dt,\qquad\EMbr  B(\omega)\triangleq\int_{-\infty}^{\infty}f(t)\sin\omega t\,dt\quad\EMbr \EMbr \Longrightarrow\quad\EMbr  f(x)\EMbr =\frac1\pi\int_0^\infty\big\{A(\omega)\cos\omega x+B(\omega)\sin\omega x\big\}\,d\omega}

\EMhold

\EMsection{انتگرال فوریه}{انتگرال فوریه}

\begin{EMimportant}{انتگرال فوریه و تبدیل‌های سینوسی و کسینوسی}

\EMdisp{f(x)\EMbr =\frac1\pi\int_0^\infty\big\{A(\omega)\cos\omega x+B(\omega)\sin\omega x\big\}\,d\omega\qquad\EMbr \EMt{(انتگرال فوریه؛ رابطهٔ سنتز)}}
\EMdisp{A(\omega)\triangleq\int_{-\infty}^{\infty}f(t)\cos\omega t\,dt\qquad\EMbr \EMt{(تبدیل کسینوسی فوریه)}}
\EMdisp{B(\omega)\triangleq\int_{-\infty}^{\infty}f(t)\sin\omega t\,dt\qquad\EMbr \EMt{(تبدیل سینوسی فوریه)}}
(روابط آنالیز)

\end{EMimportant}

\begin{EMtheorem}{}

شرط کافی وجود انتگرال فوریه برای تابع \EMim{f(x)} آن است که این تابع به صورت تکه‌ای پیوسته بوده، در هر نقطه مشتق راست و چپ وجود داشته و تابع به صورت مطلق انتگرال‌پذیر باشد، یعنی:
\EMdisp{\int_{-\infty}^{\infty}|f(x)|\,dx<\infty}
مقدار انتگرال فوریه در نقاط ناپیوستگی برابر است با متوسط حد چپ و راست تابع \EMim{f(x)} در آن نقاط.

\end{EMtheorem}

\EMsection{صورت مختلط انتگرال فوریه (تبدیل فوریه)}{صورت مختلط انتگرال فوریه (تبدیل فوریه)}

\EMdisp{f(x)\EMbr =\frac1\pi\int_0^\infty\left\{\int_{-\infty}^{\infty}f(t)\cos\omega(t-x)\,dt\right\}d\omega}
عبارت \EMim{\displaystyle\int_{-\infty}^{\infty}f(t)\cos\omega(t-x)\,dt} تابعی زوج بر حسب \EMim{\omega} است، بنابراین
\EMdisp{f(x)\EMbr =\frac1{2\pi}\int_{-\infty}^{\infty}\left\{\int_{-\infty}^{\infty}f(t)\cos\omega(t-x)\,dt\right\}d\omega}
و عبارت \EMim{\displaystyle\int_{-\infty}^{\infty}f(t)\sin\omega(t-x)\,dt} تابعی فرد بر حسب \EMim{\omega} است، بنابراین
\EMdisp{\frac{-i}{2\pi}\int_{-\infty}^{\infty}\left\{\int_{-\infty}^{\infty}f(t)\sin\omega(t-x)\,dt\right\}d\omega\EMbr =0}
\EMdisp{f(x)\EMbr =\frac1{2\pi}\int_{-\infty}^{\infty}\left\{\int_{-\infty}^{\infty}f(t)\cos\omega(t-x)\,dt\right\}d\omega-\frac i{2\pi}\int_{-\infty}^{\infty}\left\{\int_{-\infty}^{\infty}f(t)\sin\omega(t-x)\,dt\right\}d\omega}
\EMdisp{f(x)\EMbr =\frac1{2\pi}\int_{-\infty}^{\infty}\left\{\int_{-\infty}^{\infty}f(t)\big(\cos\omega(t-x)-i\sin\omega(t-x)\big)\,dt\right\}d\omega}
\EMdisp{f(x)\EMbr =\frac1{2\pi}\int_{-\infty}^{\infty}\left\{\int_{-\infty}^{\infty}f(t)e^{-i\omega(t-x)}\,dt\right\}d\omega\EMbr =\frac1{2\pi}\int_{-\infty}^{\infty}\left\{\int_{-\infty}^{\infty}f(t)e^{-i\omega t}\,dt\right\}e^{i\omega x}\,d\omega}

\begin{EMimportant}{تبدیل فوریه}

\EMdisp{f(x)\EMbr =\frac1{2\pi}\int_{-\infty}^{\infty}F(\omega)e^{i\omega x}\,d\omega\qquad\EMbr \EMt{(انتگرال فوریه؛ رابطهٔ سنتز)}}
\EMdisp{F(\omega)\EMbr =\int_{-\infty}^{\infty}f(t)e^{-i\omega t}\,dt\qquad\EMbr \EMt{(تبدیل فوریه؛ رابطهٔ آنالیز)}}

\end{EMimportant}

صورت مختلط سری فوریه (یادآوری):
\EMdisp{f(x)\EMbr =\sum_{n=-\infty}^{\infty}c_ne^{i\frac{2\pi}Tnx}\quad\EMbr \EMt{(سری فوریه؛ رابطهٔ سنتز)},\qquad\EMbr  c_n\EMbr =\frac1T\int_Tf(x)e^{-i\frac{2\pi}Tnx}\,dx\quad\EMbr \EMt{(ضرایب سری فوریه؛ رابطهٔ آنالیز)}}

\EMsection{معرفی تابع \EMim{\operatorname{Si}(u)}}{معرفی تابع operatornameSi(u)}

\EMdisp{\operatorname{Si}(x)\triangleq\int_0^x\frac{\sin w}{w}\,dw}

\begin{itemize}
\tightlist
\item
  فرد؛
\item
  دارای مقادیر مجانبی در بی‌نهایت.
\end{itemize}

\EMfigure{m10-si}{تابع \EMim{\operatorname{Si}(u)} (خط ممتد) و تابع \EMim{\sin u/u}
(خط‌چین)}

\begin{EMtheorem}{لم}

\EMdisp{\int_0^\infty\frac{\sin\alpha w}{w}\,dw\EMbr =\begin{cases}\dfrac\pi2&\alpha>0\\\noalign{\vskip 1mm} 0&\alpha=0\\\noalign{\vskip 1mm} -\dfrac\pi2&\alpha<0\end{cases}}

\end{EMtheorem}

\begin{EMproof}

برای \EMim{\alpha=0} درستی لم واضح است. با تغییر متغیر \EMim{\tau=\alpha w}:
\EMdisp{\int_0^x\frac{\sin\alpha w}{w}\,dw\EMbr =\int_0^{\alpha x}\frac{\sin\tau}{\tau/\alpha}\,\frac{d\tau}{\alpha}\EMbr =\int_0^{\alpha x}\frac{\sin\tau}{\tau}\,d\tau\EMbr =\operatorname{Si}(\alpha x),\qquad\EMbr  \operatorname{Si}(\alpha x)\EMbr =-\operatorname{Si}(-\alpha x)}
پس کافی است لم را فقط برای مقادیر مثبت \EMim{\alpha} ثابت کنیم.
\EMdisp{F(\alpha)\triangleq\int_0^\infty e^{-\alpha t}\frac{\sin t}{t}\,dt,\qquad\EMbr  F(0)\EMbr =\operatorname{Si}(+\infty)\EMbr =\;?}
\EMdisp{F'(\alpha)\EMbr =\int_0^\infty-t\,e^{-\alpha t}\frac{\sin t}t\,dt\EMbr =-\int_0^\infty e^{-\alpha t}\sin t\,dt\EMbr =-\left[\frac{e^{-\alpha t}}{1+\alpha^2}\big(-\alpha\sin t-\cos t\big)\right]_0^\infty\EMbr =\frac{-1}{1+\alpha^2}}
\EMdisp{\alpha>0:\quad\EMbr  F(\alpha)\EMbr =-\tan^{-1}\alpha+C,\qquad\EMbr  F(\alpha\to\infty)\EMbr =0\quad\EMbr \EMbr \Longrightarrow\quad\EMbr  C\EMbr =F(\infty)+\tan^{-1}\infty\EMbr =\frac\pi2}
\EMdisp{F(\alpha)\EMbr =-\tan^{-1}\alpha+\frac\pi2\quad\EMbr \EMbr \Longrightarrow\quad\EMbr  F(0)\EMbr =\operatorname{Si}(+\infty)\EMbr =\frac\pi2}

\end{EMproof}

\EMhold

\EMsection{مثال‌های انتگرال فوریه}{مثال‌های انتگرال فوریه}

\begin{EMexample}{}

مطلوب است بیان تابع پالس مستطیلی زیر به صورت انتگرال فوریه.
\EMdisp{f(x)\EMbr =\begin{cases}1&|x|<1\\ 0&|x|>1\end{cases}}

چون تابع \EMim{f(x)} زوج است، تبدیل سینوسی فوریه صفر است: \EMim{B(\omega)=0}.
\EMdisp{A(\omega)\EMbr =\int_{-\infty}^{\infty}f(t)\cos\omega t\,dt\EMbr =\int_{-1}^{1}\cos\omega t\,dt\EMbr =\frac2\omega\sin\omega}
\EMdisp{f(x)\EMbr =\frac1\pi\int_0^\infty\frac2\omega\sin\omega\cos\omega x\,d\omega\EMbr =\frac1\pi\int_0^\infty\frac1\omega\big[\sin(1-x)\omega+\sin(1+x)\omega\big]\,d\omega}
\EMdisp{=\frac1\pi\left[\int_0^\infty\frac{\sin(1-x)\omega}\omega\,d\omega+\int_0^\infty\frac{\sin(1+x)\omega}\omega\,d\omega\right]\EMbr =\begin{cases}0&x<-1\\ 1/2&x=-1\\ 1&-1<x<1\\ 1/2&x=1\\ 0&x>1\end{cases}}
توجه: مقدار انتگرال فوریه در نقاط ناپیوستگی (با استفاده از لم بالا).

\end{EMexample}

\EMfigure{m10-pulse}{پالس مستطیلی \EMim{f(x)=1} برای \EMim{|x|<1}}

\begin{EMexample}{}

مطلوب است بیان تابع نمایی زیر به صورت انتگرال فوریه.
\EMdisp{f(x)\EMbr =\begin{cases}0&x<0\\ e^{-x}&x>0\end{cases}}

\EMdisp{A(\omega)\EMbr =\int_{-\infty}^{\infty}f(x)\cos\omega x\,dx\EMbr =\int_0^\infty e^{-x}\cos\omega x\,dx\EMbr =\frac{e^{-x}}{1+\omega^2}\big(-\cos\omega x+\omega\sin\omega x\big)\Big|_0^\infty\EMbr =\frac1{1+\omega^2}}
\EMdisp{B(\omega)\EMbr =\int_{-\infty}^{\infty}f(x)\sin\omega x\,dx\EMbr =\int_0^\infty e^{-x}\sin\omega x\,dx\EMbr =\frac{e^{-x}}{1+\omega^2}\big(-\sin\omega x-\omega\cos\omega x\big)\Big|_0^\infty\EMbr =\frac\omega{1+\omega^2}}
\EMdisp{f(x)\EMbr =\frac1\pi\int_0^\infty\big\{A(\omega)\cos\omega x+B(\omega)\sin\omega x\big\}\,d\omega\EMbr =\frac1\pi\int_0^\infty\frac{\cos\omega x+\omega\sin\omega x}{1+\omega^2}\,d\omega}
\EMdisp{f(x\EMbr =0)\EMbr =\frac1\pi\int_0^\infty\frac1{1+\omega^2}\,d\omega\EMbr =\frac1\pi\Big[\tan^{-1}\omega\Big]_0^\infty\EMbr =\frac1\pi\times\frac\pi2\EMbr =\frac12}

\end{EMexample}

\begin{EMexample}{}

مطلوب است بیان تابع زیر به صورت انتگرال فوریه.
\EMdisp{f(x)\EMbr =e^{-|x|}}

چون تابع \EMim{f(x)} زوج است، تبدیل سینوسی فوریه صفر است.
\EMdisp{A(\omega)\EMbr =\int_{-\infty}^{\infty}f(x)\cos\omega x\,dx\EMbr =2\int_0^\infty e^{-x}\cos\omega x\,dx\EMbr =\frac{2e^{-x}}{1+\omega^2}\big(-\cos\omega x+\omega\sin\omega x\big)\Big|_0^\infty\EMbr =\frac2{1+\omega^2}}
\EMdisp{B(\omega)\EMbr =\int_{-\infty}^{\infty}e^{-|x|}\sin\omega x\,dx\EMbr =0}
\EMdisp{f(x)\EMbr =\frac1\pi\int_0^\infty\big\{A(\omega)\cos\omega x+B(\omega)\sin\omega x\big\}\,d\omega\EMbr =\frac2\pi\int_0^\infty\frac{\cos\omega x}{1+\omega^2}\,d\omega}
\EMdisp{f(x\EMbr =0)\EMbr =\frac2\pi\int_0^\infty\frac1{1+\omega^2}\,d\omega\EMbr =\frac2\pi\Big[\tan^{-1}\omega\Big]_0^\infty\EMbr =\frac2\pi\times\frac\pi2\EMbr =1}

\end{EMexample}

\EMfigure{m10-exp-abs}{تابع \EMim{f(x)=e^{-|x|}}}

\begin{EMexample}{}

مطلوب است بیان تابع زیر به صورت انتگرال فوریه.
\EMdisp{f(x)\EMbr =\begin{cases}\sin x&|x|\le\pi\\ 0&|x|>\pi\end{cases}}

چون تابع \EMim{f(x)} فرد است، تبدیل کسینوسی فوریه صفر است.
\EMdisp{A(\omega)\EMbr =\int_{-\pi}^{\pi}\sin x\cos\omega x\,dx\EMbr =0}
\EMdisp{B(\omega)\EMbr =\int_{-\pi}^{\pi}\sin x\sin\omega x\,dx\EMbr =2\int_0^\pi\sin x\sin\omega x\,dx\EMbr =\int_0^\pi\big[\cos(1-\omega)x-\cos(1+\omega)x\big]dx\EMbr =\left[\frac{\sin(1-\omega)x}{1-\omega}-\frac{\sin(1+\omega)x}{1+\omega}\right]_0^\pi\EMbr =\frac{2\sin(\omega\pi)}{1-\omega^2}}
\EMdisp{f(x)\EMbr =\frac1\pi\int_0^\infty\big\{A(\omega)\cos\omega x+B(\omega)\sin\omega x\big\}\,d\omega\EMbr =\frac2\pi\int_0^\infty\frac{\sin(\omega\pi)\sin\omega x}{1-\omega^2}\,d\omega}

\end{EMexample}

\EMfigure{m10-sine-pulse}{تابع \EMim{f(x)=\sin x} برای \EMim{|x|\le\pi} و صفر در بیرون}

\begin{EMexample}{}

مطلوب است بیان تابع زیر به صورت انتگرال فوریه.
\EMdisp{f(x)\EMbr =\begin{cases}x^2&|x|\le a\\ 0&|x|>a\end{cases}}

چون تابع \EMim{f(x)} زوج است، تبدیل سینوسی فوریه صفر است.
\EMdisp{B(\omega)\EMbr =\int_{-a}^{a}x^2\sin\omega x\,dx\EMbr =0}
\EMdisp{A(\omega)\EMbr =\int_{-\infty}^{\infty}f(x)\cos\omega x\,dx\EMbr =2\int_0^ax^2\cos\omega x\,dx\EMbr =2\left[\frac{x^2}\omega\sin\omega x+\frac{2x}{\omega^2}\cos\omega x-\frac2{\omega^3}\sin\omega x\right]_0^a\EMbr =2\left[\left(\frac{a^2}\omega-\frac2{\omega^3}\right)\sin\omega a+\frac{2a}{\omega^2}\cos\omega a\right]}
\EMdisp{f(x)\EMbr =\frac2\pi\int_0^\infty\left[\left(\frac{a^2}\omega-\frac2{\omega^3}\right)\sin\omega a+\frac{2a}{\omega^2}\cos\omega a\right]\cos\omega x\,d\omega}

\end{EMexample}

\EMsection{بعضی خواص انتگرال فوریه}{بعضی خواص انتگرال فوریه}

کلیهٔ خواص از روابط سنتز و آنالیز نشأت می‌گیرند:
\EMdisp{f(x)\EMbr =\frac1\pi\int_0^\infty\big\{A(\omega)\cos\omega x+B(\omega)\sin\omega x\big\}\,d\omega,\qquad\EMbr  A(\omega)\triangleq\int_{-\infty}^{\infty}f(x)\cos\omega x\,dx,\qquad\EMbr  B(\omega)\triangleq\int_{-\infty}^{\infty}f(x)\sin\omega x\,dx}

\EMsubsection{تغییر مقیاس زمانی}{تغییر مقیاس زمانی}

\EMdisp{f(x)\to f(ax),\quad\EMbr  a\in\mathbb{R}\quad\EMbr \EMbr \Longrightarrow\quad\EMbr \begin{cases}A(\omega)\to\dfrac1{|a|}A\!\left(\dfrac\omega a\right)\\\noalign{\vskip 2mm} B(\omega)\to\dfrac1{|a|}B\!\left(\dfrac\omega a\right)\end{cases}}

\begin{EMproof}

\EMdisp{A_1(\omega)\EMbr =\int_{-\infty}^{\infty}f(ax)\cos\omega x\,dx\ \ \begin{cases}\underset{u=ax}{\overset{a>0}{=}}\ \displaystyle\int_{-\infty}^{\infty}f(u)\cos\frac\omega au\,\frac{du}a=\frac1aA\!\left(\frac\omega a\right)\\\noalign{\vskip 3mm}\underset{u=ax}{\overset{a<0}{=}}\ \displaystyle\int_{+\infty}^{-\infty}f(u)\cos\frac\omega au\,\frac{du}a=-\frac1a\int_{-\infty}^{\infty}f(u)\cos\frac\omega au\,du=-\frac1aA\!\left(\frac\omega a\right)\end{cases}\ \EMbr \Longrightarrow\ A_1(\omega)\EMbr =\frac1{|a|}A\!\left(\frac\omega a\right)}

\end{EMproof}

\EMsubsection{تساوی پارسوال}{تساوی پارسوال}

\EMdisp{\int_{-\infty}^{\infty}f^2(x)\,dx\EMbr =\frac1\pi\int_0^\infty\big\{A^2(\omega)+B^2(\omega)\big\}\,d\omega}

\begin{EMproof}

\EMdisp{\begin{aligned}
\int_{-\infty}^{\infty}f^2(x)\,dx&=\int_{-\infty}^{\infty}f(x)\left(\frac1\pi\int_0^\infty\big\{A(\omega)\cos\omega x+B(\omega)\sin\omega x\big\}\,d\omega\right)dx\\
&=\frac1\pi\int_0^\infty\left\{A(\omega)\left[\int_{-\infty}^{\infty}f(x)\cos\omega x\,dx\right]+B(\omega)\left[\int_{-\infty}^{\infty}f(x)\sin\omega x\,dx\right]\right\}d\omega\\
&=\frac1\pi\int_0^\infty\big\{A^2(\omega)+B^2(\omega)\big\}\,d\omega
\end{aligned}}

\end{EMproof}

\EMsubsection{مشتق زمانی}{مشتق زمانی}

\EMdisp{f(x)\EMbr =\frac1\pi\int_0^\infty\big\{A(\omega)\cos\omega x+B(\omega)\sin\omega x\big\}\,d\omega\quad\EMbr \EMbr \Longrightarrow\quad\EMbr  f'(x)\EMbr =\frac1\pi\int_0^\infty\big\{\omega B(\omega)\cos\omega x-\omega A(\omega)\sin\omega x\big\}\,d\omega}
(با استفاده از قاعدهٔ مشتق‌گیری از انتگرال:)
\EMdisp{\frac\partial{\partial z}\int_{a(z)}^{b(z)}f(x,z)\,dx\EMbr =\int_{a(z)}^{b(z)}\frac{\partial f(x,z)}{\partial z}\,dx+f\big(b(z),z\big)\frac{db(z)}{dz}-f\big(a(z),z\big)\frac{da(z)}{dz}}

\EMsubsection{مشتق فرکانسی}{مشتق فرکانسی}

\EMdisp{\frac{dA(\omega)}{d\omega}\EMbr =A'(\omega)\EMbr =\int_{-\infty}^{\infty}-xf(x)\sin\omega x\,dx\quad\EMbr \EMbr \Longrightarrow\quad\EMbr \int_{-\infty}^{\infty}\big(xf(x)\big)\sin\omega x\,dx\EMbr =-A'(\omega)}
\EMdisp{\frac{dB(\omega)}{d\omega}\EMbr =B'(\omega)\EMbr =\int_{-\infty}^{\infty}xf(x)\cos\omega x\,dx\quad\EMbr \EMbr \Longrightarrow\quad\EMbr \int_{-\infty}^{\infty}\big(xf(x)\big)\cos\omega x\,dx\EMbr =B'(\omega)}
\EMdisp{xf(x)\EMbr =\frac1\pi\int_0^\infty\big\{B'(\omega)\cos\omega x-A'(\omega)\sin\omega x\big\}\,d\omega}

\begin{EMexercise}{}

تساوی‌های زیر را ثابت کنید:

\begin{enumerate}
\def\labelenumi{\arabic{enumi}.}
\tightlist
\item
  \EMim{\displaystyle\int_0^\infty\frac{\cos xw+w\sin xw}{1+w^2}\,dw=\begin{cases}0&x<0\\ \pi/2&x=0\\ \pi e^{-x}&x>0\end{cases}}
\item
  \EMim{\displaystyle\int_0^\infty\frac{\sin\pi w\,\sin xw}{1-w^2}\,dw=\begin{cases}\dfrac\pi2\sin x&0\le x\le\pi\\\noalign{\vskip 1mm} 0&x>\pi\end{cases}}
\item
  \EMim{\displaystyle\int_0^\infty\frac{1-\cos\pi w}{w}\sin xw\,dw=\begin{cases}\dfrac12\pi&0<x<\pi\\\noalign{\vskip 1mm} 0&x>\pi\end{cases}}
\item
  \EMim{\displaystyle\int_0^\infty\frac{\cos\frac12\pi w}{1-w^2}\cos xw\,dw=\begin{cases}\dfrac12\pi\cos x&0<|x|<\frac12\pi\\\noalign{\vskip 1mm} 0&|x|\ge\frac12\pi\end{cases}}
\item
  \EMim{\displaystyle\int_0^\infty\frac{\sin w-w\cos w}{w^2}\sin xw\,dw=\begin{cases}\dfrac12\pi x&0<x<1\\\noalign{\vskip 1mm} \dfrac14\pi&x=1\\\noalign{\vskip 1mm} 0&x>1\end{cases}}
\end{enumerate}

\end{EMexercise}

\EMpartpage{3}{۳}{معادلات دیفرانسیل با مشتقات جزئی و تبدیل لاپلاس}{Partial Differential Equations and the Laplace Transform}{فصل‌های ۱۱ تا ۱۴}
\EMsetmodule{11}{معادلهٔ موجی یک‌بعدی و روش جداسازی متغیرها}{The One-Dimensional Wave Equation and Separation of Variables}
\EMchapterpage{11}{11}{معادلهٔ موجی یک‌بعدی و روش جداسازی متغیرها}{The One-Dimensional Wave Equation and Separation of Variables}
\begin{EMminitoc}
\EMtocitem{1}{مقدمه}
\EMtocitem{2}{معادلهٔ دیفرانسیل موجی یک‌بعدی}
\EMtocitem{3}{حل معادلهٔ دیفرانسیل موجی یک‌بعدی}
\EMtocitem{4}{روش جداسازی متغیرها: مثال‌ها}
\EMtocitem{5}{معادلهٔ موجی یک‌بعدی با شرایط مرزی غیرصفر}
\EMtocitem{6}{معادلهٔ موجی یک‌بعدی با سرعت اولیهٔ صفر (امواج متحرک)}
\end{EMminitoc}
\EMchapterend
\EMhold

\EMsection{مقدمه}{مقدمه}

\begin{EMdefinition}{معادلهٔ دیفرانسیل}

معادله‌ای که در آن تابع و مشتقات وجود داشته باشند.

\end{EMdefinition}

\begin{EMdefinition}{معادلهٔ دیفرانسیل معمولی}

معادله‌ای که بر حسب فقط یک متغیر مستقل و مشتقات تابع فقط بر حسب آن متغیر نوشته شده باشد.

\end{EMdefinition}

\begin{EMdefinition}{معادلهٔ دیفرانسیل با مشتقات جزئی}

معادلهٔ دیفرانسیلی که تابع به چند متغیر مستقل وابسته بوده، مشتقات جزئی تابع بر حسب متغیرهای مختلف در معادله وجود داشته باشد.

\end{EMdefinition}

\textbf{معادلهٔ دیفرانسیل موج:}
\EMdisp{\frac{\partial^2u}{\partial t^2}\EMbr =c^2\frac{\partial^2u}{\partial x^2},\qquad\EMbr  \frac{\partial^2u}{\partial t^2}\EMbr =c^2\left(\frac{\partial^2u}{\partial x^2}+\frac{\partial^2u}{\partial y^2}\right),\qquad\EMbr  \frac{\partial^2u}{\partial t^2}\EMbr =c^2\left(\frac{\partial^2u}{\partial x^2}+\frac{\partial^2u}{\partial y^2}+\frac{\partial^2u}{\partial z^2}\right)}

\textbf{معادلهٔ دیفرانسیل انتقال حرارت:}
\EMdisp{\frac{\partial u}{\partial t}\EMbr =c^2\frac{\partial^2u}{\partial x^2},\qquad\EMbr  \frac{\partial u}{\partial t}\EMbr =c^2\left(\frac{\partial^2u}{\partial x^2}+\frac{\partial^2u}{\partial y^2}\right),\qquad\EMbr  \frac{\partial u}{\partial t}\EMbr =c^2\left(\frac{\partial^2u}{\partial x^2}+\frac{\partial^2u}{\partial y^2}+\frac{\partial^2u}{\partial z^2}\right)}

\textbf{معادلهٔ دیفرانسیل لاپلاس:}
\EMdisp{\frac{\partial^2u}{\partial x^2}+\frac{\partial^2u}{\partial y^2}\EMbr =0,\qquad\EMbr  \frac{\partial^2u}{\partial x^2}+\frac{\partial^2u}{\partial y^2}+\frac{\partial^2u}{\partial z^2}\EMbr =0,\qquad\EMbr  \nabla^2u\EMbr =0}

\EMsection{معادلهٔ دیفرانسیل موجی یک‌بعدی}{معادلهٔ دیفرانسیل موجی یک‌بعدی}

مدلسازی یک تار مرتعش (به‌دست آوردن معادلهٔ حاکم بر \EMim{u(x,t)}).

\textbf{فرضیات فیزیکی:}

\begin{itemize}
\tightlist
\item
  جرم به صورت یکنواخت در طول تار توزیع شده (تار همگن)، کاملاً انعطاف‌پذیر است.
\item
  تار چنان محکم کشیده شده است که نیروی جاذبه در آن قابل چشم‌پوشی است.
\item
  میزان انحراف تار از حالت تعادل بسیار کوچک است.
\end{itemize}

\EMfigure{m11-string}{قطعهٔ \EMim{PQ} از تار مرتعش: کشش‌های \EMim{T_1} و \EMim{T_2} در دو انتها
با زوایای \EMim{\alpha} و \EMim{\beta}}

\EMdisp{T_1\cos\alpha\EMbr =T_2\cos\beta\EMbr =T\EMbr =\text{const.}}
\EMdisp{T_2\sin\beta-T_1\sin\alpha\EMbr =\rho\,\Delta s\,\frac{\partial^2u}{\partial t^2}}
که در آن \EMim{\rho} چگالی جرمی تار و \EMim{\Delta s} دیفرانسیل طول تار است. با تقسیم بر \EMim{T}:
\EMdisp{T\big(\tan\beta-\tan\alpha\big)\EMbr =\rho\,\Delta s\,\frac{\partial^2u}{\partial t^2}\quad\EMbr \EMbr \Longrightarrow\quad\EMbr  T\left(\frac{\partial u}{\partial x}\bigg|_{(x+\Delta x,t)}-\frac{\partial u}{\partial x}\bigg|_{(x,t)}\right)\EMbr =\rho\,\Delta s\,\frac{\partial^2u}{\partial t^2}}

با استفاده از قضیهٔ مقدار میانی (\EMim{f(b)-f(a)=f'(l)(b-a)}، \EMim{l\in[a,b]}):
\EMdisp{T\left(\frac{\partial u}{\partial x}\bigg|_{(x+\Delta x,t)}-\frac{\partial u}{\partial x}\bigg|_{(x,t)}\right)\EMbr =T\frac{\partial^2u}{\partial x^2}\bigg|_{(l,t)}\Delta x\EMbr =\rho\,\Delta s\,\frac{\partial^2u}{\partial t^2},\qquad\EMbr  x<l<x+\Delta x}
طبق فرض \EMim{\displaystyle\lim_{\Delta x\to0}\frac{\Delta x}{\Delta s}=1}، پس:
\EMdisp{T\frac{\partial^2u}{\partial x^2}\EMbr =\rho\frac{\partial^2u}{\partial t^2},\qquad\EMbr  c^2\triangleq\frac T\rho\quad\EMbr \EMbr \Longrightarrow\quad\EMbr \frac{\partial^2u}{\partial t^2}\EMbr =c^2\frac{\partial^2u}{\partial x^2}}

برای حل یکتای این معادله نیاز به شرایط مرزی و شرایط اولیه داریم:

\begin{itemize}
\tightlist
\item
  شرایط مرزی (\lr{Boundary Conditions}): \EMim{u(0,t)=u(L,t)=0}
\item
  شرایط اولیه (\lr{Initial Conditions}): \EMim{u(x,0)=f(x)} و \EMim{u_t(x,0)=g(x)}، که در آن \EMim{u_t\triangleq\dfrac{\partial u}{\partial t}}
\end{itemize}

\EMsection{حل معادلهٔ دیفرانسیل موجی یک‌بعدی}{حل معادلهٔ دیفرانسیل موجی یک‌بعدی}

\EMdisp{\begin{cases}\dfrac{\partial^2u}{\partial t^2}=c^2\dfrac{\partial^2u}{\partial x^2}\\\noalign{\vskip 2mm} u(0,t)=u(L,t)=0\\ u(x,0)=f(x)\\ u_t(x,0)=g(x)\end{cases}}

در اینجا برای حل این معادله از روش \textbf{جداسازی متغیرها} استفاده می‌کنیم. فرض:
\EMdisp{u(x,t)\EMbr =X(x)\times T(t)}
\EMdisp{\frac{\partial^2u}{\partial x^2}\EMbr =X''(x)T(t),\qquad\EMbr  \frac{\partial^2u}{\partial t^2}\EMbr =X(x)T''(t)\quad\EMbr \EMbr \Longrightarrow\quad\EMbr  X(x)T''(t)\EMbr =c^2X''(x)T(t)\quad\EMbr \EMbr \Longrightarrow\quad\EMbr \frac{T''}{c^2T}\EMbr =\frac{X''}{X}\EMbr =k}

\EMim{k} باید حتماً یک عدد منفی باشد. بررسی حالت‌ها:

\textbf{اگر \EMim{k=0}:}
\EMdisp{\frac{X''}{X}\EMbr =0\ \EMbr \Rightarrow\ X''\EMbr =0\ \EMbr \Rightarrow\ X(x)\EMbr =Ax+B,\qquad\EMbr  X(0)\EMbr =0\EMbr \Rightarrow B\EMbr =0,\quad\EMbr  X(L)\EMbr =0\EMbr \Rightarrow A\EMbr =0\ \EMbr \Rightarrow\ X(x)\EMbr =0}

\textbf{اگر \EMim{k>0} (\EMim{k=\lambda^2}):}
\EMdisp{\frac{X''}{X}\EMbr =\lambda^2\ \EMbr \Rightarrow\ X''-\lambda^2X\EMbr =0\ \EMbr \Rightarrow\ X(x)\EMbr =Ae^{\lambda x}+Be^{-\lambda x},\qquad\EMbr  X(0)\EMbr =0\EMbr \Rightarrow A+B\EMbr =0,\quad\EMbr  X(L)\EMbr =0\EMbr \Rightarrow Ae^{\lambda L}+Be^{-\lambda L}\EMbr =0}
\EMdisp{\Longrightarrow\ A\EMbr =B\EMbr =0\ \EMbr \Rightarrow\ X(x)\EMbr =0}

بنابراین \EMim{k} باید حتماً یک عدد منفی باشد (\EMim{k=-\lambda^2}):
\EMdisp{\frac{T''}{c^2T}\EMbr =\frac{X''}{X}\EMbr =-\lambda^2\quad\EMbr \EMbr \Longrightarrow\quad\EMbr \begin{cases}X''+\lambda^2X=0\\ T''+c^2\lambda^2T=0\end{cases}}
\EMdisp{X(x)\EMbr =A\cos\lambda x+B\sin\lambda x,\qquad\EMbr  X(0)\EMbr =0\EMbr \Rightarrow A\EMbr =0}
\EMdisp{X(L)\EMbr =0\ \EMbr \Rightarrow\ B\sin\lambda L\EMbr =0\ \EMbr \Rightarrow\ \lambda L\EMbr =n\pi\ \ (n\in\mathbb{Z})\ \EMbr \Rightarrow\ \lambda\EMbr =\frac{n\pi}L\EMbr =\lambda_n}
\EMdisp{X_n(x)\EMbr =B_n\sin\lambda_nx,\qquad\EMbr  \lambda_n\EMbr =\frac{n\pi}L\ \ (n\in\mathbb{Z})}
\EMdisp{T''+c^2\lambda_n^2T\EMbr =0\ \EMbr \Rightarrow\ T_n(t)\EMbr =C_n\cos(c\lambda_nt)+D_n\sin(c\lambda_nt)}
\EMdisp{u_n(x,t)\EMbr =X_n(x)\times T_n(t)\EMbr =B_n\sin\lambda_nx\times\big(C_n\cos(c\lambda_nt)+D_n\sin(c\lambda_nt)\big)}
با برهم‌نهی جواب‌ها:
\EMdisp{u(x,t)\EMbr =\sum_{n=1}^{\infty}\left[a_n\cos\!\left(\frac{cn\pi}Lt\right)+b_n\sin\!\left(\frac{cn\pi}Lt\right)\right]\sin\frac{n\pi}Lx}

\textbf{اعمال شرایط اولیه.} شرط \EMim{u(x,0)=f(x)}:
\EMdisp{u(x,0)\EMbr =\sum_{n=1}^{\infty}a_n\sin\frac{n\pi}Lx\EMbr =f(x)}
تابع \EMim{f(x)} با یک تابع تناوبی و فرد برابر شده است، بنابراین برای محاسبهٔ ضرایب مجهول فوریه باید \EMim{f(x)} را یک تابع تناوبی (با دورهٔ تناوب \EMim{2L}) و فرد فرض کنیم. به این صورت خواهیم داشت:
\EMdisp{a_n\EMbr =\frac2{2L}\int_{-L}^{L}f(x)\sin\!\left(\frac{2\pi}{2L}nx\right)dx\EMbr =\frac1L\int_{-L}^{L}f(x)\sin\!\left(\frac{n\pi}Lx\right)dx\EMbr =\frac2L\int_0^Lf(x)\sin\!\left(\frac{n\pi}Lx\right)dx}
شرط \EMim{u_t(x,0)=g(x)}:
\EMdisp{u_t(x,t)\EMbr =\sum_{n=1}^{\infty}\left[-a_n\frac{cn\pi}L\sin\!\left(\frac{cn\pi}Lt\right)+b_n\frac{cn\pi}L\cos\!\left(\frac{cn\pi}Lt\right)\right]\sin\frac{n\pi}Lx}
\EMdisp{u_t(x,0)\EMbr =\sum_{n=1}^{\infty}b_n\frac{cn\pi}L\sin\frac{n\pi}Lx\EMbr =g(x)\quad\EMbr \EMbr \Longrightarrow\quad\EMbr  b_n\EMbr =\frac2{cn\pi}\int_0^Lg(x)\sin\!\left(\frac{n\pi}Lx\right)dx}

\begin{EMimportant}{جواب معادلهٔ موجی یک‌بعدی (روش جداسازی متغیرها)}

\EMdisp{u(x,t)\EMbr =\sum_{n=1}^{\infty}\left[a_n\cos\!\left(\frac{cn\pi}Lt\right)+b_n\sin\!\left(\frac{cn\pi}Lt\right)\right]\sin\frac{n\pi}Lx}
\EMdisp{a_n\EMbr =\frac2L\int_0^Lf(x)\sin\!\left(\frac{n\pi}Lx\right)dx,\qquad\EMbr  b_n\EMbr =\frac2{cn\pi}\int_0^Lg(x)\sin\!\left(\frac{n\pi}Lx\right)dx}

\end{EMimportant}

\begin{EMexample}{مثال ۱}

جواب معادلهٔ موجی یک‌بعدی زیر را به‌دست آورید.
\EMdisp{\frac{\partial^2u}{\partial t^2}\EMbr =\frac{\partial^2u}{\partial x^2},\qquad\EMbr  u(0,t)\EMbr =u(\pi,t)\EMbr =0,\qquad\EMbr  u(x,0)\EMbr =k\sin2x,\qquad\EMbr  u_t(x,0)\EMbr =0}

با توجه به اینکه سرعت اولیهٔ تار مرتعش صفر است، \EMim{g(x)=0\Rightarrow b_n=0}. بنابراین (\EMim{c=1}، \EMim{L=\pi}):
\EMdisp{u(x,t)\EMbr =\sum_{n=1}^{\infty}a_n\cos\!\left(\frac{cn\pi}Lt\right)\sin\frac{n\pi}Lx\EMbr =\sum_{n=1}^{\infty}a_n\cos nt\,\sin nx}
\EMdisp{a_n\EMbr =\frac2L\int_0^Lf(x)\sin\!\left(\frac{n\pi}Lx\right)dx\EMbr =\frac2\pi\int_0^\pi k\sin2x\sin nx\,dx\EMbr =\begin{cases}0&n\neq2\\ k&n=2\end{cases}}
\EMdisp{u(x,t)\EMbr =k\sin2x\cos2t}

\end{EMexample}

\EMfigure{m11-mode2}{مود دوم ارتعاش تار: \EMim{u(x,t)=k\sin2x\cos2t} در لحظه‌های مختلف}

\begin{EMexample}{مثال ۲}

جواب معادلهٔ موجی یک‌بعدی زیر را به‌دست آورید.
\EMdisp{\frac{\partial^2u}{\partial t^2}\EMbr =\frac{\partial^2u}{\partial x^2},\qquad\EMbr  u(0,t)\EMbr =u(\pi,t)\EMbr =0,\qquad\EMbr  u(x,0)\EMbr =f(x),\qquad\EMbr  u_t(x,0)\EMbr =0}
که در آن \EMim{f(x)} تابع مثلثی شکل زیر است: \EMim{f(x)=\dfrac kax} برای \EMim{0\le x\le a} و \EMim{f(x)=\dfrac k{a-\pi}(x-\pi)} برای \EMim{a\le x\le\pi}.

\EMim{u_t(x,0)=g(x)=0\Rightarrow b_n=0} و
\EMdisp{u(x,t)\EMbr =\sum_{n=1}^{\infty}a_n\cos nt\,\sin nx}
\EMdisp{\begin{aligned}
a_n&=\frac2\pi\int_0^\pi f(x)\sin nx\,dx=\frac2\pi\int_0^a\frac ka\,x\sin nx\,dx+\frac2\pi\int_a^\pi\frac k{a-\pi}(x-\pi)\sin nx\,dx\\
&=\frac{2k}{a\pi}\left[x\left(\frac{-1}n\cos nx\right)-(1)\left(\frac{-1}{n^2}\sin nx\right)\right]_0^a+\frac{2k}{(a-\pi)\pi}\left[(x-\pi)\left(\frac{-1}n\cos nx\right)-(1)\left(\frac{-1}{n^2}\sin nx\right)\right]_a^\pi\\
&=\cdots=\frac{2k\sin na}{an^2(\pi-a)}
\end{aligned}}
\EMdisp{u(x,t)\EMbr =\frac{2k}{a(\pi-a)}\sum_{n=1}^{\infty}\frac{\sin na}{n^2}\cos nt\,\sin nx}

\end{EMexample}

\EMfigure{m11-triangle}{توزیع اولیهٔ مثلثی \EMim{f(x)} با قلهٔ \EMim{k} در \EMim{x=a}}

\begin{EMexample}{مثال ۳}

جواب معادلهٔ موجی یک‌بعدی زیر را به‌دست آورید.
\EMdisp{\frac{\partial^2u}{\partial t^2}\EMbr =\frac{\partial^2u}{\partial x^2},\qquad\EMbr  u(0,t)\EMbr =u(\pi,t)\EMbr =0,\qquad\EMbr  u(x,0)\EMbr =kx(\pi-x),\qquad\EMbr  u_t(x,0)\EMbr =0}

\EMim{g(x)=0\Rightarrow b_n=0} و
\EMdisp{u(x,t)\EMbr =\sum_{n=1}^{\infty}a_n\cos nt\,\sin nx}
\EMdisp{a_n\EMbr =\frac2\pi\int_0^\pi kx(\pi-x)\sin nx\,dx\EMbr =\frac{2k}\pi\Big[(-x^2+\pi x)\Big(\frac{-1}n\cos nx\Big)-(-2x+\pi)\Big(\frac{-1}{n^2}\sin nx\Big)+(-2)\Big(\frac1{n^3}\cos nx\Big)\Big]_0^\pi\EMbr =\cdots\EMbr =\frac{4k\big(1-(-1)^n\big)}{\pi n^3}}
\EMdisp{u(x,t)\EMbr =\frac{4k}\pi\sum_{n=1}^{\infty}\frac{1-(-1)^n}{n^3}\cos nt\,\sin nx}

\end{EMexample}

\EMfigure{m11-parabola}{توزیع اولیهٔ سهمی‌شکل \EMim{f(x)=kx(\pi-x)}}

\begin{EMexercise}{}

(الف) مقدار نسبت \EMim{\dfrac{a_1^2}{a_1^2+a_2^2+\cdots}} را به دست آورید.

(ب) نسبت \EMim{a_{2n-1}/a_{2n+1}} را بر حسب \EMim{n\in\mathbb{N}} رسم کنید.

\end{EMexercise}

\EMhold

\EMsection{روش جداسازی متغیرها: مثال‌ها}{روش جداسازی متغیرها: مثال‌ها}

\begin{EMexample}{مثال ۱}

مطلوب است حل معادلهٔ دیفرانسیل مقابل به روش جداسازی متغیرها: \EMim{xu_x=yu_y}.

\EMdisp{x\frac{\partial u}{\partial x}\EMbr =y\frac{\partial u}{\partial y},\qquad\EMbr  u(x,y)\EMbr =X(x)\cdot Y(y)}
\EMdisp{xX'Y\EMbr =yXY'\ \EMbr \Rightarrow\ \frac{xX'}X\EMbr =\frac{yY'}Y\EMbr =k}
\EMdisp{\frac{X'}X\EMbr =\frac kx\ \EMbr \Rightarrow\ \ln X\EMbr =k\ln x+c_1\ \EMbr \Rightarrow\ X(x)\EMbr =C_1x^k,\qquad\EMbr  \frac{Y'}Y\EMbr =\frac ky\ \EMbr \Rightarrow\ \ln Y\EMbr =k\ln y+c_2\ \EMbr \Rightarrow\ Y(y)\EMbr =C_2y^k}
\EMdisp{u(x,y)\EMbr =X(x)Y(y)\EMbr =C_1x^kC_2y^k\EMbr =C(xy)^k}

\end{EMexample}

\begin{EMexample}{مثال ۲}

مطلوب است حل معادلهٔ دیفرانسیل مقابل به روش جداسازی متغیرها: \EMim{u_{xx}+u_{yy}=0}.

\EMdisp{u(x,y)\EMbr =X(x)\cdot Y(y)\ \EMbr \Rightarrow\ X''Y+XY''\EMbr =0\ \EMbr \Rightarrow\ \frac{X''}X\EMbr =-\frac{Y''}Y\EMbr =k}
\EMdisp{\begin{cases}X''-kX=0\ \Rightarrow\ X(x)=Ae^{\sqrt kx}+Be^{-\sqrt kx}\\ Y''+kY=0\ \Rightarrow\ Y(y)=Ce^{\sqrt{-k}\,y}+De^{-\sqrt{-k}\,y}\end{cases}}
\EMdisp{u(x,y)\EMbr =X(x)Y(y)\EMbr =\big(Ae^{\sqrt kx}+Be^{-\sqrt kx}\big)\big(Ce^{\sqrt{-k}\,y}+De^{-\sqrt{-k}\,y}\big)}

\end{EMexample}

\begin{EMexample}{مثال ۳}

مطلوب است حل معادلهٔ دیفرانسیل مقابل به روش جداسازی متغیرها: \EMim{x^2u_{xy}=-3y^2u}.

\EMdisp{x^2\frac{\partial^2u}{\partial x\partial y}+3y^2u\EMbr =0,\qquad\EMbr  u(x,y)\EMbr =X(x)\cdot Y(y)\ \EMbr \Rightarrow\ x^2X'Y'+3y^2XY\EMbr =0\ \EMbr \Rightarrow\ \frac{x^2X'}X\EMbr =\frac{-3y^2Y}{Y'}\EMbr =k}
\EMdisp{\frac{X'}X\EMbr =\frac k{x^2}\ \EMbr \Rightarrow\ \ln X\EMbr =\frac{-k}x+c_1\ \EMbr \Rightarrow\ X(x)\EMbr =C_1e^{-k/x}}
\EMdisp{\frac{Y'}Y\EMbr =\frac{-3y^2}k\ \EMbr \Rightarrow\ \ln Y\EMbr =-\frac{y^3}k+c_2\ \EMbr \Rightarrow\ Y(y)\EMbr =C_2e^{-y^3/k}}
\EMdisp{u(x,y)\EMbr =X(x)Y(y)\EMbr =C_1e^{-k/x}\times C_2e^{-y^3/k}\EMbr =Ce^{-(k/x+y^3/k)}}

\end{EMexample}

\EMhold

\EMsection{معادلهٔ موجی یک‌بعدی با شرایط مرزی غیرصفر}{معادلهٔ موجی یک‌بعدی با شرایط مرزی غیرصفر}

\begin{EMexample}{مسئله}

جواب معادلهٔ دیفرانسیل موجی یک‌بعدی زیر را به‌دست آورید.
\EMdisp{\frac{\partial^2u}{\partial t^2}\EMbr =c^2\frac{\partial^2u}{\partial x^2},\qquad\EMbr  u(0,t)\EMbr =P_0,\quad\EMbr  u(L,t)\EMbr =P_1,\quad\EMbr  u(x,0)\EMbr =f(x),\quad\EMbr  u_t(x,0)\EMbr =g(x)}

\EMdisp{u(x,t)\EMbr =w(x,t)+v(x)\quad\EMbr \EMbr \Longrightarrow\quad\EMbr \begin{cases}w(0,t)+v(0)=P_0\\ w(L,t)+v(L)=P_1\end{cases}}
\EMdisp{\frac{\partial^2w}{\partial t^2}\EMbr =c^2\left(\frac{\partial^2w}{\partial x^2}+\frac{\partial^2v}{\partial x^2}\right)}
تابع \EMim{v(x)} را چنان انتخاب می‌کنیم که مسئلهٔ \EMim{w} شرایط مرزی همگن داشته باشد:
\EMdisp{\begin{cases}\dfrac{\partial^2v}{\partial x^2}=0\\ v(0)=P_0\\ v(L)=P_1\end{cases}\quad\EMbr \EMbr \Longrightarrow\quad\EMbr  v(x)\EMbr =\left(\frac{P_1-P_0}L\right)x+P_0}
\EMdisp{\begin{cases}\dfrac{\partial^2w}{\partial t^2}=c^2\dfrac{\partial^2w}{\partial x^2}\\ w(0,t)=w(L,t)=0\end{cases}\quad\EMbr \EMbr \Longrightarrow\quad\EMbr  w(x,t)\EMbr =\sum_{n=1}^{\infty}\left[A_n\cos\!\left(\frac{cn\pi}Lt\right)+B_n\sin\!\left(\frac{cn\pi}Lt\right)\right]\sin\frac{n\pi}Lx}
\EMdisp{u(x,t)\EMbr =\underbrace{\sum_{n=1}^{\infty}\left[A_n\cos\!\left(\frac{cn\pi}Lt\right)+B_n\sin\!\left(\frac{cn\pi}Lt\right)\right]\sin\frac{n\pi}Lx}_{w(x,t)}+\underbrace{\left(\frac{P_1-P_0}L\right)x+P_0}_{v(x)}}
شرط اولیهٔ \EMim{u(x,0)=f(x)}:
\EMdisp{\sum_{n=1}^{\infty}A_n\sin\frac{n\pi}Lx\EMbr =f(x)-\left(\frac{P_1-P_0}L\right)x-P_0\triangleq f_1(x)}
\EMdisp{A_n\EMbr =\frac2L\int_0^Lf_1(x)\sin\frac{n\pi}Lx\,dx\EMbr =\frac2L\int_0^L\left[f(x)-\left(\frac{P_1-P_0}L\right)x-P_0\right]\sin\frac{n\pi}Lx\,dx}

\end{EMexample}

\begin{EMexercise}{}

\EMim{B_n} را برای حالت (الف) \EMim{g(x)=0} و (ب) \EMim{g(x)} دلخواه به دست آورید.

\end{EMexercise}

\EMsection{معادلهٔ موجی یک‌بعدی با سرعت اولیهٔ صفر (امواج متحرک)}{معادلهٔ موجی یک‌بعدی با سرعت اولیهٔ صفر (امواج متحرک)}

اگر سرعت اولیه صفر باشد (\EMim{g(x)=0}) آنگاه خواهیم داشت:
\EMdisp{u(x,t)\EMbr =\sum_{n=1}^{\infty}a_n\cos\!\left(\frac{cn\pi}Lt\right)\sin\frac{n\pi}Lx,\qquad\EMbr  u(x,0)\EMbr =\sum_{n=1}^{\infty}a_n\sin\frac{n\pi}Lx\EMbr =f(x)\equiv\tilde f(x)}
با استفاده از \EMim{\sin\alpha\cos\beta=\frac12\big(\sin(\alpha+\beta)+\sin(\alpha-\beta)\big)}:
\EMdisp{\begin{aligned}
u(x,t)&=\sum_{n=1}^{\infty}a_n\frac12\left[\sin\!\left(\frac{n\pi}Lx+\frac{cn\pi}Lt\right)+\sin\!\left(\frac{n\pi}Lx-\frac{cn\pi}Lt\right)\right]\\
&=\frac12\sum_{n=1}^{\infty}a_n\left[\sin\frac{n\pi}L(x+ct)+\sin\frac{n\pi}L(x-ct)\right]\\
&=\frac12\sum_{n=1}^{\infty}a_n\sin\frac{n\pi}L(x+ct)+\frac12\sum_{n=1}^{\infty}a_n\sin\frac{n\pi}L(x-ct)=\frac12\Big[\tilde f(x+ct)+\tilde f(x-ct)\Big]
\end{aligned}}

\begin{EMimportant}{امواج متحرک}

\EMdisp{u(x,t)\EMbr =\frac12\Big[\tilde f(x+ct)+\tilde f(x-ct)\Big]}
\EMim{\tilde f} ادامهٔ فرد و متناوب (با دورهٔ \EMim{2L}) تابع اولیهٔ \EMim{f} است؛ جواب برهم‌نهی دو موج با نیمی از دامنهٔ اولیه است که با سرعت \EMim{c} در دو جهت مخالف حرکت می‌کنند.

\end{EMimportant}

\EMfigure{m11-traveling}{ادامهٔ فرد و متناوب \EMim{\tilde f(x)} و دو موج
\EMim{\frac12\tilde f(x+ct)} و \EMim{\frac12\tilde f(x-ct)} که در دو جهت
مخالف حرکت می‌کنند}

\begin{EMexample}{}

جواب معادلهٔ موجی یک‌بعدی زیر را با فرض سرعت اولیهٔ صفر در زمان‌های مختلف رسم کنید.
\EMdisp{\frac{\partial^2u}{\partial t^2}\EMbr =c^2\frac{\partial^2u}{\partial x^2},\qquad\EMbr  u(0,t)\EMbr =u(L,t)\EMbr =0,\qquad\EMbr  u(x,0)\EMbr =f(x),\qquad\EMbr  u_t(x,0)\EMbr =0}
\EMdisp{f(x)\EMbr =\begin{cases}\dfrac{2k}Lx&0<x<\dfrac L2\\\noalign{\vskip 7mm} \dfrac{2k}L(L-x)&\dfrac L2<x<L\end{cases}}

با رسم \EMim{\frac12\tilde f(x+ct)} و \EMim{\frac12\tilde f(x-ct)} و جمع آنها در زمان‌های \EMim{t=0}، \EMim{t=L/5c}، \EMim{t=2L/5c}، \EMim{t=L/2c} و \EMim{t=3L/5c} شکل جواب به دست می‌آید.

\end{EMexample}

\EMfigure{m11-triangle-time}{جواب مثال برای \EMim{t=0}، \EMim{L/5c}، \EMim{2L/5c}، \EMim{L/2c} و
\EMim{3L/5c}؛ در \EMim{t=L/2c} شکل تار صاف است}

\EMsetmodule{12}{روش دالامبر و معادلهٔ انتقال حرارت یک‌بعدی}{d'Alembert's Method and the One-Dimensional Heat Equation}
\EMchapterpage{12}{12}{روش دالامبر و معادلهٔ انتقال حرارت یک‌بعدی}{d'Alembert's Method and the One-Dimensional Heat Equation}
\begin{EMminitoc}
\EMtocitem{1}{روش دالامبر در حل معادلهٔ موجی یک‌بعدی}
\EMtocitem{2}{معادلهٔ دیفرانسیل انتقال حرارت}
\EMtocitem{3}{حل معادلهٔ دیفرانسیل انتقال حرارت یک‌بعدی}
\EMtocitem{4}{معادلهٔ انتقال حرارت یک‌بعدی بر روی مفتول نامتناهی}
\EMtocitem{5}{تابع خطا}
\end{EMminitoc}
\EMchapterend
\EMsection{روش دالامبر در حل معادلهٔ موجی یک‌بعدی}{روش دالامبر در حل معادلهٔ موجی یک‌بعدی}

\EMdisp{\begin{cases}\dfrac{\partial^2u}{\partial t^2}=c^2\dfrac{\partial^2u}{\partial x^2}\\\noalign{\vskip 2mm} u(0,t)=u(L,t)=0\\ u(x,0)=f(x)\\ u_t(x,0)=g(x)\end{cases}}
با استفاده از تغییر متغیرها به شکل زیر معادلهٔ دیفرانسیل را بازنویسی و حل می‌کنیم.
\EMdisp{u(x,t)\EMbr \Rightarrow u(\varphi,\psi),\qquad\EMbr \begin{cases}\varphi=x+ct\\ \psi=x-ct\end{cases}}
\EMdisp{\frac{\partial u}{\partial x}\EMbr =\frac{\partial u}{\partial\varphi}\frac{\partial\varphi}{\partial x}+\frac{\partial u}{\partial\psi}\frac{\partial\psi}{\partial x}\EMbr =\frac{\partial u}{\partial\varphi}+\frac{\partial u}{\partial\psi}}
\EMdisp{\frac{\partial u}{\partial t}\EMbr =\frac{\partial u}{\partial\varphi}\frac{\partial\varphi}{\partial t}+\frac{\partial u}{\partial\psi}\frac{\partial\psi}{\partial t}\EMbr =c\frac{\partial u}{\partial\varphi}-c\frac{\partial u}{\partial\psi}}
\EMdisp{\frac{\partial^2u}{\partial x^2}\EMbr =\frac\partial{\partial x}\left(\frac{\partial u}{\partial\varphi}+\frac{\partial u}{\partial\psi}\right)\EMbr =\left(\frac{\partial^2u}{\partial\varphi^2}+\frac{\partial^2u}{\partial\varphi\partial\psi}\right)\cdot1+\left(\frac{\partial^2u}{\partial\psi\partial\varphi}+\frac{\partial^2u}{\partial\psi^2}\right)\cdot1\EMbr =\frac{\partial^2u}{\partial\varphi^2}+2\frac{\partial^2u}{\partial\varphi\partial\psi}+\frac{\partial^2u}{\partial\psi^2}}
(با استفاده از \EMim{\dfrac{\partial^2u}{\partial\varphi\partial\psi}=\dfrac{\partial^2u}{\partial\psi\partial\varphi}}.)

\begin{EMexercise}{}

به طریق مشابه تساوی زیر را ثابت کنید:
\EMdisp{\frac{\partial^2u}{\partial t^2}\EMbr =c^2\left(\frac{\partial^2u}{\partial\varphi^2}-2\frac{\partial^2u}{\partial\varphi\partial\psi}+\frac{\partial^2u}{\partial\psi^2}\right)}

\end{EMexercise}

با قرار دادن دو رابطهٔ بالا در معادلهٔ دیفرانسیل موجی یک‌بعدی:
\EMdisp{c^2\left(\frac{\partial^2u}{\partial\varphi^2}-2\frac{\partial^2u}{\partial\varphi\partial\psi}+\frac{\partial^2u}{\partial\psi^2}\right)\EMbr =c^2\left(\frac{\partial^2u}{\partial\varphi^2}+2\frac{\partial^2u}{\partial\varphi\partial\psi}+\frac{\partial^2u}{\partial\psi^2}\right)\quad\EMbr \EMbr \Longrightarrow\quad\EMbr \frac{\partial^2u}{\partial\varphi\partial\psi}\EMbr =0}

\textbf{حل عمومی:}
\EMdisp{\frac{\partial u}{\partial\varphi}\EMbr =\phi_1(\varphi)\ \EMbr \Rightarrow\ u(\varphi,\psi)\EMbr =\int\phi_1(\varphi)\,d\varphi+\Psi(\psi)\EMbr =\Phi(\varphi)+\Psi(\psi)\quad\EMbr \EMbr \Longrightarrow\quad\EMbr  u(x,t)\EMbr =\Phi(x+ct)+\Psi(x-ct)}

\textbf{اعمال شرایط اولیه:}
\EMdisp{u(x,0)\EMbr =\Phi(x)+\Psi(x)\EMbr =f(x)}
\EMdisp{u_t(x,t)\EMbr =\frac{\partial\Phi}{\partial\varphi}\frac{\partial\varphi}{\partial t}+\frac{\partial\Psi}{\partial\psi}\frac{\partial\psi}{\partial t}\EMbr =c\Phi'(\varphi)-c\Psi'(\psi)\EMbr =c\Phi'(x+ct)-c\Psi'(x-ct)}
\EMdisp{u_t(x,0)\EMbr =c\Phi'(x)-c\Psi'(x)\EMbr =c\big[\Phi'(x)-\Psi'(x)\big]\EMbr =g(x)}
\EMdisp{\begin{cases}\Phi(x)+\Psi(x)=f(x)\\ \Phi'(x)-\Psi'(x)=\dfrac1cg(x)\end{cases}\EMbr \Rightarrow\begin{cases}\Phi'(x)+\Psi'(x)=f'(x)\\ \Phi'(x)-\Psi'(x)=\dfrac1cg(x)\end{cases}\EMbr \Rightarrow\begin{cases}\Phi'(x)=\dfrac12f'(x)+\dfrac1{2c}g(x)\\\noalign{\vskip 2mm} \Psi'(x)=\dfrac12f'(x)-\dfrac1{2c}g(x)\end{cases}}
\EMdisp{\begin{cases}\Phi(x)=\dfrac12f(x)+\dfrac1{2c}\displaystyle\int_0^xg(s)\,ds+k_1\\\noalign{\vskip 2mm} \Psi(x)=\dfrac12f(x)-\dfrac1{2c}\displaystyle\int_0^xg(s)\,ds+k_2\end{cases}}
\EMdisp{u(x,t)\EMbr =\frac12\left(f(x+ct)+\frac1c\int_0^{x+ct}g(s)\,ds\right)+\frac12\left(f(x-ct)-\frac1c\int_0^{x-ct}g(s)\,ds\right)+k_1+k_2}

\begin{EMimportant}{فرمول دالامبر}

\EMdisp{u(x,t)\EMbr =\frac12\Big[f(x+ct)+f(x-ct)\Big]+\frac1{2c}\int_{x-ct}^{x+ct}g(s)\,ds+K,\qquad\EMbr  u(0,0)\EMbr =0\ \EMbr \Rightarrow\ K\EMbr =0}
\EMdisp{u(x,t)\EMbr =\frac12\Big[f(x+ct)+f(x-ct)\Big]+\frac1{2c}\int_{x-ct}^{x+ct}g(s)\,ds}

\end{EMimportant}

\begin{EMexercise}{}

ثابت کنید که اگر بخواهیم شرایط مرزی برقرار باشند بایستی در رابطهٔ بالا توابع \EMim{f(x)} و \EMim{g(x)} هر دو تناوبی با دورهٔ تناوب \EMim{2L} بوده و فرد باشند.

\end{EMexercise}

\EMsection{معادلهٔ دیفرانسیل انتقال حرارت}{معادلهٔ دیفرانسیل انتقال حرارت}

\EMdisp{\frac{\partial u}{\partial t}\EMbr =c^2\frac{\partial^2u}{\partial x^2},\qquad\EMbr  \frac{\partial u}{\partial t}\EMbr =c^2\left(\frac{\partial^2u}{\partial x^2}+\frac{\partial^2u}{\partial y^2}\right),\qquad\EMbr  \frac{\partial u}{\partial t}\EMbr =c^2\left(\frac{\partial^2u}{\partial x^2}+\frac{\partial^2u}{\partial y^2}+\frac{\partial^2u}{\partial z^2}\right)}

\EMkeepfig{m12-rod}{3}

\EMsubsection{معادلهٔ دیفرانسیل انتقال حرارت یک‌بعدی}{معادلهٔ دیفرانسیل انتقال حرارت یک‌بعدی}

\EMfigure{m12-rod}{میلهٔ یک‌بعدی \EMim{0\le x\le L} با توزیع دمای \EMim{u(x,t)}}

برای حل یکتای این معادله نیاز به شرایط مرزی و شرایط اولیه داریم:

\begin{itemize}
\tightlist
\item
  شرایط مرزی (\lr{Boundary Conditions}): \EMim{u(0,t)=u(L,t)=0}
\item
  شرایط اولیه (\lr{Initial Conditions}): \EMim{u(x,0)=f(x)}
\end{itemize}

\EMsection{حل معادلهٔ دیفرانسیل انتقال حرارت یک‌بعدی}{حل معادلهٔ دیفرانسیل انتقال حرارت یک‌بعدی}

\EMdisp{\begin{cases}\dfrac{\partial u}{\partial t}=c^2\dfrac{\partial^2u}{\partial x^2}\\\noalign{\vskip 2mm} u(0,t)=u(L,t)=0\\ u(x,0)=f(x)\end{cases}}
در اینجا برای حل این معادله از روش \textbf{جداسازی متغیرها} استفاده می‌کنیم. فرض: \EMim{u(x,t)=X(x)\times T(t)}.
\EMdisp{\frac{\partial^2u}{\partial x^2}\EMbr =X''(x)T(t),\qquad\EMbr  \frac{\partial u}{\partial t}\EMbr =X(x)T'(t)\quad\EMbr \EMbr \Longrightarrow\quad\EMbr  X(x)T'(t)\EMbr =c^2X''(x)T(t)\quad\EMbr \EMbr \Longrightarrow\quad\EMbr \frac{T'}{c^2T}\EMbr =\frac{X''}X\EMbr =k}
\EMim{k} باید حتماً یک عدد منفی باشد (برای \EMim{k=0} و \EMim{k>0} دقیقاً مانند مسئلهٔ موجی \EMim{X(x)=0} به دست می‌آید). با \EMim{k=-\lambda^2}:
\EMdisp{\frac{T'}{c^2T}\EMbr =\frac{X''}X\EMbr =-\lambda^2\quad\EMbr \EMbr \Longrightarrow\quad\EMbr \begin{cases}X''+\lambda^2X=0\\ T'+c^2\lambda^2T=0\end{cases}}
\EMdisp{X(x)\EMbr =A\cos\lambda x+B\sin\lambda x,\quad\EMbr  X(0)\EMbr =0\EMbr \Rightarrow A\EMbr =0,\qquad\EMbr  X(L)\EMbr =0\EMbr \Rightarrow B\sin\lambda L\EMbr =0\EMbr \Rightarrow\lambda\EMbr =\frac{n\pi}L\EMbr =\lambda_n}
\EMdisp{X_n(x)\EMbr =B_n\sin\lambda_nx,\qquad\EMbr  T'+c^2\lambda_n^2T\EMbr =0\ \EMbr \Rightarrow\ T_n(t)\EMbr =C_ne^{-(c\lambda_n)^2t}}
\EMdisp{u_n(x,t)\EMbr =X_n(x)\times T_n(t)\EMbr =B_n\sin\!\left(\frac{n\pi}Lx\right)\times C_ne^{-\left(\frac{cn\pi}L\right)^2t}}
\EMdisp{u(x,t)\EMbr =\sum_{n=1}^{\infty}A_ne^{-(c\lambda_n)^2t}\sin(\lambda_nx),\qquad\EMbr  \lambda_n\EMbr =\frac{n\pi}L}
شرط اولیه \EMim{u(x,0)=f(x)}:
\EMdisp{u(x,0)\EMbr =\sum_{n=1}^{\infty}A_n\sin\lambda_nx\EMbr =f(x)}
تابع \EMim{f(x)} با یک تابع تناوبی و فرد برابر شده است، بنابراین برای محاسبهٔ ضرایب مجهول فوریه باید \EMim{f(x)} را یک تابع تناوبی (با دورهٔ تناوب \EMim{2L}) و فرد فرض کنیم:
\EMdisp{A_n\EMbr =\frac2{2L}\int_{-L}^{L}f(x)\sin\!\left(\frac{2\pi}{2L}nx\right)dx\EMbr =\frac1L\int_{-L}^{L}f(x)\sin\!\left(\frac{n\pi}Lx\right)dx\EMbr =\frac2L\int_0^Lf(x)\sin\!\left(\frac{n\pi}Lx\right)dx}

\begin{EMimportant}{جواب معادلهٔ گرمای یک‌بعدی (روش جداسازی متغیرها)}

\EMdisp{u(x,t)\EMbr =\sum_{n=1}^{\infty}A_ne^{-(c\lambda_n)^2t}\sin(\lambda_nx),\qquad\EMbr  \lambda_n\EMbr =\frac{n\pi}L,\qquad\EMbr  A_n\EMbr =\frac2L\int_0^Lf(x)\sin\!\left(\frac{n\pi}Lx\right)dx}

\end{EMimportant}

\begin{EMexample}{مثال ۱}

جواب معادلهٔ انتقال حرارت یک‌بعدی زیر را به‌دست آورید.
\EMdisp{\frac{\partial u}{\partial t}\EMbr =\frac{\partial^2u}{\partial x^2},\qquad\EMbr  u(0,t)\EMbr =u(10,t)\EMbr =0,\qquad\EMbr  u(x,0)\EMbr =f(x)\EMbr =\begin{cases}x&0\le x<5\\ 10-x&5\le x\le10\end{cases}}

\EMdisp{\begin{aligned}
A_n&=\frac2L\int_0^Lf(x)\sin\!\left(\frac{n\pi}Lx\right)dx=0.2\int_0^5x\sin\!\left(\frac{n\pi}{10}x\right)dx+0.2\int_5^{10}(10-x)\sin\!\left(\frac{n\pi}{10}x\right)dx\\
&=0.2\left[(x)\left(\frac{-10}{n\pi}\cos\frac{n\pi}{10}x\right)-(1)\left(\frac{-100}{n^2\pi^2}\sin\frac{n\pi}{10}x\right)\right]_0^5\\
&\quad+0.2\left[(10-x)\left(\frac{-10}{n\pi}\cos\frac{n\pi}{10}x\right)-(-1)\left(\frac{-100}{n^2\pi^2}\sin\frac{n\pi}{10}x\right)\right]_5^{10}=\frac{40}{n^2\pi^2}\sin\frac{n\pi}2
\end{aligned}}
\EMdisp{u(x,t)\EMbr =\sum_{n=1}^{\infty}\left(\frac{40}{n^2\pi^2}\sin\frac{n\pi}2\right)e^{-\left(\frac{cn\pi}{10}\right)^2t}\sin\!\left(\frac{n\pi}{10}x\right)}
(با \EMim{c=1}.)

\end{EMexample}

\begin{EMexample}{مثال ۲ (شرایط مرزی غیرصفر)}

جواب معادلهٔ انتقال حرارت یک‌بعدی زیر را به‌دست آورید.
\EMdisp{\frac{\partial u}{\partial t}\EMbr =c^2\frac{\partial^2u}{\partial x^2},\qquad\EMbr  u(0,t)\EMbr =T_0,\qquad\EMbr  u(L,t)\EMbr =T_1,\qquad\EMbr  u(x,0)\EMbr =f(x)}

\EMdisp{u(x,t)\EMbr =w(x,t)+v(x)\quad\EMbr \EMbr \Longrightarrow\quad\EMbr \begin{cases}w(0,t)+v(0)=T_0\\ w(L,t)+v(L)=T_1\end{cases}}
\EMdisp{\frac{\partial w}{\partial t}\EMbr =c^2\left(\frac{\partial^2w}{\partial x^2}+\frac{\partial^2v}{\partial x^2}\right)}
\EMdisp{\begin{cases}\dfrac{\partial^2v}{\partial x^2}=0\\ v(0)=T_0\\ v(L)=T_1\end{cases}\EMbr \Rightarrow\ v(x)\EMbr =\left(\frac{T_1-T_0}L\right)x+T_0,\qquad\EMbr \begin{cases}\dfrac{\partial w}{\partial t}=c^2\dfrac{\partial^2w}{\partial x^2}\\ w(0,t)=w(L,t)=0\end{cases}\EMbr \Rightarrow\ w(x,t)\EMbr =\sum_{n=1}^{\infty}A_ne^{-(c\lambda_n)^2t}\sin(\lambda_nx)}
\EMdisp{u(x,t)\EMbr =\underbrace{\sum_{n=1}^{\infty}A_ne^{-(c\lambda_n)^2t}\sin(\lambda_nx)}_{w(x,t)}+\underbrace{\left(\frac{T_1-T_0}L\right)x+T_0}_{v(x)},\qquad\EMbr  \lambda_n\EMbr =\frac{n\pi}L}
شرط اولیه:
\EMdisp{\sum_{n=1}^{\infty}A_n\sin\frac{n\pi}Lx\EMbr =f(x)-\left(\frac{T_1-T_0}L\right)x-T_0\triangleq f_1(x)\ \EMbr \Longrightarrow\ A_n\EMbr =\frac2L\int_0^L\left[f(x)-\left(\frac{T_1-T_0}L\right)x-T_0\right]\sin\frac{n\pi}Lx\,dx}

\end{EMexample}

\begin{EMexample}{مثال ۳ (مرزهای عایق)}

جواب معادلهٔ انتقال حرارت یک‌بعدی زیر را به‌دست آورید.
\EMdisp{\frac{\partial u}{\partial t}\EMbr =c^2\frac{\partial^2u}{\partial x^2},\qquad\EMbr  u_x(0,t)\EMbr =u_x(L,t)\EMbr =0,\qquad\EMbr  u(x,0)\EMbr =f(x)}

برای حل این مسئله از روش \textbf{جداسازی متغیرها} استفاده می‌کنیم: \EMim{u(x,t)=X(x)\times T(t)}.
\EMdisp{\frac{T'}{c^2T}\EMbr =\frac{X''}X\EMbr =-\lambda^2\ \EMbr \Rightarrow\ \begin{cases}X''+\lambda^2X=0\ \Rightarrow\ X(x)=A\cos\lambda x+B\sin\lambda x\\ T'+c^2\lambda^2T=0\end{cases}}
\EMdisp{X'(x)\EMbr =-A\lambda\sin\lambda x+B\lambda\cos\lambda x,\qquad\EMbr  X'(0)\EMbr =0\EMbr \Rightarrow B\EMbr =0,\qquad\EMbr  X'(L)\EMbr =0\EMbr \Rightarrow-A\lambda\sin\lambda L\EMbr =0\EMbr \Rightarrow\lambda L\EMbr =n\pi\EMbr \Rightarrow\lambda_n\EMbr =\frac{n\pi}L}
\EMdisp{X_n(x)\EMbr =A_n\cos\lambda_nx,\qquad\EMbr  T_n(t)\EMbr =C_ne^{-(c\lambda_n)^2t}}
\EMdisp{u(x,t)\EMbr =\sum_{n=1}^{\infty}B_ne^{-(c\lambda_n)^2t}\cos(\lambda_nx),\qquad\EMbr  \lambda_n\EMbr =\frac{n\pi}L}
شرط اولیه \EMim{u(x,0)=\sum_{n=1}^{\infty}B_n\cos\!\left(\dfrac{n\pi}Lx\right)=f(x)}:
\EMdisp{B_n\EMbr =\frac2L\int_0^Lf(x)\cos\!\left(\frac{n\pi}Lx\right)dx}

\end{EMexample}

\EMsection{معادلهٔ انتقال حرارت یک‌بعدی بر روی مفتول نامتناهی}{معادلهٔ انتقال حرارت یک‌بعدی بر روی مفتول نامتناهی}

\EMdisp{\begin{cases}\dfrac{\partial u}{\partial t}=c^2\dfrac{\partial^2u}{\partial x^2}\\\noalign{\vskip 2mm} u(\infty,t)=u(-\infty,t)=0\\ u(x,0)=f(x),\quad-\infty<x<\infty\end{cases}}
با فرض \EMim{u(x,t)=X(x)\times T(t)}:
\EMdisp{X(x)T'(t)\EMbr =c^2X''(x)T(t)\ \EMbr \Rightarrow\ \frac{T'}{c^2T}\EMbr =\frac{X''}X\EMbr =-\lambda^2\ \EMbr \Rightarrow\ \begin{cases}X(x)=A_1(\lambda)\cos\lambda x+B_1(\lambda)\sin\lambda x=X_\lambda(x)\\ T(t)=C_1(\lambda)e^{-c^2\lambda^2t}=T_\lambda(t)\end{cases}}
\EMdisp{u_\lambda(x,t)\EMbr =\big[A(\lambda)\cos\lambda x+B(\lambda)\sin\lambda x\big]e^{-c^2\lambda^2t}\quad\EMbr \EMbr \Longrightarrow\quad\EMbr  u(x,t)\EMbr =\frac1\pi\int_0^\infty\big[A(\lambda)\cos\lambda x+B(\lambda)\sin\lambda x\big]e^{-c^2\lambda^2t}\,d\lambda}
صحت معادله:
\EMdisp{\frac{\partial^2u}{\partial x^2}\EMbr =\frac1\pi\int_0^\infty-\lambda^2\big[A\cos\lambda x+B\sin\lambda x\big]e^{-c^2\lambda^2t}\,d\lambda,\qquad\EMbr  \frac{\partial u}{\partial t}\EMbr =\frac1\pi\int_0^\infty-c^2\lambda^2\big[A\cos\lambda x+B\sin\lambda x\big]e^{-c^2\lambda^2t}\,d\lambda\ \EMbr \Rightarrow\ \frac{\partial u}{\partial t}\EMbr =c^2\frac{\partial^2u}{\partial x^2}\ \EMcheck }
\textbf{اعمال شرط اولیه:} \EMim{u(x,0)=\dfrac1\pi\displaystyle\int_0^\infty\big[A(\lambda)\cos\lambda x+B(\lambda)\sin\lambda x\big]d\lambda=f(x)}، یعنی انتگرال فوریه (فصل ۱۰):
\EMdisp{A(\lambda)\EMbr =\mathfrak{I}_C\{f(x)\}\EMbr =\int_{-\infty}^{\infty}f(x)\cos\lambda x\,dx,\qquad\EMbr  B(\lambda)\EMbr =\mathfrak{I}_S\{f(x)\}\EMbr =\int_{-\infty}^{\infty}f(x)\sin\lambda x\,dx}
\EMdisp{\begin{aligned}
u(x,t)&=\frac1\pi\int_0^\infty\left[\left(\int_{-\infty}^{\infty}f(\nu)\cos\lambda\nu\,d\nu\right)\cos\lambda x+\left(\int_{-\infty}^{\infty}f(\nu)\sin\lambda\nu\,d\nu\right)\sin\lambda x\right]e^{-c^2\lambda^2t}\,d\lambda\\
&=\frac1\pi\int_0^\infty\left\{\int_{-\infty}^{\infty}f(\nu)\big[\cos\lambda\nu\cos\lambda x+\sin\lambda\nu\sin\lambda x\big]d\nu\right\}e^{-c^2\lambda^2t}\,d\lambda\\
&=\frac1\pi\int_0^\infty\left\{\int_{-\infty}^{\infty}f(\nu)\cos\lambda(x-\nu)\,d\nu\right\}e^{-c^2\lambda^2t}\,d\lambda
\end{aligned}}

\begin{EMimportant}{شکل اول جواب}

\EMdisp{u(x,t)\EMbr =\frac1\pi\int_{-\infty}^{\infty}f(\nu)\left\{\int_0^\infty e^{-c^2\lambda^2t}\cos\lambda(x-\nu)\,d\lambda\right\}d\nu}

\end{EMimportant}

\EMhold

\EMsubsection{لم‌های ریاضی}{لم‌های ریاضی}

\begin{EMtheorem}{لم (الف)}

\EMdisp{I\EMbr =\int_0^\infty e^{-x^2}\,dx\EMbr =\frac{\sqrt\pi}2}

\end{EMtheorem}

\begin{EMproof}

\EMdisp{I^2\EMbr =\int_0^\infty e^{-x^2}dx\times\int_0^\infty e^{-y^2}dy\EMbr =\int_0^\infty\!\!\int_0^\infty e^{-(x^2+y^2)}\,dx\,dy\ \overset{\rho^2=x^2+y^2}{=}\ \int_0^{\pi/2}\!\!\int_0^\infty\rho\,e^{-\rho^2}\,d\rho\,d\theta\EMbr =\int_0^{\pi/2}\left[\frac{-1}2e^{-\rho^2}\right]_0^\infty d\theta\EMbr =\frac\pi4}

\end{EMproof}

\begin{EMtheorem}{لم (ب)}

\EMdisp{\int_0^\infty e^{-s^2}\cos(2bs)\,ds\EMbr =\frac{\sqrt\pi}2e^{-b^2}}

\end{EMtheorem}

\begin{EMproof}

\EMdisp{I(b)\EMbr =\int_0^\infty e^{-s^2}\cos(2bs)\,ds\ \EMbr \Rightarrow\ \frac{dI(b)}{db}\EMbr =-\int_0^\infty2se^{-s^2}\sin(2bs)\,ds}
با انتگرال‌گیری جزء به جزء (\EMim{du=-2se^{-s^2}ds}، \EMim{v=\sin2bs}):
\EMdisp{\frac{dI(b)}{db}\EMbr =e^{-s^2}\sin2bs\Big|_0^\infty-2b\int_0^\infty e^{-s^2}\cos2bs\,ds\EMbr =-2b\,I(b)\ \EMbr \Rightarrow\ I(b)\EMbr =ke^{-b^2},\qquad\EMbr  I(0)\EMbr =k\EMbr =\int_0^\infty e^{-s^2}ds\EMbr =\frac{\sqrt\pi}2}

\end{EMproof}

با تغییر متغیر \EMim{c\lambda\sqrt t=s}:
\EMdisp{\int_0^\infty e^{-c^2\lambda^2t}\cos\lambda(x-\nu)\,d\lambda\EMbr =\frac1{c\sqrt t}\int_0^\infty e^{-s^2}\cos\!\left[2\frac{(x-\nu)}{2c\sqrt t}s\right]ds\EMbr =\frac1{c\sqrt t}\frac{\sqrt\pi}2e^{-\left(\frac{x-\nu}{2c\sqrt t}\right)^2}}

\begin{EMimportant}{شکل دوم جواب}

\EMdisp{u(x,t)\EMbr =\frac1{2c\sqrt{\pi t}}\int_{-\infty}^{\infty}f(\nu)\,e^{-\frac{(x-\nu)^2}{4c^2t}}\,d\nu}

\end{EMimportant}

با تغییر متغیر \EMim{z=\dfrac{\nu-x}{2c\sqrt t}} (یعنی \EMim{\nu=x+2c\sqrt t\,z} و \EMim{d\nu=2c\sqrt t\,dz}):

\begin{EMimportant}{شکل سوم جواب}

\EMdisp{u(x,t)\EMbr =\frac1{\sqrt\pi}\int_{-\infty}^{\infty}f\big(x+2c\sqrt t\,z\big)\,e^{-z^2}\,dz}

\end{EMimportant}

\EMsection{تابع خطا}{تابع خطا}

در ریاضی اغلب تابع خطا را به صورت زیر تعریف می‌کنند:
\EMdisp{\operatorname{erf}(x)\triangleq\frac2{\sqrt\pi}\int_0^xe^{-t^2}\,dt}
خواص تابع خطا:

\begin{itemize}
\tightlist
\item
  خاصیت تقارن فرد: \EMim{\operatorname{erf}(-x)=-\operatorname{erf}(x)}
\item
  خاصیت مجانبی: \EMim{\operatorname{erf}(\infty)=1}
\item
  بسط مک‌لورن این تابع به صورت زیر است (با بسط \EMim{e^x} شروع و این تساوی را ثابت کنید):
  \EMdisp{\operatorname{erf}(x)\EMbr =\frac2{\sqrt\pi}\left[x-\frac{x^3}{1!\times3}+\frac{x^5}{2!\times5}-\frac{x^7}{3!\times7}+-\cdots\right]\EMbr =\frac2{\sqrt\pi}\sum_{n=1}^{\infty}(-1)^{n+1}\frac{x^{2n-1}}{(n-1)!\,(2n-1)}}
\end{itemize}

\EMfigure{m12-erf}{نمودار تابع خطا \EMim{\operatorname{erf}(x)} برای \EMim{-2\le x\le2}}

\begin{EMexample}{مثال ۱}

جواب معادلهٔ بالا را با فرض شرط اولیهٔ مقابل به دست آورید.
\EMdisp{f(x)\EMbr =\begin{cases}U_0&|x|<1\\ 0&|x|\ge1\end{cases}}

از شکل دوم جواب استفاده می‌کنیم:
\EMdisp{u(x,t)\EMbr =\frac{U_0}{2c\sqrt{\pi t}}\int_{-1}^{1}e^{-\frac{(x-\nu)^2}{4c^2t}}\,d\nu\ \overset{z=\frac{\nu-x}{2c\sqrt t}}{=}\ \frac{U_0}{\sqrt\pi}\int_{-\frac{1+x}{2c\sqrt t}}^{\frac{1-x}{2c\sqrt t}}e^{-z^2}\,dz}
\EMdisp{=\frac{U_0}{\sqrt\pi}\left\{\int_{-\frac{1+x}{2c\sqrt t}}^{0}e^{-z^2}dz+\int_0^{\frac{1-x}{2c\sqrt t}}e^{-z^2}dz\right\}\EMbr =\frac{U_0}{\sqrt\pi}\frac{\sqrt\pi}2\left\{\frac2{\sqrt\pi}\int_0^{\frac{1+x}{2c\sqrt t}}e^{-z^2}dz+\frac2{\sqrt\pi}\int_0^{\frac{1-x}{2c\sqrt t}}e^{-z^2}dz\right\}}
\EMdisp{u(x,t)\EMbr =\frac{U_0}2\left[\operatorname{erf}\!\left(\frac{1+x}{2c\sqrt t}\right)+\operatorname{erf}\!\left(\frac{1-x}{2c\sqrt t}\right)\right]}

\end{EMexample}

\EMfigure{m12-heat-pulse}{توزیع دما \EMim{u(x,t)} در لحظه‌های \EMim{t=0}، \EMim{\frac18}،
\EMim{\frac12}، \EMim{1}، \EMim{2} و \EMim{8} برای پالس اولیهٔ
\EMim{U_0=100} و \EMim{c=1}}

\begin{EMexample}{مثال ۲}

جواب معادلهٔ بالا را با فرض شرط اولیهٔ مقابل به دست آورید.
\EMdisp{f(x)\EMbr =\begin{cases}1&x>0\\ 0&x\le0\end{cases}}

از شکل دوم جواب استفاده می‌کنیم:
\EMdisp{u(x,t)\EMbr =\frac1{2c\sqrt{\pi t}}\int_0^\infty e^{-\frac{(x-\nu)^2}{4c^2t}}\,d\nu\ \overset{z=\frac{\nu-x}{2c\sqrt t}}{=}\ \frac1{\sqrt\pi}\int_{-\frac x{2c\sqrt t}}^{\infty}e^{-z^2}\,dz\EMbr =\frac1{\sqrt\pi}\left\{\int_{-\frac x{2c\sqrt t}}^{0}e^{-z^2}dz+\int_0^\infty e^{-z^2}dz\right\}}
\EMdisp{u(x,t)\EMbr =\frac12\left[\operatorname{erf}\!\left(\frac x{2c\sqrt t}\right)+1\right]}

\end{EMexample}

\EMsetmodule{13}{معادلات دو‌بعدی موج و گرما: مستطیل و دایره}{Two-Dimensional Wave and Heat Equations: Rectangle and Disc}
\EMchapterpage{13}{13}{معادلات دو‌بعدی موج و گرما: مستطیل و دایره}{Two-Dimensional Wave and Heat Equations: Rectangle and Disc}
\begin{EMminitoc}
\EMtocitem{1}{معادلهٔ موجی دو‌بعدی}
\EMtocitem{2}{معادلهٔ موجی دو‌بعدی با شرایط مرزی مستطیل}
\EMtocitem{3}{معادلهٔ گرمای دو‌بعدی با شرایط مرزی مستطیل}
\EMtocitem{4}{بیان لاپلاسین در مختصات قطبی}
\EMtocitem{5}{معادلهٔ موجی دو‌بعدی با شرایط مرزی دایره}
\end{EMminitoc}
\EMchapterend
\EMsection{معادلهٔ موجی دو‌بعدی}{معادلهٔ موجی دو‌بعدی}

مدلسازی یک غشاء مرتعش (به‌دست آوردن معادلهٔ حاکم بر \EMim{u(x,y,t)}).

\textbf{فرضیات فیزیکی:}

\begin{itemize}
\tightlist
\item
  جرم به صورت یکنواخت در طول غشاء توزیع شده (غشاء همگن)، کاملاً انعطاف‌پذیر است.
\item
  غشاء چنان محکم کشیده شده‌است که نیروی جاذبه در آن قابل چشم‌پوشی است.
\item
  میزان انحراف غشاء از حالت تعادل بسیار کوچک است.
\end{itemize}

\EMfigure{m13-membrane}{المان \EMim{\Delta x\times\Delta y} از غشاء مرتعش و نیروهای کششی
\EMim{T\Delta x} و \EMim{T\Delta y} وارد بر آن}

با فرایند مدلسازی مشابه حالت یک‌بعدی، معادلهٔ ارتعاش غشاء دوبعدی به دست می‌آید:
\EMdisp{\frac{\partial^2u}{\partial t^2}\EMbr =c^2\left(\frac{\partial^2u}{\partial x^2}+\frac{\partial^2u}{\partial y^2}\right)\EMbr =c^2\nabla^2u}
برای حل یکتای این معادله نیاز به شرایط مرزی و شرایط اولیه داریم:

\begin{itemize}
\tightlist
\item
  شرایط مرزی: \EMim{u(x,y,t)=0} روی مرزهای غشاء
\item
  شرایط اولیه: \EMim{u(x,y,0)=f(x,y)} و \EMim{u_t(x,y,0)=g(x,y)}، که در آن \EMim{u_t\triangleq\dfrac{\partial u}{\partial t}}
\end{itemize}

\EMkeepfig{m13-rectangle}{3}

\EMsection{معادلهٔ موجی دو‌بعدی با شرایط مرزی مستطیل}{معادلهٔ موجی دو‌بعدی با شرایط مرزی مستطیل}

\EMfigure{m13-rectangle}{ناحیهٔ مستطیلی \EMim{0\le x\le a}، \EMim{0\le y\le b}}

\EMdisp{\frac{\partial^2u}{\partial t^2}\EMbr =c^2\left(\frac{\partial^2u}{\partial x^2}+\frac{\partial^2u}{\partial y^2}\right)}

\begin{itemize}
\tightlist
\item
  شرایط مرزی: \EMim{u(0,y,t)=u(a,y,t)=0} و \EMim{u(x,0,t)=u(x,b,t)=0}
\item
  شرایط اولیه: \EMim{u(x,y,0)=f(x,y)} و \EMim{u_t(x,y,0)=g(x,y)}
\end{itemize}

\EMsubsection{حل معادله}{حل معادله}

در اینجا برای حل این معادله از روش \textbf{جداسازی متغیرها} استفاده می‌کنیم. فرض: \EMim{u(x,y,t)=X(x)\times Y(y)\times T(t)}.
\EMdisp{\frac{\partial^2u}{\partial x^2}\EMbr =X''YT,\qquad\EMbr \frac{\partial^2u}{\partial y^2}\EMbr =XY''T,\qquad\EMbr \frac{\partial^2u}{\partial t^2}\EMbr =XYT''\quad\EMbr \EMbr \Longrightarrow\quad\EMbr  XYT''\EMbr =c^2\big(X''YT+XY''T\big)}
\EMdisp{\frac{T''}{c^2T}\EMbr =\frac{X''}X+\frac{Y''}Y\EMbr =-\lambda^2\ \EMbr \Rightarrow\ \frac{Y''}Y\EMbr =-\lambda^2-\frac{X''}X\EMbr =-\rho^2\ \EMbr \Rightarrow\ \frac{X''}X\EMbr =-\lambda^2+\rho^2\EMbr =-\mu^2}
\EMdisp{\begin{cases}X''+\mu^2X=0\\ Y''+\rho^2Y=0\\ T''+c^2\lambda^2T=0\end{cases}\qquad\EMbr \lambda^2\EMbr =\mu^2+\rho^2}
\textbf{جواب در راستای \EMim{x}:}
\EMdisp{X(x)\EMbr =A_1\cos\mu x+A_2\sin\mu x,\quad\EMbr  X(0)\EMbr =0\EMbr \Rightarrow A_1\EMbr =0,\quad\EMbr  X(a)\EMbr =0\EMbr \Rightarrow A_2\sin\mu a\EMbr =0\EMbr \Rightarrow\mu a\EMbr =n\pi\EMbr \Rightarrow\mu_n\EMbr =\frac{n\pi}a}
\EMdisp{X_n(x)\EMbr =A_n\sin\mu_nx,\qquad\EMbr  \mu_n\EMbr =\frac{n\pi}a}
\textbf{جواب در راستای \EMim{y}:}
\EMdisp{Y(y)\EMbr =B_1\cos\rho y+B_2\sin\rho y,\quad\EMbr  Y(0)\EMbr =0\EMbr \Rightarrow B_1\EMbr =0,\quad\EMbr  Y(b)\EMbr =0\EMbr \Rightarrow B_2\sin\rho b\EMbr =0\EMbr \Rightarrow\rho b\EMbr =m\pi\EMbr \Rightarrow\rho_m\EMbr =\frac{m\pi}b}
\EMdisp{Y_m(y)\EMbr =B_m\sin\rho_my,\qquad\EMbr  \rho_m\EMbr =\frac{m\pi}b}
\textbf{جواب در راستای زمان:}
\EMdisp{\lambda_{nm}^2\EMbr =\mu_n^2+\rho_m^2\EMbr =\left(\frac{n\pi}a\right)^2+\left(\frac{m\pi}b\right)^2,\qquad\EMbr  T_{nm}(t)\EMbr =C_1\cos(c\lambda_{nm}t)+C_2\sin(c\lambda_{nm}t)}
\EMdisp{u_{nm}(x,y,t)\EMbr =\big[A_{nm}\cos(c\lambda_{nm}t)+B_{nm}\sin(c\lambda_{nm}t)\big]\sin\!\left(\frac{n\pi}ax\right)\sin\!\left(\frac{m\pi}by\right),\qquad\EMbr  n,m\EMbr =1,2,\dots}
\EMdisp{u(x,y,t)\EMbr =\sum_{n=1}^{\infty}\sum_{m=1}^{\infty}\big[A_{nm}\cos(c\lambda_{nm}t)+B_{nm}\sin(c\lambda_{nm}t)\big]\sin\!\left(\frac{n\pi}ax\right)\sin\!\left(\frac{m\pi}by\right)}

\textbf{اعمال شرط اولیه \EMim{u(x,y,0)=f(x,y)}:}
\EMdisp{f(x,y)\EMbr =\sum_{n=1}^{\infty}\sum_{m=1}^{\infty}A_{nm}\sin\!\left(\frac{n\pi}ax\right)\sin\!\left(\frac{m\pi}by\right)\EMbr =\sum_{n=1}^{\infty}K_n(y)\sin\frac{n\pi}ax,\qquad\EMbr  K_n(y)\triangleq\sum_{m=1}^{\infty}A_{nm}\sin\frac{m\pi}by}
\EMdisp{K_n(y)\EMbr =\frac2a\int_0^af(x,y)\sin\frac{n\pi}ax\,dx,\qquad\EMbr  A_{nm}\EMbr =\frac2b\int_0^bK_n(y)\sin\frac{m\pi}by\,dy}
\EMdisp{A_{nm}\EMbr =\frac4{ab}\int_0^b\left\{\int_0^af(x,y)\sin\frac{n\pi}ax\,dx\right\}\sin\frac{m\pi}by\,dy}

\begin{EMexercise}{}

ثابت کنید:
\EMdisp{B_{nm}\EMbr =\frac4{abc\lambda_{nm}}\int_0^b\left\{\int_0^ag(x,y)\sin\frac{n\pi}ax\,dx\right\}\sin\frac{m\pi}by\,dy}

\end{EMexercise}

\begin{EMimportant}{جواب معادلهٔ موجی دو‌بعدی در مستطیل}

\EMdisp{u(x,y,t)\EMbr =\sum_{n=1}^{\infty}\sum_{m=1}^{\infty}\big[A_{nm}\cos(c\lambda_{nm}t)+B_{nm}\sin(c\lambda_{nm}t)\big]\sin\!\left(\frac{n\pi}ax\right)\sin\!\left(\frac{m\pi}by\right),\qquad\EMbr \lambda_{nm}^2\EMbr =\left(\frac{n\pi}a\right)^2+\left(\frac{m\pi}b\right)^2}
\EMdisp{A_{nm}\EMbr =\frac4{ab}\int_0^b\left\{\int_0^af(x,y)\sin\frac{n\pi}ax\,dx\right\}\sin\frac{m\pi}by\,dy,\qquad\EMbr  B_{nm}\EMbr =\frac4{abc\lambda_{nm}}\int_0^b\left\{\int_0^ag(x,y)\sin\frac{n\pi}ax\,dx\right\}\sin\frac{m\pi}by\,dy}

\end{EMimportant}

\begin{EMexample}{}

مطلوب است جواب معادلهٔ حرکت دو‌بعدی بالا با فرض \EMim{g(x,y)=0} و \EMim{f(x,y)=xy}.

\EMdisp{g(x,y)\EMbr =0\ \EMbr \Rightarrow\ B_{nm}\EMbr =0}
\EMdisp{A_{nm}\EMbr =\frac4{ab}\int_0^b\left\{\int_0^axy\sin\frac{n\pi}ax\,dx\right\}\sin\frac{m\pi}by\,dy\EMbr =\frac4{ab}\left\{\left[\int_0^by\sin\frac{m\pi}by\,dy\right]\times\left[\int_0^ax\sin\frac{n\pi}ax\,dx\right]\right\}}
\EMdisp{=\frac4{ab}\left[y\left(\frac{-b}{m\pi}\cos\frac{m\pi}by\right)-(1)\left(\frac{-b^2}{m^2\pi^2}\sin\frac{m\pi}by\right)\right]_0^b\times\left[x\left(\frac{-a}{n\pi}\cos\frac{n\pi}ax\right)-(1)\left(\frac{-a^2}{n^2\pi^2}\sin\frac{n\pi}ax\right)\right]_0^a}
\EMdisp{=\frac4{ab}\cdot\frac{-b^2}{m\pi}(-1)^m\cdot\frac{-a^2}{n\pi}(-1)^n\EMbr =\frac{4ab}{nm\pi^2}(-1)^{n+m}}
\EMdisp{u(x,y,t)\EMbr =\sum_{n=1}^{\infty}\sum_{m=1}^{\infty}\frac{4ab}{nm\pi^2}(-1)^{n+m}\cos(c\lambda_{nm}t)\sin\!\left(\frac{n\pi}ax\right)\sin\!\left(\frac{m\pi}by\right)}

\end{EMexample}

\EMsection{معادلهٔ گرمای دو‌بعدی با شرایط مرزی مستطیل}{معادلهٔ گرمای دو‌بعدی با شرایط مرزی مستطیل}

\EMdisp{\frac{\partial u}{\partial t}\EMbr =c^2\left(\frac{\partial^2u}{\partial x^2}+\frac{\partial^2u}{\partial y^2}\right)\EMbr =c^2\nabla^2u}

\begin{itemize}
\tightlist
\item
  شرط مرزی: \EMim{u(x,y,t)=0} روی مرزهای غشاء
\item
  شرط اولیه: \EMim{u(x,y,0)=f(x,y)}
\end{itemize}

\EMsubsection{حل معادله}{حل معادله}

فرض: \EMim{u(x,y,t)=X(x)\times Y(y)\times T(t)}.
\EMdisp{\frac{\partial^2u}{\partial x^2}\EMbr =X''YT,\qquad\EMbr \frac{\partial^2u}{\partial y^2}\EMbr =XY''T,\qquad\EMbr \frac{\partial u}{\partial t}\EMbr =XYT'\quad\EMbr \EMbr \Longrightarrow\quad\EMbr  XYT'\EMbr =c^2\big(X''YT+XY''T\big)}
\EMdisp{\frac{T'}{c^2T}\EMbr =\frac{X''}X+\frac{Y''}Y\EMbr =-\lambda^2,\qquad\EMbr  \frac{Y''}Y\EMbr =-\lambda^2-\frac{X''}X\EMbr =-\rho^2,\qquad\EMbr \frac{X''}X\EMbr =-\lambda^2+\rho^2\EMbr =-\mu^2}
\EMdisp{\begin{cases}X''+\mu^2X=0\\ Y''+\rho^2Y=0\\ T'+c^2\lambda^2T=0\end{cases}\qquad\EMbr \lambda^2\EMbr =\mu^2+\rho^2}
مشابه حالت موجی (\EMim{X_n(x)=A_n\sin\mu_nx} با \EMim{\mu_n=n\pi/a} و \EMim{Y_m(y)=B_m\sin\rho_my} با \EMim{\rho_m=m\pi/b}) و با
\EMdisp{T_{nm}(t)\EMbr =C_1e^{-c^2\lambda_{nm}^2t},\qquad\EMbr \lambda_{nm}^2\EMbr =\mu_n^2+\rho_m^2\EMbr =\left(\frac{n\pi}a\right)^2+\left(\frac{m\pi}b\right)^2}
\EMdisp{u_{nm}(x,y,t)\EMbr =A_{nm}\sin\!\left(\frac{n\pi}ax\right)\sin\!\left(\frac{m\pi}by\right)e^{-c^2\lambda_{nm}^2t},\qquad\EMbr  n\EMbr =1,2,\dots,\ m\EMbr =1,2,\dots}
\EMdisp{u(x,y,t)\EMbr =\sum_{n=1}^{\infty}\sum_{m=1}^{\infty}A_{nm}\sin\!\left(\frac{n\pi}ax\right)\sin\!\left(\frac{m\pi}by\right)e^{-c^2\lambda_{nm}^2t}}
\textbf{اعمال شرط اولیه:} دقیقاً مانند حالت موجی
\EMdisp{A_{nm}\EMbr =\frac4{ab}\int_0^b\left\{\int_0^af(x,y)\sin\frac{n\pi}ax\,dx\right\}\sin\frac{m\pi}by\,dy}

\EMkeepfig{m13-polar}{5}

\EMsection{بیان لاپلاسین در مختصات قطبی}{بیان لاپلاسین در مختصات قطبی}

هدف آن است که \EMim{\nabla^2u=u_{xx}+u_{yy}} در مختصات قطبی بر حسب \EMim{r,\theta} نوشته شود.

\EMfigure{m13-polar}{مختصات قطبی: \EMim{x=r\cos\theta}، \EMim{y=r\sin\theta}}

\EMdisp{r\EMbr =\sqrt{x^2+y^2},\quad\EMbr  x\EMbr =r\cos\theta,\qquad\EMbr \theta\EMbr =\tan^{-1}\!\left(\frac yx\right),\quad\EMbr  y\EMbr =r\sin\theta}
\EMdisp{u_x\EMbr =u_rr_x+u_\theta\theta_x\ \EMbr \Rightarrow\ u_{xx}\EMbr =\big(u_rr_x+u_\theta\theta_x\big)_x\EMbr =u_{rr}r_x^2+u_{\theta\theta}\theta_x^2+2r_x\theta_xu_{r\theta}+r_{xx}u_r+\theta_{xx}u_\theta}
\EMdisp{u_y\EMbr =u_rr_y+u_\theta\theta_y\ \EMbr \Rightarrow\ u_{yy}\EMbr =u_{rr}r_y^2+u_{\theta\theta}\theta_y^2+2r_y\theta_yu_{r\theta}+r_{yy}u_r+\theta_{yy}u_\theta}
\EMdisp{r_x\EMbr =\frac xr,\quad\EMbr  r_{xx}\EMbr =\frac{y^2}{r^3},\quad\EMbr \theta_x\EMbr =\frac{-y}{r^2},\quad\EMbr \theta_{xx}\EMbr =\frac{2xy}{r^4},\qquad\EMbr  r_y\EMbr =\frac yr,\quad\EMbr  r_{yy}\EMbr =\frac{x^2}{r^3},\quad\EMbr \theta_y\EMbr =\frac x{r^2},\quad\EMbr \theta_{yy}\EMbr =\frac{-2xy}{r^4}}
\EMdisp{u_{xx}\EMbr =\frac{x^2}{r^2}u_{rr}+\frac{y^2}{r^4}u_{\theta\theta}+\frac{-2xy}{r^3}u_{r\theta}+\frac{y^2}{r^3}u_r+\frac{2xy}{r^4}u_\theta}
\EMdisp{u_{yy}\EMbr =\frac{y^2}{r^2}u_{rr}+\frac{x^2}{r^4}u_{\theta\theta}+\frac{2xy}{r^3}u_{r\theta}+\frac{x^2}{r^3}u_r+\frac{-2xy}{r^4}u_\theta}
با جمع این دو رابطه:

\begin{EMimportant}{لاپلاسین در مختصات قطبی}

\EMdisp{\nabla^2u\EMbr =u_{rr}+\frac1{r^2}u_{\theta\theta}+\frac1ru_r}

\end{EMimportant}

\EMkeepfig{m13-disc}{3}

\EMsection{معادلهٔ موجی دو‌بعدی با شرایط مرزی دایره}{معادلهٔ موجی دو‌بعدی با شرایط مرزی دایره}

\EMfigure{m13-disc}{غشاء دایره‌ای به شعاع \EMim{R}}

\EMdisp{\frac{\partial^2u}{\partial t^2}\EMbr =c^2\nabla^2u,\qquad\EMbr  u(R,t)\EMbr =0\ \ \EMt{(شرط مرزی)},\qquad\EMbr  u(r,0)\EMbr =f(r),\ \ u_t(r,0)\EMbr =g(r)\ \ \EMt{(شرایط اولیه)}}

\EMsubsection{حل معادله}{حل معادله}

با توجه به تقارن دایروی در این مسئله، بهتر است آن را در مختصات قطبی حل کنیم. بر روی غشاء دایروی، ارتعاش تابعی از \EMim{\theta} نیست، لذا داریم:
\EMdisp{\begin{cases}\dfrac{\partial^2u}{\partial t^2}=c^2\left(u_{rr}+\dfrac1ru_r\right)\\\noalign{\vskip 2mm} u(R,t)=0,\quad u(r,0)=f(r),\quad u_t(r,0)=g(r)\end{cases}}
حال با فرض تفکیک‌پذیری متغیرها \EMim{u(r,t)=P(r)\times T(t)} داریم:
\EMdisp{PT''\EMbr =c^2\left(P''T+\frac1rP'T\right)\ \EMbr \Rightarrow\ \frac{T''}{c^2T}\EMbr =\frac{P''}P+\frac1r\frac{P'}P\EMbr =-\lambda^2\ \EMbr \Rightarrow\ \begin{cases}P''+\dfrac1rP'+\lambda^2P=0\\\noalign{\vskip 1mm} T''+c^2\lambda^2T=0\end{cases}}

\EMkeepfig{m13-bessel}{5}

\EMsubsection{معادلهٔ دیفرانسیل بسل از مرتبهٔ \EMim{n}}{معادلهٔ دیفرانسیل بسل از مرتبهٔ n}

\EMdisp{y''+\frac1xy'+\left(1-\frac{n^2}{x^2}\right)y\EMbr =0\quad\EMbr \EMbr \Longrightarrow\quad\EMbr  y(x)\EMbr =AJ_n(x)+BY_n(x)}
که در آن \EMim{J_n} توابع بسل نوع اول و از مرتبهٔ \EMim{n} و \EMim{Y_n} توابع بسل نوع دوم و از مرتبهٔ \EMim{n} هستند.

\EMfigure{m13-bessel}{توابع بسل نوع اول \EMim{J_0,\dots,J_4} و نوع دوم \EMim{Y_0,\dots,Y_4}
برای \EMim{0<x\le10}}

با تغییر متغیر \EMim{s=\lambda r}:
\EMdisp{P'\EMbr =\frac{dP}{dr}\EMbr =\frac{dP}{ds}\frac{ds}{dr}\EMbr =\lambda\frac{dP}{ds},\qquad\EMbr  P''\EMbr =\lambda^2\frac{d^2P}{ds^2},\qquad\EMbr \frac1r\lambda\EMbr =\frac{\lambda^2}s}
\EMdisp{\lambda^2\frac{d^2P}{ds^2}+\frac1r\lambda\frac{dP}{ds}+\lambda^2P\EMbr =0\ \EMbr \Rightarrow\ \frac{d^2P}{ds^2}+\frac1s\frac{dP}{ds}+P\EMbr =0}
این معادلهٔ دیفرانسیل بسل از مرتبهٔ صفر است (\EMim{n=0})، پس:
\EMdisp{P(r)\EMbr =AJ_0(\lambda r)+BY_0(\lambda r)}
\textbf{اعمال شرایط مرزی:}

\begin{itemize}
\tightlist
\item
  چون \EMim{P(0)\neq\infty} و \EMim{Y_0} در مبدأ بی‌کران است: \EMim{D=0} (ضریب \EMim{Y_0} صفر است) و \EMim{P(r)=CJ_0(\lambda r)}.
\item
  \EMim{P(R)=0} نتیجه می‌دهد \EMim{J_0(\lambda R)=0}، یعنی \EMim{\lambda R=\alpha_m} که در آن \EMim{\alpha_m} ریشه‌های \EMim{J_0} هستند:
  \EMdisp{\lambda\EMbr =\lambda_m\EMbr =\frac{\alpha_m}R,\quad\EMbr  m\EMbr =1,2,\dots\quad\EMbr \EMbr \Longrightarrow\quad\EMbr  P_m(r)\EMbr =C_mJ_0\!\left(\frac{\alpha_m}Rr\right)}
\end{itemize}

\EMfigure{m13-j0-zeros}{تابع \EMim{J_0(x)} و ریشه‌های آن \EMim{\alpha_1,\alpha_2,\alpha_3}}

\EMdisp{T''+c^2\lambda_m^2T\EMbr =0\ \EMbr \Rightarrow\ T_m(t)\EMbr =A_m\cos(c\lambda_mt)+B_m\sin(c\lambda_mt)}

\begin{EMimportant}{جواب معادلهٔ موجی دو‌بعدی در دایره}

\EMdisp{u(r,t)\EMbr =\sum_{m=1}^{\infty}\big(A_m\cos(c\lambda_mt)+B_m\sin(c\lambda_mt)\big)J_0\!\left(\frac{\alpha_m}Rr\right),\qquad\EMbr \lambda_m\EMbr =\frac{\alpha_m}R}

\end{EMimportant}

\begin{EMremark}{تحقیق (تا ۱ نمرهٔ اضافه)}

گزارشی کامل از معادلهٔ دیفرانسیل بسل از مرتبهٔ \EMim{n} و توابع بسل نوع اول و دوم و خواص آنها و نیز حل کامل معادلهٔ دیفرانسیل موجی دوبعدی با شرایط مرزی دایره تهیه کنید (گزارش تایپ شده و در فرمت \lr{Word} باشد).

\end{EMremark}

\EMsetmodule{14}{تبدیل لاپلاس}{The Laplace Transform}
\EMchapterpage{14}{14}{تبدیل لاپلاس}{The Laplace Transform}
\begin{EMminitoc}
\EMtocitem{1}{تبدیل لاپلاس یک‌طرفه}
\EMtocitem{2}{استفاده از تبدیل لاپلاس در حل معادلات دیفرانسیل}
\end{EMminitoc}
\EMchapterend
\EMhold

\EMsection{تبدیل لاپلاس یک‌طرفه}{تبدیل لاپلاس یک‌طرفه}

\begin{EMdefinition}{تبدیل لاپلاس یک‌طرفه}

\EMdisp{F(s)\EMbr =\mathcal{L}\{f(t)\}\triangleq\int_0^\infty f(t)\,e^{-st}\,dt\qquad\EMbr \EMbr \Longleftrightarrow\qquad\EMbr  f(t)\EMbr =\mathcal{L}^{-1}\{F(s)\}}

\end{EMdefinition}

تبدیل لاپلاس یکی از ابزارهای مهم در کاربردهای مختلف مهندسی و ریاضی است. یکی از کاربردهای تبدیل لاپلاس حل معادلات دیفرانسیل است (فصل ۶ کتاب مرجع درس را مطالعه فرمایید).

تبدیل لاپلاس یک عملگر خطی است:
\EMdisp{\mathcal{L}\{af(t)+bg(t)\}\EMbr =a\,\mathcal{L}\{f(t)\}+b\,\mathcal{L}\{g(t)\}}

\begin{EMexample}{مثال ۱}

\EMdisp{f(t)\EMbr =\begin{cases}1&t\ge0\\ 0&t<0\end{cases}\ \EMbr \Longrightarrow\ \mathcal{L}\{f(t)\}\EMbr =F(s)\EMbr =\;?}

\EMdisp{\mathcal{L}\{f(t)\}\EMbr =\int_0^\infty1\times e^{-st}\,dt\EMbr =-\frac1se^{-st}\Big|_0^\infty\EMbr =\frac1s}

\end{EMexample}

\begin{EMexample}{مثال ۲}

\EMdisp{f(t)\EMbr =\begin{cases}e^{at}&t\ge0\\ 0&t<0\end{cases}\ \EMbr \Longrightarrow\ \mathcal{L}\{f(t)\}\EMbr =F(s)\EMbr =\;?}

\EMdisp{\mathcal{L}\{f(t)\}\EMbr =\int_0^\infty e^{at}\times e^{-st}\,dt\EMbr =\int_0^\infty e^{-(s-a)t}\,dt\EMbr =-\frac1{s-a}e^{-(s-a)t}\Big|_0^\infty\EMbr =\frac1{s-a}}

\end{EMexample}

\begin{EMexample}{مثال ۳}

\EMdisp{\mathcal{L}\{\cosh(at)\}\EMbr =\mathcal{L}\left\{\tfrac12e^{at}+\tfrac12e^{-at}\right\}\EMbr =\tfrac12\mathcal{L}\{e^{at}\}+\tfrac12\mathcal{L}\{e^{-at}\}\EMbr =\frac12\frac1{s-a}+\frac12\frac1{s+a}\EMbr =\frac s{s^2-a^2}}

\end{EMexample}

\begin{EMexample}{مثال ۴}

\EMdisp{\mathcal{L}\{\sinh(at)\}\EMbr =\mathcal{L}\left\{\tfrac12e^{at}-\tfrac12e^{-at}\right\}\EMbr =\tfrac12\mathcal{L}\{e^{at}\}-\tfrac12\mathcal{L}\{e^{-at}\}\EMbr =\frac12\frac1{s-a}-\frac12\frac1{s+a}\EMbr =\frac a{s^2-a^2}}

\end{EMexample}

\begin{EMexample}{مثال ۵}

با استفاده از \EMim{\displaystyle\int e^{ax}\cos nx\,dx=\frac{e^{ax}}{a^2+n^2}\big(a\cos nx+n\sin nx\big)}:
\EMdisp{\mathcal{L}\{\cos\omega t\}\EMbr =\int_0^\infty\cos\omega t\,e^{-st}\,dt\EMbr =\frac{e^{-st}}{s^2+\omega^2}\big(-s\cos\omega t+\omega\sin\omega t\big)\Big|_0^\infty\EMbr =\frac s{s^2+\omega^2}}

\end{EMexample}

\begin{EMexample}{مثال ۶}

با استفاده از \EMim{\displaystyle\int e^{ax}\sin nx\,dx=\frac{e^{ax}}{a^2+n^2}\big(a\sin nx-n\cos nx\big)}:
\EMdisp{\mathcal{L}\{\sin\omega t\}\EMbr =\int_0^\infty\sin\omega t\,e^{-st}\,dt\EMbr =\frac{e^{-st}}{s^2+\omega^2}\big(-s\sin\omega t-\omega\cos\omega t\big)\Big|_0^\infty\EMbr =\frac\omega{s^2+\omega^2}}

\end{EMexample}

\EMsubsection{جدول تبدیل‌های لاپلاس (جدول ۶٫۱ کتاب؛ جدول مفصل‌تر در فصل ۶٫۹ کتاب)}{جدول تبدیل‌های لاپلاس (جدول ۶٫۱ کتاب؛ جدول مفصل‌تر در فصل ۶٫۹ کتاب)}

{\def\LTcaptype{none} % do not increment counter
\Needspace{15\baselineskip}
\begin{longtable}[c]{@{}
  >{\centering\arraybackslash}p{(\linewidth - 12\tabcolsep) * \real{0.1515}}
  >{\centering\arraybackslash}p{(\linewidth - 12\tabcolsep) * \real{0.1515}}
  >{\centering\arraybackslash}p{(\linewidth - 12\tabcolsep) * \real{0.1515}}
  >{\raggedright\arraybackslash}p{(\linewidth - 12\tabcolsep) * \real{0.0909}}
  >{\centering\arraybackslash}p{(\linewidth - 12\tabcolsep) * \real{0.1515}}
  >{\centering\arraybackslash}p{(\linewidth - 12\tabcolsep) * \real{0.1515}}
  >{\centering\arraybackslash}p{(\linewidth - 12\tabcolsep) * \real{0.1515}}@{}}
\toprule\noalign{}
\rowcolor{EMindigoL}\begin{minipage}[b]{\linewidth}\centering
\end{minipage} & \begin{minipage}[b]{\linewidth}\centering
\EMim{f(t)}
\end{minipage} & \begin{minipage}[b]{\linewidth}\centering
\EMim{\mathcal{L}(f)}
\end{minipage} & \begin{minipage}[b]{\linewidth}\raggedright
\end{minipage} & \begin{minipage}[b]{\linewidth}\centering
\end{minipage} & \begin{minipage}[b]{\linewidth}\centering
\EMim{f(t)}
\end{minipage} & \begin{minipage}[b]{\linewidth}\centering
\EMim{\mathcal{L}(f)}
\end{minipage} \\
\midrule\noalign{}
\endhead
\bottomrule\noalign{}
\endlastfoot
1 & \EMim{1} & \EMim{1/s} & & 7 & \EMim{\cos\omega t} & \EMim{\dfrac{s}{s^2+\omega^2}} \\
2 & \EMim{t} & \EMim{1/s^2} & & 8 & \EMim{\sin\omega t} & \EMim{\dfrac{\omega}{s^2+\omega^2}} \\
3 & \EMim{t^2} & \EMim{2!/s^3} & & 9 & \EMim{\cosh at} & \EMim{\dfrac{s}{s^2-a^2}} \\
4 & \EMim{t^n} \EMim{(n=0,1,\dots)} & \EMim{\dfrac{n!}{s^{n+1}}} & & 10 & \EMim{\sinh at} & \EMim{\dfrac{a}{s^2-a^2}} \\
5 & \EMim{t^a} (\EMim{a} مثبت) & \EMim{\dfrac{\Gamma(a+1)}{s^{a+1}}} & & 11 & \EMim{e^{at}\cos\omega t} & \EMim{\dfrac{s-a}{(s-a)^2+\omega^2}} \\
6 & \EMim{e^{at}} & \EMim{\dfrac1{s-a}} & & 12 & \EMim{e^{at}\sin\omega t} & \EMim{\dfrac{\omega}{(s-a)^2+\omega^2}} \\
\end{longtable}
}

\EMdisp{\Gamma(x)\EMbr =\int_0^\infty e^{-t}t^{x-1}\,dt}

\EMsubsection{خواص تبدیل لاپلاس یک‌طرفه (جدول فصل ۶٫۸ کتاب)}{خواص تبدیل لاپلاس یک‌طرفه (جدول فصل ۶٫۸ کتاب)}

با تعریف تابع پلهٔ واحد
\EMdisp{u(t)\triangleq\begin{cases}1&t\ge0\\ 0&t<0\end{cases}}

{\def\LTcaptype{none} % do not increment counter
\Needspace{22\baselineskip}
\begin{longtable}[c]{@{}
  >{\raggedright\arraybackslash}p{(\linewidth - 2\tabcolsep) * \real{0.4444}}
  >{\centering\arraybackslash}p{(\linewidth - 2\tabcolsep) * \real{0.5556}}@{}}
\toprule\noalign{}
\rowcolor{EMindigoL}\begin{minipage}[b]{\linewidth}\raggedright
عنوان خاصیت
\end{minipage} & \begin{minipage}[b]{\linewidth}\centering
خاصیت
\end{minipage} \\
\midrule\noalign{}
\endhead
\bottomrule\noalign{}
\endlastfoot
تعریف تبدیل لاپلاس & \EMim{F(s)=\mathcal{L}\{f(t)\}\triangleq\displaystyle\int_0^\infty f(t)e^{-st}\,dt} \\
معکوس تبدیل لاپلاس & \EMim{f(t)=\mathcal{L}^{-1}\{F(s)\}} \\
خطی بودن & \EMim{\mathcal{L}\{af(t)+bg(t)\}=a\mathcal{L}\{f(t)\}+b\mathcal{L}\{g(t)\}} \\
شیفت زمانی & \EMim{\mathcal{L}\{f(t-a)u(t-a)\}=e^{-as}F(s)} \\
شیفت در حوزهٔ \EMim{s} & \EMim{\mathcal{L}\{e^{at}f(t)\}=F(s-a)} \\
تغییر مقیاس زمانی & \EMim{\mathcal{L}\{f(at)\}=\dfrac1aF(s/a)} \\
مشتق زمانی & \EMim{\mathcal{L}\{f^{(n)}(t)\}=s^nF(s)-s^{n-1}f(0)-s^{n-2}f^{(1)}(0)-\cdots-f^{(n-1)}(0)} \\
انتگرال زمانی & \EMim{\mathcal{L}\left\{\displaystyle\int_0^tf(\tau)\,d\tau\right\}=\dfrac1sF(s)} \\
مشتق در حوزهٔ \EMim{s} & \EMim{\mathcal{L}\{tf(t)\}=-F'(s)} \\
انتگرال در حوزهٔ \EMim{s} & \EMim{\mathcal{L}\left\{\dfrac{f(t)}t\right\}=\displaystyle\int_s^\infty F(\nu)\,d\nu} \\
\end{longtable}
}

\EMsection{استفاده از تبدیل لاپلاس در حل معادلات دیفرانسیل}{استفاده از تبدیل لاپلاس در حل معادلات دیفرانسیل}

مراحل حل یک معادلهٔ دیفرانسیل با شرایط اولیه با استفاده از تبدیل لاپلاس:

\begin{enumerate}
\def\labelenumi{\arabic{enumi}.}
\tightlist
\item
  تبدیل مسئلهٔ مقدار اولیه به یک مسئلهٔ جبری؛
\item
  محاسبهٔ جواب مسئلهٔ جبری؛
\item
  حل مسئلهٔ مقدار اولیه (بازگرداندن جواب به حوزهٔ زمان).
\end{enumerate}

\EMfigure{m14-solve-scheme}{مراحل حل مسئلهٔ مقدار اولیه با تبدیل لاپلاس: مسئلهٔ مقدار اولیه
\EMim{\xrightarrow{1}} مسئلهٔ جبری \EMim{\xrightarrow{2}} جواب مسئلهٔ جبری
\EMim{\xrightarrow{3}} جواب مسئلهٔ مقدار اولیه}

\EMhold

\EMsubsection{معادلات دیفرانسیل معمولی}{معادلات دیفرانسیل معمولی}

\begin{EMexample}{}

معادلهٔ دیفرانسیل \EMim{y''-y=t}، \EMim{y(0)=1}، \EMim{y'(0)=1} را با استفاده از تبدیل لاپلاس حل کنید.

\EMdisp{s^2Y-sy(0)-y'(0)-Y\EMbr =\frac1{s^2}\ \EMbr \Longrightarrow\ (s^2-1)Y(s)\EMbr =s+1+\frac1{s^2}}
\EMdisp{Y(s)\EMbr =\frac{s+1}{s^2-1}+\frac1{s^2(s^2-1)}\EMbr =\frac1{s-1}+\left(\frac1{s^2-1}-\frac1{s^2}\right)}
\EMdisp{y(t)\EMbr =\mathcal{L}^{-1}\{Y(s)\}\EMbr =\mathcal{L}^{-1}\left\{\frac1{s-1}\right\}+\mathcal{L}^{-1}\left\{\frac1{s^2-1}\right\}-\mathcal{L}^{-1}\left\{\frac1{s^2}\right\}\EMbr =e^tu(t)+\sinh t\,u(t)-t\,u(t)}

\end{EMexample}

\EMhold

\EMsubsection{معادلات دیفرانسیل با مشتقات جزئی}{معادلات دیفرانسیل با مشتقات جزئی}

\begin{EMexample}{}

معادلهٔ دیفرانسیل با مشتقات جزئی زیر را با استفاده از تبدیل لاپلاس حل کنید.
\EMdisp{\frac{\partial w}{\partial x}+x\frac{\partial w}{\partial t}\EMbr =0,\qquad\EMbr  w(x,0)\EMbr =0,\qquad\EMbr  w(0,t)\EMbr =t\,u(t)}

از طرفین نسبت به \EMim{t} تبدیل لاپلاس می‌گیریم:
\EMdisp{\mathcal{L}_t\left\{\frac{\partial w}{\partial x}+x\frac{\partial w}{\partial t}\right\}\EMbr =\mathcal{L}_t\left\{\frac{\partial w}{\partial x}\right\}+x\,\mathcal{L}_t\left\{\frac{\partial w}{\partial t}\right\}\EMbr =0}
\EMdisp{\frac\partial{\partial x}w(x,s)+x\big[s\,w(x,s)-w(x,0)\big]\EMbr =0,\qquad\EMbr  w(x,0)\EMbr =0}
\EMdisp{\frac\partial{\partial x}w(x,s)+xs\,w(x,s)\EMbr =0\ \EMbr \Rightarrow\ \frac{dw}w\EMbr =-sx\,dx\ \EMbr \Rightarrow\ \ln w\EMbr =-s\frac{x^2}2+\ln C\ \EMbr \Rightarrow\ w(x,s)\EMbr =Ce^{-\frac{x^2}2s}}
از شرط \EMim{w(0,t)=t\,u(t)} داریم \EMim{w(0,s)=\mathcal{L}\{w(0,t)\}=1/s^2=C}:
\EMdisp{w(x,s)\EMbr =\frac1{s^2}e^{-\frac{x^2}2s}}
با استفاده از خاصیت شیفت زمانی \EMim{\mathcal{L}\{f(t-a)u(t-a)\}=e^{-as}F(s)}:
\EMdisp{w(x,t)\EMbr =\left(t-\frac{x^2}2\right)\times u\!\left(t-\frac{x^2}2\right)}

\end{EMexample}

\end{document}
